% Lutte contre les accidents cardiovasculaires — SVT 1ère TI — Cours signé OnBuch+
% Programme officiel MINESEC. Compiler avec : tectonic NCT03-lutte-accidents-cardiovasculaires.tex
\newcommand{\DOCMATIERE}{SVT}
\newcommand{\DOCNIVEAU}{1ère TI}
\newcommand{\DOCMODULE}{Module 2 : L'Éducation à la Santé}
\newcommand{\DOCLECON}{NCT03}
\newcommand{\DOCTITRE}{Lutte contre les accidents cardiovasculaires}
% OnBuch+ — préambule « cours signés » ENRICHI (compilé avec tectonic / XeLaTeX)
% Système visuel « Pop » d'OnBuch+ : fond crème (jamais blanc), bordures encre
% épaisses, ombres « dures » décalées, titres Archivo Black en capitales.
% Version MATHÉMATIQUES : graphes (pgfplots/tikz), boîtes pédagogiques
% (définition, propriété, méthode, attention, le savais-tu, exercice résolu,
% exercice, TP, bilan), page de garde et signature OnBuch+ ancrée en bas.
%
% Macros attendues dans main.tex AVANT \input{preamble} :
%   \DOCMATIERE  \DOCNIVEAU  \DOCMODULE  \DOCLECON (numéro)  \DOCTITRE
\documentclass[11pt]{article}

\usepackage{fontspec}
\usepackage[french]{babel}
\usepackage[a4paper,margin=2cm,top=3.4cm,bottom=2.8cm,headheight=1.4cm,footskip=1.3cm]{geometry}
\usepackage{amsmath,amssymb,amsfonts}
\usepackage{mathtools} % \xmapsto, \coloneqq... souvent utilisés par les modèles
\usepackage{cancel} % simplifications barrées (bilans de forces)
% \abs{z} (module, valeur absolue) et \norm{u} (norme), utilisés par les modèles
\providecommand{\abs}{}\let\abs\relax\DeclarePairedDelimiter{\abs}{\lvert}{\rvert}
\providecommand{\norm}{}\let\norm\relax\DeclarePairedDelimiter{\norm}{\lVert}{\rVert}
\usepackage[table]{xcolor} % [table] : fournit aussi \rowcolors (alternance de lignes)
\usepackage{pagecolor}
\usepackage{tikz}
\usetikzlibrary{arrows.meta,positioning,calc,shapes.geometric,shapes.misc,shapes.arrows,decorations.pathmorphing,decorations.pathreplacing,decorations.markings,patterns,shadows,mindmap,trees,fit,backgrounds,matrix,intersections,angles,quotes}
\usepackage{pgfplots}
\usepackage[version=4]{mhchem} % équations nucléaires \ce{^{235}_{92}U}
\usepackage[european,siunitx]{circuitikz} % schémas électriques
\pgfplotsset{compat=1.18}
\usepgfplotslibrary{fillbetween,groupplots} % aires entre courbes (calcul intégral)
\usepackage{enumitem}
\usepackage{fancyhdr}
% [most] charge toutes les bibliothèques tcolorbox (listings, minted, raster,
% documentation...) et gonfle inutilement la mémoire TeX sur un document long
% (~300 boîtes) : on ne charge que ce qui est réellement utilisé.
\usepackage[skins,breakable]{tcolorbox}
\usepackage{graphicx}
\usepackage{colortbl}
\usepackage{booktabs}
\usepackage{tabularx}
\usepackage{diagbox} % case d'en-tête diagonale des tableaux à double entrée
% Colonnes extensibles centrée / gauche / droite (C, L, R), souvent utilisées par les modèles
\newcolumntype{C}{>{\centering\arraybackslash}X}
\newcolumntype{L}{>{\raggedright\arraybackslash}X}
\newcolumntype{R}{>{\raggedleft\arraybackslash}X}
\usepackage{array}
\usepackage{multicol}
\usepackage{siunitx}
% parse-numbers=false : accepte aussi \SI{2,5 \times 10^{-3}}{...}
\sisetup{locale=FR,output-decimal-marker={,},group-minimum-digits=4,parse-numbers=false}
\DeclareSIUnit\ohm{\ensuremath{\symup{\Omega}}} % U+2126 absent de la police
\DeclareSIUnit\division{div} % réglages d'oscilloscope (ms/div, V/div)
\DeclareSIUnit\tr{tr} % tours (vitesse de rotation, ex. \SI{1500}{\tr\per\minute})
\DeclareSIUnit\ch{ch} % cheval-vapeur (puissance moteur, unité usuelle hors SI)
\DeclareSIUnit\franccfa{FCFA} % franc CFA (coûts, contexte camerounais)
\DeclareSIUnit\dioptre{\ensuremath{\delta}} % vergence d'une lentille (dioptrie)
\usepackage{qrcode}
% \degree : commande fréquemment utilisée par les modèles pour un angle ou une
% température en dehors de \SI{}{} (non fournie par les paquets chargés).
\providecommand{\degree}{^{\circ}}
% \qed (fin de démonstration), utilisé par les modèles sans amsthm
\providecommand{\qed}{\hfill$\square$}
% \barre : double barre des valeurs interdites dans un tableau de variations (tkz-tab)
\providecommand{\barre}{\ensuremath{\|}}
% environnement proof (amsthm non chargé) utilisé par les modèles
\ifdefined\proof\else\newenvironment{proof}[1][Démonstration]{\par\noindent\textit{#1.}\ }{\qed\par}\fi
\usepackage{lastpage}
\usepackage{hyperref}
\hypersetup{hidelinks}

% ============================================================
% Palette « Pop » OnBuch+ — mêmes hex que la classe P (plus_ui.dart)
% ============================================================
\definecolor{poporange}{HTML}{F59321}
\definecolor{popdark}{HTML}{E07A0C}
\definecolor{popcream}{HTML}{FFF6E8}
\definecolor{popink}{HTML}{1C1714}
\definecolor{poppurple}{HTML}{7A5AE0}
\definecolor{popgreen}{HTML}{2E9E6A}
\definecolor{poppink}{HTML}{E0457B}
\definecolor{popblue}{HTML}{2F7DE1}
\definecolor{popgold}{HTML}{C98A12}
\definecolor{popmuted}{HTML}{8A7F76}
\definecolor{popline}{HTML}{EADFD2}
\definecolor{popgray}{HTML}{8A7F76}
\colorlet{popgrey}{popgray}
% Alias tolérants (le modèle nomme parfois les couleurs différemment)
\colorlet{popwhite}{white}
\colorlet{popblack}{popink}
\colorlet{popred}{poppink}
\colorlet{popyellow}{popgold}
% Teintes claires pour fonds de tableaux / zones
\colorlet{popblueL}{popblue!10!white}
\colorlet{poporangeL}{poporange!14!white}
\colorlet{popgreenL}{popgreen!12!white}
\colorlet{poppurpleL}{poppurple!12!white}
\colorlet{poppinkL}{poppink!12!white}
\colorlet{popgoldL}{popgold!14!white}

\pagecolor{popcream}

% ============================================================
% Typographie
% ============================================================
\defaultfontfeatures{Ligatures=TeX}
\setmainfont{PlusJakartaSans-Medium.ttf}[Path=../../../fonts/,BoldFont=PlusJakartaSans-Bold.ttf]
% Maths en sans-serif (Fira Math) avec chiffres et lettres droites de Plus
% Jakarta, pour que les formules se fondent dans le texte.
\usepackage[math-style=ISO,bold-style=ISO]{unicode-math}
\setmathfont{FiraMath-Regular.otf}[Path=../../../fonts/]
\setmathfont{PlusJakartaSans-Medium.ttf}[Path=../../../fonts/,range={up/num,up/{Latin,latin},\mathup}]
\usepackage{icomma} % virgule décimale sans espace en mode math
% Caractères mathématiques Unicode absents des polices de texte (Plus Jakarta,
% Archivo Black) : rendus par leur équivalent mathématique, en texte comme en
% formule — sinon ils s'affichent en carré vide (ex. « Arithmétique dans ℤ »).
\usepackage{newunicodechar}
\newunicodechar{ℝ}{\ensuremath{\mathbb{R}}}
\newunicodechar{ℤ}{\ensuremath{\mathbb{Z}}}
\newunicodechar{ℂ}{\ensuremath{\mathbb{C}}}
\newunicodechar{ℕ}{\ensuremath{\mathbb{N}}}
\newunicodechar{ℚ}{\ensuremath{\mathbb{Q}}}
\newunicodechar{𝔻}{\ensuremath{\mathbb{D}}}
\newunicodechar{α}{\ensuremath{\alpha}}
\newunicodechar{β}{\ensuremath{\beta}}
\newunicodechar{γ}{\ensuremath{\gamma}}
\newunicodechar{δ}{\ensuremath{\delta}}
\newunicodechar{ε}{\ensuremath{\varepsilon}}
\newunicodechar{θ}{\ensuremath{\theta}}
\newunicodechar{λ}{\ensuremath{\lambda}}
\newunicodechar{μ}{\ensuremath{\mu}}
\newunicodechar{σ}{\ensuremath{\sigma}}
\newunicodechar{τ}{\ensuremath{\tau}}
\newunicodechar{φ}{\ensuremath{\varphi}}
\newunicodechar{ψ}{\ensuremath{\psi}}
\newunicodechar{ω}{\ensuremath{\omega}}
\newunicodechar{Σ}{\ensuremath{\Sigma}}
% majuscules produites par \MakeUppercase dans les titres (α-aminés -> Α)
\newunicodechar{Α}{\ensuremath{\alpha}}
\newunicodechar{Β}{\ensuremath{\beta}}
\newunicodechar{Γ}{\ensuremath{\Gamma}}
\newunicodechar{Δ}{\ensuremath{\Delta}}
\newunicodechar{Ε}{\ensuremath{\varepsilon}}
\newunicodechar{Θ}{\ensuremath{\Theta}}
\newunicodechar{Λ}{\ensuremath{\Lambda}}
\newunicodechar{Μ}{\ensuremath{\mu}}
\newunicodechar{Τ}{\ensuremath{\tau}}
\newunicodechar{Φ}{\ensuremath{\Phi}}
\newunicodechar{Ψ}{\ensuremath{\Psi}}
\newunicodechar{Ω}{\ensuremath{\Omega}}
\newunicodechar{↦}{\ensuremath{\mapsto}}
\newunicodechar{∈}{\ensuremath{\in}}
\newunicodechar{∉}{\ensuremath{\notin}}
\newunicodechar{∧}{\ensuremath{\wedge}}
\newunicodechar{⇔}{\ensuremath{\Leftrightarrow}}
\newunicodechar{⟺}{\ensuremath{\Longleftrightarrow}}
\newunicodechar{⇒}{\ensuremath{\Rightarrow}}
\newunicodechar{⟹}{\ensuremath{\Longrightarrow}}
\newunicodechar{∘}{\ensuremath{\circ}}
\newunicodechar{∪}{\ensuremath{\cup}}
\newunicodechar{∩}{\ensuremath{\cap}}
\newunicodechar{≡}{\ensuremath{\equiv}}
\newunicodechar{⊂}{\ensuremath{\subset}}
\newunicodechar{⊥}{\ensuremath{\perp}}
\newunicodechar{∃}{\ensuremath{\exists}}
\newunicodechar{∀}{\ensuremath{\forall}}
\newunicodechar{∛}{\ensuremath{\sqrt[3]{\,}}}
\newunicodechar{∜}{\ensuremath{\sqrt[4]{\,}}}
\newunicodechar{⃗}{\ensuremath{{}^{\to}}}
\newunicodechar{ⁿ}{\ifmmode^{n}\else\textsuperscript{\ensuremath{n}}\fi}
\newunicodechar{ˣ}{\ifmmode^{x}\else\textsuperscript{\ensuremath{x}}\fi}
\newunicodechar{ᵇ}{\ifmmode^{b}\else\textsuperscript{\ensuremath{b}}\fi}
\newunicodechar{ʳ}{\ifmmode^{r}\else\textsuperscript{\ensuremath{r}}\fi}
\newunicodechar{ᵘ}{\ifmmode^{u}\else\textsuperscript{\ensuremath{u}}\fi}
\newunicodechar{ʸ}{\ifmmode^{y}\else\textsuperscript{\ensuremath{y}}\fi}
\newunicodechar{ᵃ}{\ifmmode^{a}\else\textsuperscript{\ensuremath{a}}\fi}
\newunicodechar{ᵉ}{\ifmmode^{e}\else\textsuperscript{\ensuremath{e}}\fi}
\newunicodechar{ᵏ}{\ifmmode^{k}\else\textsuperscript{\ensuremath{k}}\fi}
\newunicodechar{ᵖ}{\ifmmode^{p}\else\textsuperscript{\ensuremath{p}}\fi}
\newunicodechar{ᴺ}{\ifmmode^{N}\else\textsuperscript{\ensuremath{N}}\fi}
\newunicodechar{ᶜ}{\ifmmode^{c}\else\textsuperscript{\ensuremath{c}}\fi}
\newunicodechar{ʰ}{\ifmmode^{h}\else\textsuperscript{\ensuremath{h}}\fi}
\newunicodechar{ᵠ}{\ifmmode^{q}\else\textsuperscript{\ensuremath{q}}\fi}
\newunicodechar{ₙ}{\ifmmode_{n}\else\textsubscript{\ensuremath{n}}\fi}
\newunicodechar{ₐ}{\ifmmode_{a}\else\textsubscript{\ensuremath{a}}\fi}
\newunicodechar{ᵢ}{\ifmmode_{i}\else\textsubscript{\ensuremath{i}}\fi}
\newunicodechar{ᵣ}{\ifmmode_{r}\else\textsubscript{\ensuremath{r}}\fi}
\newunicodechar{ₖ}{\ifmmode_{k}\else\textsubscript{\ensuremath{k}}\fi}
\newunicodechar{ᵦ}{\ifmmode_{\beta}\else\textsubscript{\ensuremath{\beta}}\fi}
\newunicodechar{ₓ}{\ifmmode_{x}\else\textsubscript{\ensuremath{x}}\fi}
\newfontfamily\popdisplay{ArchivoBlack-Regular.ttf}[Path=../../../fonts/]
\newfontfamily\poplogo{SpaceGrotesk-Bold.ttf}[Path=../../../fonts/]
\newfontfamily\popsemi{PlusJakartaSans-SemiBold.ttf}[Path=../../../fonts/]
\newfontfamily\popxbold{PlusJakartaSans-ExtraBold.ttf}[Path=../../../fonts/]
\newfontfamily\popxboldls{PlusJakartaSans-ExtraBold.ttf}[Path=../../../fonts/,LetterSpace=3.0]

\setlist[enumerate]{leftmargin=1.7em,itemsep=5pt,topsep=4pt}
\setlist[itemize]{leftmargin=1.7em,itemsep=4pt,topsep=4pt,label={\color{poporange}\textbullet}}
\setlength{\parindent}{0pt}
\setlength{\parskip}{5pt}
\renewcommand{\arraystretch}{1.25}

% pgfplots : style Pop par défaut
\pgfplotsset{
  every axis/.append style={axis line style={popink,line width=1pt},tick style={popink},
    grid style={popline,line width=0.5pt},label style={font=\small},tick label style={font=\footnotesize}},
  popcurve/.style={poporange,line width=1.8pt,no markers,smooth},
}
% Même style disponible dans un \draw TikZ simple (hors pgfplots)
\tikzset{popcurve/.style={poporange,line width=1.8pt}}
\tikzset{popgrid/.style={popline,very thin}}
\colorlet{popcurve}{poporange} % le nom du style est parfois utilisé comme couleur

% ============================================================
% Pill / squiggle
% ============================================================
\newcommand{\poppill}[3]{%
  \tikz[baseline=-0.55ex]\node[draw=popink,line width=1.2pt,rounded corners=2.6pt,%
    fill=#2,inner xsep=6.5pt,inner ysep=3pt]{%
    \popxboldls\fontsize{7}{7}\selectfont\color{#3}\MakeUppercase{#1}};%
}
\newcommand{\popsquiggle}[1]{%
  \tikz[baseline=-0.4ex]{
    \draw[#1,line width=2.6pt,line cap=round]
      plot [smooth] coordinates {(0,0.09) (0.34,0.01) (0.68,0.21) (1.02,0.01) (1.36,0.10)};
  }%
}
% Mot-clé surligné (marqueur orange sous le texte)
\newcommand{\cle}[1]{{\popxbold\color{popdark}#1}}

% ============================================================
% En-tête / pied
% ============================================================
\pagestyle{fancy}
\fancyhf{}
\renewcommand{\headrulewidth}{0pt}
\renewcommand{\footrulewidth}{0pt}
\lhead{\raisebox{-0.15ex}{\poplogo\fontsize{13}{13}\selectfont\color{popink}OnBuch\color{poporange}+}}
\rhead{\poppill{\DOCMATIERE\ \textbullet\ \DOCNIVEAU\ \textbullet\ Leçon \DOCLECON}{white}{popink}}
\lfoot{\popxbold\fontsize{7.6}{7.6}\selectfont\color{popmuted}\MakeUppercase{Cours signé OnBuch+}\ \textbullet\ {\popsemi\DOCTITRE}}
\rfoot{\tikz[baseline=-0.6ex]\node[rounded corners=8.5pt,fill=popink,minimum height=17pt,minimum width=17pt,inner xsep=5pt,inner sep=0]{\popxbold\fontsize{8}{8}\selectfont\color{popcream}\thepage\,/\,\pageref*{LastPage}};}

% ============================================================
% Titres
% ============================================================
\newcommand{\chaptitre}[2]{%
  \begin{tcolorbox}[enhanced,colback=poporange,colframe=popink,boxrule=2.6pt,arc=11pt,
      drop shadow=popink,left=15pt,right=15pt,top=12pt,bottom=11pt,before skip=3pt,after skip=18pt]
    \poppill{#2}{popink}{popcream}\par\vspace{9pt}
    {\raggedright\hyphenpenalty=10000\exhyphenpenalty=10000\popdisplay\fontsize{22}{24}\selectfont\color{popcream}\MakeUppercase{#1}\par}
  \end{tcolorbox}
}
% Section numérotée : pastille orange avec le numéro + titre Archivo
\newcounter{popsec}
\newcommand{\coursec}[1]{%
  \stepcounter{popsec}\setcounter{popsub}{0}%
  \par\vspace{16pt}\noindent
  \begin{minipage}{\linewidth}
  \tikz[baseline=-0.7ex]\node[draw=popink,line width=1.6pt,fill=poporange,rounded corners=4pt,minimum size=22pt,inner sep=0]{\popdisplay\fontsize{12}{12}\selectfont\color{popcream}\thepopsec};%
  \hspace{8pt}{\popdisplay\fontsize{15}{17}\selectfont\color{popink}\MakeUppercase{#1}}\par
  \vspace{4pt}\noindent{\color{poporange}\rule{36pt}{3.6pt}}
  \end{minipage}\par\nopagebreak\vspace{8pt}
}
\newcounter{popsub}
\newcommand{\courssub}[1]{%
  \stepcounter{popsub}%
  \par\vspace{9pt}\noindent{\popxbold\fontsize{11.5}{13}\selectfont\color{popink}{\color{poporange}\thepopsec.\thepopsub}\hspace{6pt}#1}\par\nopagebreak\vspace{4pt}
}

% ============================================================
% Boîtes pédagogiques — même recette : fond blanc, bordure épaisse colorée,
% ombre dure encre, badge plein. Titre optionnel [#1].
% ============================================================
\tcbset{popbase/.style={breakable,enhanced,colback=white,boxrule=2.3pt,arc=9pt,
  shadow={3pt}{-3pt}{0pt}{popink},left=14pt,right=14pt,top=11pt,bottom=11pt,before skip=11pt,after skip=15pt,
  fontupper=\popsemi\fontsize{10.6}{14.5}\selectfont\color{popink},
  fontlower=\popsemi\fontsize{10.6}{14.5}\selectfont\color{popink}}}
% Coche dessinée (le glyphe ✓ n'existe pas dans les polices du document)
\newcommand{\popcheck}[1]{\tikz[baseline=-0.5ex]\draw[#1,line width=1.6pt,line cap=round,line join=round](-0.13,0.0)--(-0.04,-0.09)--(0.13,0.1);}
\providecommand{\coche}{\popcheck{popgreen}}
\providecommand{\resultat}[1]{\par\smallskip\noindent\fcolorbox{popgreen}{popgreenL}{\parbox{\dimexpr\linewidth-2\fboxsep-2\fboxrule\relax}{\textbf{Résultat.} #1}}\par\smallskip}
\newcommand{\popboxhead}[3]{\poppill{#1}{#2}{white}\ifx\relax#3\relax\else\hspace{6pt}{\popxbold\fontsize{10.6}{12}\selectfont\color{popink}#3}\fi\par\vspace{7pt}}

\newtcolorbox{aretenir}[1][]{popbase,colframe=popgreen,before upper={\popboxhead{À retenir}{popgreen}{#1}}}
\newtcolorbox{exemplebox}[1][]{popbase,colframe=poporange,before upper={\popboxhead{Exemple}{poporange}{#1}}}
\newtcolorbox{definition}[1][]{popbase,colframe=popblue,before upper={\popboxhead{Définition}{popblue}{#1}}}
\newtcolorbox{propriete}[1][]{popbase,colframe=poppurple,before upper={\popboxhead{Propriété}{poppurple}{#1}}}
\newtcolorbox{methode}[1][]{popbase,colframe=popgold,before upper={\popboxhead{Méthode}{popgold}{#1}}}
\newtcolorbox{attention}[1][]{popbase,colframe=poppink,before upper={\popboxhead{Attention}{poppink}{#1}}}
\newtcolorbox{experience}[1][]{popbase,colframe=popblue,colback=popblueL,before upper={\popboxhead{Expérience}{popblue}{#1}}}
% « Le savais-tu ? » : carte violette pleine (contexte, culture, Cameroun)
\newtcolorbox{savaistu}[1][]{popbase,colback=poppurple,colframe=popink,
  fontupper=\popsemi\fontsize{10.4}{14.2}\selectfont\color{white},
  before upper={\poppill{Le savais-tu\,?}{white}{poppurple}\ifx\relax#1\relax\else\hspace{6pt}{\popxbold\color{white}#1}\fi\par\vspace{7pt}}}
% Exercice résolu : énoncé en haut, \tcblower puis la solution sur fond crème
\newcounter{popexo}\newcounter{popexr}
\newtcolorbox{exoresolu}[1][]{popbase,colframe=poporange,colbacklower=popcream,
  segmentation style={popink,line width=1.2pt,dashed},
  before upper={\stepcounter{popexr}\popboxhead{Exercice résolu \thepopexr}{poporange}{#1}},
  before lower={\poppill{Solution}{popink}{popcream}\par\vspace{6pt}}}
% Exercice (non résolu en place) — la correction est dans la partie Corrigés
\newtcolorbox{exercice}[1][]{popbase,colframe=popink,
  before upper={\stepcounter{popexo}\popboxhead{Exercice \thepopexo}{popink}{#1}}}
% Corrigé
\newtcolorbox{corrige}[1][]{popbase,colframe=popgreen,colback=popgreenL,
  before upper={\popboxhead{Corrigé}{popgreen}{#1}}}
% Objectifs / prérequis / bilan
\newtcolorbox{objectifs}{popbase,colframe=popink,colback=poporangeL,
  before upper={\popboxhead{Objectifs de la leçon}{popink}{}}}
\newtcolorbox{prerequis}{popbase,colframe=popink,colback=popblueL,
  before upper={\popboxhead{Prérequis}{popblue}{}}}
\newtcolorbox{bilan}{popbase,colframe=popink,colback=white,boxrule=2.8pt,
  before upper={\popboxhead{Fiche bilan}{poporange}{L'essentiel à connaître pour l'examen}}}
% Figure encadrée avec légende
\newtcolorbox{popfigure}[1][]{enhanced,breakable=false,colback=white,colframe=popline,boxrule=1.4pt,arc=8pt,
  left=10pt,right=10pt,top=10pt,bottom=8pt,before skip=10pt,after skip=12pt,halign=center,
  lower separated=false,halign lower=center,
  fontlower=\popsemi\fontsize{9}{11.5}\selectfont\color{popmuted},
  #1}
\newcommand{\legende}[1]{\par\vspace{6pt}{\popsemi\fontsize{9}{11.5}\selectfont\color{popmuted}#1}}

% ============================================================
% Signature OnBuch+
% ============================================================
\newcommand{\poplogoinline}[1]{{\poplogo\fontsize{#1}{#1}\selectfont\color{popink}OnBuch\color{poporange}+}}
\newcommand{\signatureonbuch}{%
  \begin{tcolorbox}[enhanced,colback=popink,colframe=popink,boxrule=2.2pt,arc=10pt,
      drop shadow=poporange,left=14pt,right=14pt,top=10pt,bottom=10pt,before skip=0pt,after skip=0pt]
    \tikz[baseline=-0.7ex]\node[circle,fill=popgreen,draw=popcream,line width=1.3pt,minimum size=22pt,inner sep=0]{\popcheck{white}};%
    \hspace{9pt}\begin{minipage}[c]{0.62\linewidth}
      {\popxbold\fontsize{12}{13}\selectfont\color{popcream}Signé\ \textbullet\ {\poplogo OnBuch\color{poporange}+}}\par\vspace{2pt}
      {\raggedright\popsemi\fontsize{8}{10.5}\selectfont\color{popline}\MakeUppercase{Équipe pédagogique OnBuch+}\\\MakeUppercase{Conforme au programme officiel MINESEC}\par}
    \end{minipage}\hfill
    {\poplogo\fontsize{22}{22}\selectfont\color{popcream}OnBuch\color{poporange}+}
  \end{tcolorbox}%
}

% ============================================================
% Page de garde
% #1 titre · #2 matière · #3 niveau · #4 module · #5 n° leçon · #6 durée ·
% #7 description / objectifs (texte ou itemize)
% ============================================================
\newcommand{\pagegarde}[7]{%
  \thispagestyle{empty}
  \newgeometry{a4paper,margin=1.7cm,top=1.6cm,bottom=1.4cm}%
  \begin{tikzpicture}[remember picture,overlay]
    \fill[poporange!18] ([xshift=-3.2cm,yshift=-3.6cm]current page.north east) circle (4.6cm);
    \fill[poppurple!14] ([xshift=2.4cm,yshift=7.6cm]current page.south west) circle (3.8cm);
  \end{tikzpicture}%
  {\poplogo\fontsize{30}{30}\selectfont\color{popink}OnBuch\color{poporange}+}\hfill
  \raisebox{0.6ex}{\poppill{Cours signé \textbullet\ Premium}{popink}{popcream}}\par
  \vspace{1pt}\popsquiggle{poppurple}\par
  \vspace{8pt}
  \begin{tcolorbox}[enhanced,colback=poporange,colframe=popink,boxrule=2.8pt,arc=13pt,
      drop shadow=popink,left=18pt,right=18pt,top=12pt,bottom=13pt,after skip=10pt]
    \poppill{#2 \textbullet\ #3}{popink}{popcream}\hspace{5pt}\poppill{#4}{white}{popink}\par\vspace{8pt}
    {\popdisplay\fontsize{14}{14}\selectfont\color{popink}LEÇON #5}\par\vspace{4pt}
    {\raggedright\hyphenpenalty=10000\exhyphenpenalty=10000\popdisplay\fontsize{25}{27}\selectfont\color{popcream}\MakeUppercase{#1}\par}
    \vspace{8pt}
    {\popxbold\fontsize{9.5}{9.5}\selectfont\color{popink}\MakeUppercase{Durée conseillée : #6}}
  \end{tcolorbox}
  \begin{tcolorbox}[enhanced,colback=white,colframe=popink,boxrule=2.2pt,arc=10pt,
      drop shadow=popink,left=16pt,right=16pt,top=10pt,bottom=10pt,after skip=8pt]
    \poppill{Dans cette leçon}{poppurple}{white}\par\vspace{5pt}
    \popsemi\fontsize{9.3}{12.6}\selectfont\color{popink}#7
  \end{tcolorbox}
  \noindent\begin{minipage}[t]{0.487\linewidth}
    \begin{tcolorbox}[enhanced,colback=popgreen,colframe=popink,boxrule=2.2pt,arc=10pt,
        drop shadow=popink,left=11pt,right=11pt,top=10pt,bottom=10pt,height=3.1cm,valign=center]
      \tcbox[colback=white,colframe=popink,boxrule=1.3pt,arc=5pt,boxsep=0pt,nobeforeafter,
        left=3pt,right=3pt,top=3pt,bottom=3pt,valign=center]{\qrcode[height=1.9cm]{https://play.google.com/store/apps/details?id=com.luvvix.onbuch&hl=fr}}%
      \hspace{8pt}\begin{minipage}[c]{0.5\linewidth}
        {\popdisplay\fontsize{11}{12.5}\selectfont\color{white}\MakeUppercase{Télécharge OnBuch}}\par\vspace{4pt}
        {\raggedright\popsemi\fontsize{8.4}{11}\selectfont\color{white}Gratuit \textbullet\ Google Play\\Tuteur IA, annales, concours, résultats\par}
      \end{minipage}
    \end{tcolorbox}
  \end{minipage}\hfill
  \begin{minipage}[t]{0.487\linewidth}
    \begin{tcolorbox}[enhanced,colback=poppurple,colframe=popink,boxrule=2.2pt,arc=10pt,
        drop shadow=popink,left=11pt,right=11pt,top=10pt,bottom=10pt,height=3.1cm,valign=center]
      \tikz[baseline=-0.7ex]\node[circle,fill=white,draw=popink,line width=1.3pt,minimum size=1.6cm,inner sep=0]{\popdisplay\fontsize{24}{24}\selectfont\color{poppurple}+};%
      \hspace{8pt}\begin{minipage}[c]{0.6\linewidth}
        {\popdisplay\fontsize{11}{12.5}\selectfont\color{white}\MakeUppercase{Passe à OnBuch+}}\par\vspace{4pt}
        {\raggedright\popsemi\fontsize{8.4}{11}\selectfont\color{white}Tous les cours signés, Tuteur IA illimité, fiches PDF — onglet OnBuch+ de l'app\par}
      \end{minipage}
    \end{tcolorbox}
  \end{minipage}
  \vspace{8pt}
  \popcontenu
  \vfill
  \signatureonbuch
  \restoregeometry
  \newpage
}

% Rangée « Ce cours contient » (page de garde)
\newcommand{\poptile}[3]{% #1 couleur #2 gros texte #3 légende
  \begin{tcolorbox}[enhanced,colback=#1,colframe=popink,boxrule=2pt,arc=9pt,shadow={2.5pt}{-2.5pt}{0pt}{popink},
      left=6pt,right=6pt,top=9pt,bottom=9pt,halign=center,valign=center,height=1.75cm,width=\linewidth,nobeforeafter]
    {\popdisplay\fontsize{12.5}{13}\selectfont\color{white}\MakeUppercase{#2}}\par\vspace{3pt}
    {\centering\popsemi\fontsize{7.8}{9.5}\selectfont\color{white}#3\par}
  \end{tcolorbox}}
\newcommand{\popcontenu}{%
  \par\noindent{\popxboldls\fontsize{7.5}{7.5}\selectfont\color{popmuted}CE COURS SIGNÉ CONTIENT}\par\vspace{7pt}
  \noindent\begin{minipage}[t]{0.235\linewidth}\poptile{popblue}{Cours}{complet, illustré, pas à pas}\end{minipage}\hfill
  \begin{minipage}[t]{0.235\linewidth}\poptile{poporange}{Méthodes}{et exercices résolus}\end{minipage}\hfill
  \begin{minipage}[t]{0.235\linewidth}\poptile{poppink}{Exercices}{type examen + corrigés}\end{minipage}\hfill
  \begin{minipage}[t]{0.235\linewidth}\poptile{popgreen}{Fiche bilan}{l'essentiel à retenir}\end{minipage}\par}

% Sommaire
\newtcolorbox{sommaire}{enhanced,colback=white,colframe=popink,boxrule=2.2pt,arc=10pt,
  drop shadow=popink,left=18pt,right=18pt,top=13pt,bottom=13pt,before skip=6pt,after skip=20pt,
  before upper={\poppill{Sommaire}{poppurple}{white}\par\vspace{9pt}\popsemi\fontsize{10.6}{14.5}\selectfont\color{popink}}}

% Signature de fin de cours — poussée tout en bas de la dernière page
\newcommand{\findecours}{%
  \par\vfill\nopagebreak
  \begin{center}\popsquiggle{poporange}\end{center}\vspace{4pt}
  \signatureonbuch
}
\AtBeginDocument{\def\llbracket{\mathopen{[\mkern-3.5mu[}}\def\rrbracket{\mathclose{]\mkern-3.5mu]}}} % intervalles d'entiers (glyphe ⟦ absent de la police)
% Glyphes absents de Fira Math : redéfinis à partir de glyphes disponibles
\AtBeginDocument{%
  \def\setminus{\mathbin{\backslash}}\let\smallsetminus\setminus
  \def\vdots{\mathord{\vbox{\baselineskip3.2pt\lineskiplimit0pt\kern4pt\hbox{.}\hbox{.}\hbox{.}}}}%
  \def\ddots{\mathinner{\mkern1mu\raise6.5pt\hbox{.}\mkern2mu\raise3.6pt\hbox{.}\mkern2mu\raise0.7pt\hbox{.}\mkern1mu}}%
  \def\triangle{\mathord{\scalebox{0.9}{$\symup{\Delta}$}}}%
  \def\nabla{\mathord{\raisebox{\depth}{\scalebox{1}[-1]{$\symup{\Delta}$}}}}%
  \def\checkmark{\popcheck{popdark}}%
  \def\star{\mathbin{\ast}}\def\bigstar{\mathord{\ast}}%
  \def\diamond{\mathbin{\scalebox{0.6}{$\Diamond$}}}%
  \def\bigcup{\mathop{\mathchoice{\raisebox{-0.3em}{\scalebox{1.7}{$\cup$}}}{\scalebox{1.25}{$\cup$}}{\cup}{\cup}}\displaylimits}%
  \def\bigcap{\mathop{\mathchoice{\raisebox{-0.3em}{\scalebox{1.7}{$\cap$}}}{\scalebox{1.25}{$\cap$}}{\cap}{\cap}}\displaylimits}%
  \def\lesseqqgtr{\mathrel{\lessgtr}}\def\bigoplus{\mathop{\oplus}\displaylimits}%
}
\providecommand{\ssi}{si et seulement si} % abréviation fréquente
\AtBeginDocument{\providecommand{\wideparen}{\overparen}} % arc de cercle

%% FIGFIX-BEGIN v4 — correctifs globaux des figures (docs/onbuch-generation/tools/figfix)
\usepackage{adjustbox}\usepackage{etoolbox}
% toute figure TikZ/pgfplots plus large que la ligne est réduite à la largeur disponible
\BeforeBeginEnvironment{tikzpicture}{\begin{adjustbox}{max width=\linewidth}}
\AfterEndEnvironment{tikzpicture}{\end{adjustbox}}
% légendes pgfplots lisibles par-dessus les courbes
\pgfplotsset{every axis/.append style={legend style={fill=white,fill opacity=0.92,text opacity=1,draw=black!15,inner sep=3pt}}}
% étiquettes d'arêtes (syntaxe edge["..."]) avec fond blanc
\tikzset{every edge quotes/.append style={fill=white,inner sep=1.2pt}}
%% FIGFIX-END
\begin{document}

\pagegarde{Lutte contre les accidents cardiovasculaires}{SVT}{1ère TI}{Module 2 : L'Éducation à la Santé}{NCT03}{2 h}{Cette leçon présente le fonctionnement simplifié du système cardiovasculaire, les principaux accidents qui le menacent (infarctus, AVC, hypertension) et leurs facteurs de risque. Elle outille l'élève pour identifier ces facteurs, interpréter une mesure de tension artérielle ou de fréquence cardiaque et proposer des gestes de prévention adaptés au contexte camerounais.
\vspace{4pt}
\begin{multicols}{2}\small
\begin{itemize}
\item À la fin de cette leçon, je sais décrire le fonctionnement simplifié du système cardiovasculaire (cœur, vaisseaux, double circulation)
\item Schématiser et légender le trajet du sang dans le cœur et les grands vaisseaux
\item Définir l'infarctus du myocarde, l'accident vasculaire cérébral (AVC) et l'hypertension artérielle
\item Identifier les facteurs de risque cardiovasculaire dans une situation de vie donnée
\item Interpréter une mesure de tension artérielle (ex. 12/8, 16/10) et une fréquence cardiaque
\item Distinguer les facteurs de risque modifiables des facteurs non modifiables
\end{itemize}\end{multicols}}

\begin{sommaire}
\begin{multicols}{2}
\begin{enumerate}
\item Le système cardiovasculaire : structure et fonctionnement
\item La tension artérielle et sa mesure
\item Les principaux accidents cardiovasculaires
\item Les facteurs de risque cardiovasculaire
\item La prévention des accidents cardiovasculaires
\item Synthèse et évaluation des acquis
\item Activité d'intégration
\item Méthodes et exercices résolus
\item Exercices
\item Corrigés des exercices
\item Fiche bilan
\end{enumerate}
\end{multicols}
\end{sommaire}

\begin{objectifs}
\begin{itemize}
\item À la fin de cette leçon, je sais décrire le fonctionnement simplifié du système cardiovasculaire (cœur, vaisseaux, double circulation)
\item Schématiser et légender le trajet du sang dans le cœur et les grands vaisseaux
\item Définir l'infarctus du myocarde, l'accident vasculaire cérébral (AVC) et l'hypertension artérielle
\item Identifier les facteurs de risque cardiovasculaire dans une situation de vie donnée
\item Interpréter une mesure de tension artérielle (ex. 12/8, 16/10) et une fréquence cardiaque
\item Distinguer les facteurs de risque modifiables des facteurs non modifiables
\item Proposer des mesures de prévention adaptées : alimentation équilibrée, activité physique, dépistage régulier
\item Sensibiliser mon entourage à la lutte contre les accidents cardiovasculaires
\end{itemize}
\end{objectifs}

\begin{prerequis}
\begin{itemize}
\item Notions sur le sang et ses constituants (classe de Seconde)
\item Organisation générale du corps humain en appareils et systèmes
\item Notion de circulation sanguine vue au collège (BEPC)
\item Lecture et interprétation de valeurs numériques simples (comparaison, moyenne)
\item Notions d'hygiène de vie et d'alimentation équilibrée (Éducation à la santé)
\end{itemize}
\end{prerequis}

\newpage

\coursec{Le système cardiovasculaire : structure et fonctionnement}

\begin{attention}[Situation d'accroche : un drame évitable]
Pendant un match de football au quartier à Yaoundé, un supporter de 52 ans s'effondre soudainement en se tenant la poitrine. Les témoins évoquent ses habitudes : consommation quotidienne de plats très salés, absence d'activité physique et tabagisme. À l'hôpital, le médecin diagnostique un \cle{infarctus du myocarde} et annonce que sa tension artérielle était de 17/11 depuis des années sans qu'il le sache. Comment un tel drame aurait-il pu être évité ? Pour répondre à cette question, il faut d'abord comprendre comment fonctionne la « machine » qui fait circuler le sang dans notre corps : le \cle{système cardiovasculaire}. Cette première partie de la leçon nous donne les bases indispensables pour, ensuite, comprendre les accidents cardiovasculaires et apprendre à les prévenir.
\end{attention}

\textbf{Problématique :} comment le cœur et les vaisseaux sanguins assurent-ils en permanence le transport du sang dans tout l'organisme, et que se passe-t-il lorsque cette machinerie se dérègle ?

\courssub{Qu'est-ce que le système cardiovasculaire ?}

\textbf{Observation préliminaire.} Lorsque tu cours derrière un ballon, tu sens ton cœur battre fort dans ta poitrine et ton pouls s'accélérer à ton poignet. Lorsque tu te coupes légèrement le doigt en épluchant des macabos, le sang coule aussitôt. Ces observations simples montrent qu'il existe dans notre corps un liquide, le \cle{sang}, mis en mouvement par un organe, le \cle{cœur}, et circulant dans des conduits, les \cle{vaisseaux sanguins}. Ces trois éléments forment un système intégré.

\begin{definition}[Le système cardiovasculaire]
Le \cle{système cardiovasculaire} (ou appareil circulatoire) est l'ensemble formé par le \textbf{cœur}, les \textbf{vaisseaux sanguins} (artères, veines et capillaires) et le \textbf{sang}. Il assure le transport dans tout l'organisme :
\begin{itemize}
\item du \textbf{dioxygène} ($O_2$) prélevé dans les poumons et des \textbf{nutriments} absorbés au niveau de l'intestin, vers toutes les cellules ;
\item des \textbf{déchets} (dioxyde de carbone $CO_2$, urée...) produits par les cellules, vers les organes d'élimination (poumons, reins).
\end{itemize}
C'est le « réseau de distribution et d'évacuation » du corps humain.
\end{definition}

\begin{attention}[Piège fréquent]
Ne confonds pas \emph{système cardiovasculaire} et \emph{système respiratoire}. Le système respiratoire (poumons) fait entrer l'$O_2$ et sortir le $CO_2$ ; le système cardiovasculaire \emph{transporte} ces gaz dans le sang. Les deux travaillent en étroite collaboration, mais ce sont deux systèmes distincts.
\end{attention}

\courssub{Le cœur : une double pompe musculaire}

\textbf{Observation.} Dissecte (en imagination ou en TP) un cœur de bœuf acheté au marché : c'est un organe ferme, de la taille d'un poing fermé, essentiellement constitué de muscle. En le coupant en deux, on découvre qu'il est \textbf{creux} : il contient des cavités séparées par une cloison.

\begin{definition}[Le cœur]
Le cœur est un organe \textbf{musculaire et creux}, situé dans la cage thoracique, entre les deux poumons, légèrement décalé vers la gauche. Sa paroi musculaire, le \cle{myocarde}, se contracte de façon rythmique et automatique, sans intervention de la volonté, toute la vie durant. Le cœur est divisé en \textbf{quatre cavités} :
\begin{itemize}
\item deux cavités supérieures, les \cle{oreillettes} (oreillette droite et oreillette gauche), qui reçoivent le sang ;
\item deux cavités inférieures, les \cle{ventricules} (ventricule droit et ventricule gauche), à paroi plus épaisse, qui propulsent le sang.
\end{itemize}
Une cloison, le \textbf{septum}, sépare complètement le cœur droit du cœur gauche : le sang ne passe jamais directement d'un côté à l'autre.
\end{definition}

\textbf{Le rôle des valvules.} Entre chaque oreillette et son ventricule, ainsi qu'à la sortie des ventricules, se trouvent des \cle{valvules}. Ce sont des replis membraneux qui s'ouvrent dans un seul sens, comme une porte à battant unique.

\begin{propriete}[Rôle des valvules — admise]
Les valvules cardiaques empêchent le \textbf{reflux} du sang : elles imposent au sang un trajet à sens unique :
\[ \text{oreillette} \longrightarrow \text{ventricule} \longrightarrow \text{artère}. \]
Sans elles, une partie du sang reviendrait en arrière à chaque contraction et le cœur perdrait son efficacité.
\end{propriete}

\begin{propriete}[Le cœur, une double pompe — admise]
Le cœur fonctionne comme \textbf{deux pompes accolées} travaillant en même temps :
\begin{itemize}
\item le \textbf{cœur droit} (oreillette droite + ventricule droit) reçoit le sang pauvre en $O_2$ venant des organes et le propulse vers les \textbf{poumons} ;
\item le \textbf{cœur gauche} (oreillette gauche + ventricule gauche) reçoit le sang riche en $O_2$ venant des poumons et le propulse vers \textbf{tout le corps}.
\end{itemize}
La paroi du ventricule gauche est plus épaisse que celle du ventricule droit, car il doit pousser le sang beaucoup plus loin (jusqu'aux pieds !).
\end{propriete}

\begin{savaistu}[Le cœur a aussi besoin d'être nourri]
Le myocarde travaille sans relâche : il a donc besoin de son propre apport en $O_2$ et nutriments. Il est vascularisé par les \textbf{artères coronaires} (ou coronariques), qui naissent à la base de l'aorte. Si l'une d'elles se bouche (par un caillot sur une plaque d'athérome), le myocarde en aval n'est plus oxygéné : c'est l'\textbf{infarctus du myocarde}, le drame de notre supporter de Yaoundé.
\end{savaistu}

\begin{popfigure}
\begin{center}
\begin{tikzpicture}[scale=1.05, every node/.style={font=\small}]
% Contour du cœur (vue de face, schématique)
\draw[very thick, poppink, fill=poppinkL] (0,0.4) .. controls (-2.6,0.9) and (-2.4,4.4) .. (-0.9,4.6)
  .. controls (-0.2,4.75) and (0,4.3) .. (0,4.1)
  .. controls (0,4.3) and (0.2,4.75) .. (0.9,4.6)
  .. controls (2.4,4.4) and (2.6,0.9) .. (0,0.4) -- cycle;
% Septum
\draw[very thick, poppink] (0,0.55) -- (0,4.05);
% Cloison oreillette/ventricule
\draw[thick, poppink] (-1.9,2.6) -- (-0.15,2.6);
\draw[thick, poppink] (0.15,2.6) -- (1.9,2.6);
% Valvules
\draw[thick, popdark] (-0.35,2.6) -- (-0.05,2.25);
\draw[thick, popdark] (0.35,2.6) -- (0.05,2.25);
% Gros vaisseaux
\draw[very thick, popblue] (0.5,4.3) -- (0.5,5.6) -- (1.5,6.1); % aorte
\draw[very thick, poppurple] (-0.6,4.3) -- (-0.6,5.2) -- (-1.6,5.7); % artère pulmonaire
\draw[very thick, poppurple] (-1.15,4.5) -- (-1.15,6.0); % veine cave sup
\draw[very thick, poppurple] (-1.35,0.9) -- (-1.35,-0.4); % veine cave inf
\draw[very thick, popblue] (1.5,4.0) -- (2.3,4.6); % veines pulmonaires
\draw[very thick, popblue] (1.6,3.5) -- (2.5,3.7);
% Flèches du trajet du sang
\draw[-{Stealth[length=3mm]}, very thick, poppurple] (-1.15,5.4) -- (-1.0,3.6);
\draw[-{Stealth[length=3mm]}, very thick, poppurple] (-1.35,0.2) -- (-1.2,1.8);
\draw[-{Stealth[length=3mm]}, very thick, poppurple] (-0.9,2.3) -- (-0.75,1.2);
\draw[-{Stealth[length=3mm]}, very thick, poppurple] (-0.7,1.6) -- (-0.62,4.6);
\draw[-{Stealth[length=3mm]}, very thick, popblue] (2.2,4.4) -- (1.0,3.6);
\draw[-{Stealth[length=3mm]}, very thick, popblue] (0.9,2.3) -- (0.8,1.1);
\draw[-{Stealth[length=3mm]}, very thick, popblue] (0.7,1.5) -- (0.55,4.9);
% Labels cavités
\node[popdark] at (-1.0,3.3) {\textbf{OD}};
\node[popdark] at (-0.85,1.5) {\textbf{VD}};
\node[popdark] at (1.0,3.3) {\textbf{OG}};
\node[popdark] at (0.85,1.5) {\textbf{VG}};
% Légendes pointées
\node[right, popblue] at (1.7,6.1) {Aorte (vers le corps)};
\node[left, poppurple] at (-1.8,5.7) {Artère pulmonaire (vers les poumons)};
\node[left, poppurple] at (-1.35,6.0) {Veine cave supérieure};
\node[left, poppurple] at (-1.55,-0.4) {Veine cave inférieure};
\node[right, popblue] at (2.5,4.15) {Veines pulmonaires};
\node[right, popdark] at (0.45,2.35) {Valvules};
\end{tikzpicture}
\end{center}
\legende{Schéma simplifié du cœur en coupe frontale (vue de face : le cœur droit apparaît à gauche du schéma). OD : oreillette droite ; VD : ventricule droit ; OG : oreillette gauche ; VG : ventricule gauche. Les flèches violettes indiquent le trajet du sang désoxygéné, les flèches bleues celui du sang oxygéné.}
\end{popfigure}

\courssub{Les vaisseaux sanguins : un réseau de conduits spécialisés}

\textbf{Observation.} Lors d'une prise de sang à l'hôpital, l'infirmier pique dans un vaisseau bleuâtre visible sous la peau du pli du coude : c'est une \textbf{veine}. En revanche, lorsqu'une artère est sectionnée (accident grave), le sang jaillit par jets puissants, au rythme des battements du cœur. Cette différence s'explique par la structure et le rôle de chaque type de vaisseau.

\begin{definition}[Les trois types de vaisseaux]
\begin{itemize}
\item Les \cle{artères} sont les vaisseaux qui transportent le sang \textbf{du cœur vers les organes}. Leur paroi est épaisse et élastique, car elles supportent la forte pression du sang éjecté par les ventricules. La plus grosse artère du corps est l'\textbf{aorte}.
\item Les \cle{veines} sont les vaisseaux qui ramènent le sang \textbf{des organes vers le cœur}. Leur paroi est plus fine ; beaucoup possèdent des valvules (les « valvules veineuses ») qui empêchent le sang de redescendre, surtout dans les jambes.
\item Les \cle{capillaires} sont des vaisseaux microscopiques, à paroi extrêmement fine (une seule couche de cellules), qui relient les artères aux veines. C'est à leur niveau que se font les \textbf{échanges} entre le sang et les cellules : $O_2$ et nutriments sortent du sang, $CO_2$ et déchets y entrent.
\end{itemize}
\end{definition}

\begin{attention}[Piège classique au Probatoire]
« Artère » ne signifie pas « sang oxygéné » et « veine » ne signifie pas « sang désoxygéné » ! L'\textbf{artère pulmonaire} transporte du sang \emph{désoxygéné} (du cœur vers les poumons) et les \textbf{veines pulmonaires} transportent du sang \emph{oxygéné} (des poumons vers le cœur). Retiens la définition par la \emph{direction du flux} : artère = le sang quitte le cœur ; veine = le sang revient au cœur.
\end{attention}

\courssub{La double circulation : deux circuits en boucle fermée}

\textbf{Question.} Le sang qui part du cœur revient-il directement ? Non. Le sang effectue en permanence \textbf{deux circuits} qui s'enchaînent, ce qui justifie l'expression « double pompe ».

\begin{definition}[La petite circulation et la grande circulation]
\begin{itemize}
\item La \cle{petite circulation} (ou circulation pulmonaire) : \textbf{cœur droit $\rightarrow$ poumons $\rightarrow$ cœur gauche}. Le ventricule droit chasse le sang désoxygéné dans l'artère pulmonaire jusqu'aux poumons ; là, au niveau des capillaires entourant les alvéoles pulmonaires, le sang se charge en $O_2$ et rejette le $CO_2$ : c'est l'\textbf{oxygénation du sang}. Le sang oxygéné revient au cœur gauche par les veines pulmonaires.
\item La \cle{grande circulation} (ou circulation générale) : \textbf{cœur gauche $\rightarrow$ organes $\rightarrow$ cœur droit}. Le ventricule gauche chasse le sang oxygéné dans l'aorte, qui se ramifie vers tous les organes (cerveau, muscles, reins, intestin...). Au niveau des capillaires, le sang distribue l'$O_2$ et les nutriments et collecte le $CO_2$ et les déchets. Le sang désoxygéné revient au cœur droit par les veines caves.
\end{itemize}
Le trajet complet s'écrit donc :
\[ \text{VD} \rightarrow \text{poumons} \rightarrow \text{OG} \rightarrow \text{VG} \rightarrow \text{organes} \rightarrow \text{OD} \rightarrow \text{VD} \rightarrow \dots \]
\end{definition}

\begin{popfigure}
\begin{center}
\begin{tikzpicture}[scale=0.95, every node/.style={font=\small}]
% Coeur
\draw[very thick, poppink, fill=poppinkL] (0,0) ellipse (1.1 and 1.4);
\draw[thick, poppink] (0,-1.35) -- (0,1.35);
\node[popdark] at (-0.55,0.9) {\textbf{Cœur}};
\node[popdark] at (-0.55,0.45) {\textbf{droit}};
\node[popdark] at (0.55,0.9) {\textbf{Cœur}};
\node[popdark] at (0.55,0.45) {\textbf{gauche}};
% Poumons
\draw[very thick, popblue, fill=popblueL] (-3.4,2.6) ellipse (1.0 and 0.8);
\draw[very thick, popblue, fill=popblueL] (-1.6,2.6) ellipse (1.0 and 0.8);
\node[popdark] at (-2.5,2.6) {\textbf{Poumons}};
% Organes
\draw[very thick, popgreen, fill=popgreenL, rounded corners] (-4.2,-3.6) rectangle (-0.8,-2.4);
\node[popdark] at (-2.5,-3.0) {\textbf{Organes du corps}};
% Petite circulation (sang désoxygéné aller, oxygéné retour)
\draw[-{Stealth[length=3mm]}, very thick, poppurple] (-0.9,0.7) .. controls (-2.6,0.9) and (-3.2,1.2) .. (-3.2,1.9);
\draw[-{Stealth[length=3mm]}, very thick, popblue] (-1.8,1.9) .. controls (-1.4,1.3) and (-0.6,1.6) .. (0.35,1.15);
\node[poppurple, align=center] at (-4.35,1.15) {petite\\circulation};
% Grande circulation
\draw[-{Stealth[length=3mm]}, very thick, popblue] (0.9,-0.5) .. controls (2.4,-1.0) and (2.4,-3.0) .. (-0.7,-3.0);
\draw[-{Stealth[length=3mm]}, very thick, poppurple] (-4.3,-3.0) .. controls (-5.6,-3.0) and (-5.6,-0.6) .. (-1.05,-0.55);
\node[popblue, align=center] at (2.45,-2.0) {grande\\circulation};
% Légende couleurs
\draw[very thick, popblue] (2.2,1.6) -- (3.0,1.6); \node[right] at (3.0,1.6) {sang oxygéné};
\draw[very thick, poppurple] (2.2,1.0) -- (3.0,1.0); \node[right] at (3.0,1.0) {sang désoxygéné};
% Echanges
\node[popdark, align=center] at (-2.5,4.0) {échanges : $O_2$ entre, $CO_2$ sort};
\node[popdark, align=center] at (-2.5,-4.3) {échanges : $O_2$ et nutriments sortent, $CO_2$ et déchets entrent};
\end{tikzpicture}
\end{center}
\legende{La double circulation. Petite circulation (cœur droit $\rightarrow$ poumons $\rightarrow$ cœur gauche) : oxygénation du sang. Grande circulation (cœur gauche $\rightarrow$ organes $\rightarrow$ cœur droit) : distribution de l'$O_2$ et des nutriments, collecte des déchets.}
\end{popfigure}

\courssub{Le cycle cardiaque : systole et diastole}

\textbf{Observation.} Pose ta main sur ta poitrine : tu sens des battements réguliers, « boum-boum... boum-boum... ». Chaque battement correspond à un \cle{cycle cardiaque}, qui comporte deux temps.

\begin{definition}[Le cycle cardiaque]
Le cycle cardiaque est l'ensemble des événements qui se déroulent pendant un battement. Il comprend :
\begin{itemize}
\item la \cle{systole} : phase de \textbf{contraction} du myocarde. Les oreillettes se contractent d'abord, puis les ventricules, qui \textbf{éjectent} le sang dans les artères (aorte et artère pulmonaire) ;
\item la \cle{diastole} : phase de \textbf{relaxation} du myocarde. Les cavités se relâchent et se \textbf{remplissent} de sang venu des veines.
\end{itemize}
Au repos, un cycle dure environ \num{0,8}~s. Le bruit « boum-boum » entendu au stéthoscope correspond à la fermeture des valvules.
\end{definition}

\begin{savaistu}[Un travailleur infatigable]
Le cœur d'un adulte bat environ \textbf{100 000 fois par jour} et pompe près de \textbf{7 500 litres de sang} quotidiennement, soit l'équivalent de plus de 100 baignoires ! Et cela sans jamais s'arrêter, de la vie intra-utérine jusqu'à la mort. Le myocarde est le muscle le plus endurant du corps — à condition d'en prendre soin.
\end{savaistu}

\courssub{La fréquence cardiaque : mesure et interprétation}

\begin{definition}[Fréquence cardiaque]
La \cle{fréquence cardiaque} (FC) est le nombre de battements du cœur par minute. Chez l'adulte au repos, elle est normalement comprise entre \textbf{60 et 100 battements/min}. Elle augmente à l'effort, sous l'effet du stress ou de la fièvre, et diminue pendant le sommeil. Chez le sportif entraîné, elle peut descendre au repos jusqu'à 50 battements/min : c'est un signe de bonne condition physique, pas une maladie.
\end{definition}

\begin{methode}[Mesurer sa fréquence cardiaque au poignet]
\begin{enumerate}
\item S'asseoir et se reposer quelques minutes (mesure « au repos »).
\item Retourner la paume de la main vers le haut ; poser \textbf{deux doigts} (index et majeur, jamais le pouce qui a son propre pouls) sur la face interne du poignet, du côté du pouce : c'est l'\textbf{artère radiale}.
\item Sentir les pulsations, puis compter les battements pendant \textbf{30 secondes}.
\item Multiplier le résultat par \textbf{2} pour obtenir la fréquence cardiaque en battements/min.
\end{enumerate}
\end{methode}

\begin{exemplebox}[Un calcul simple]
Un élève de Première C compte \textbf{35 pulsations en 30 s} au repos. Sa fréquence cardiaque est :
\[ FC = 35 \times 2 = 70 \text{ battements/min}. \]
Cette valeur est comprise entre 60 et 100 battements/min : elle est donc \textbf{normale}.
\end{exemplebox}

\begin{experience}[TP : effet de l'effort sur la fréquence cardiaque]
\textbf{Matériel :} un chronomètre (ou téléphone), un escabeau ou une marche d'escalier.
\textbf{Protocole :}
\begin{enumerate}
\item Mesurer la FC au repos, assis, pendant 30 s (multiplier par 2).
\item Effectuer des montées-descentes de marche pendant 3 minutes.
\item Mesurer la FC immédiatement à l'arrêt de l'effort, puis toutes les minutes pendant 5 minutes de récupération.
\item Reporter les valeurs dans un tableau et tracer la courbe $FC = f(t)$.
\end{enumerate}
\textbf{Observations attendues :} la FC augmente fortement pendant l'effort (les muscles réclament plus d'$O_2$), puis diminue progressivement pendant la récupération jusqu'à retrouver sa valeur de repos.
\textbf{Interprétation :} le cœur adapte son débit aux besoins de l'organisme. Un retour rapide à la normale traduit un cœur en bonne santé.
\end{experience}

\begin{popfigure}
\begin{center}
\begin{tikzpicture}
\begin{axis}[
  width=\linewidth, height=6.5cm,
  xlabel={Temps (min)}, ylabel={Fréquence cardiaque (battements/min)},
  xmin=0, xmax=12, ymin=50, ymax=160,
  xtick={0,1,2,3,4,5,6,7,8,9,10,11,12},
  ytick={60,80,100,120,140},
  grid=major, grid style={popline!50},
  legend style={at={(0.98,0.35)}, anchor=east, font=\small},
]
\addplot[popcurve, popblue, very thick, mark=*, mark size=1.5pt] coordinates {
  (0,72) (1,95) (2,120) (3,138) (4,128) (5,110) (6,95) (7,84) (8,78) (9,74) (10,72)
};
\addlegendentry{FC mesurée chez un élève}
\addplot[poporange, dashed, thick, domain=0:12] {100};
\addlegendentry{Limite supérieure normale au repos (100)}
\draw[popmuted, dashed] (axis cs:3,50) -- (axis cs:3,138);
\node[popdark, font=\small, anchor=west] at (axis cs:3.1,60) {fin de l'effort};
\end{axis}
\end{tikzpicture}
\end{center}
\legende{Évolution de la fréquence cardiaque d'un élève au repos (0--1 min), pendant un effort de 3 min (montées de marche) puis pendant la récupération. La FC monte jusqu'à 138 battements/min puis redescend progressivement vers sa valeur de repos (72 battements/min) : le cœur est sain.}
\end{popfigure}

\begin{exoresolu}[Exploiter des mesures de fréquence cardiaque]
\textbf{Énoncé.} Lors d'un TP, trois élèves mesurent leur fréquence cardiaque au repos en comptant les pulsations de l'artère radiale pendant 30 s : Aïcha compte 32 pulsations, Bernard compte 47 pulsations, Carine compte 38 pulsations.
\begin{enumerate}
\item Calculer la fréquence cardiaque de chaque élève.
\item Indiquer, en justifiant, si chaque valeur est normale.
\item Après 3 minutes d'effort, Aïcha mesure 145 battements/min. Expliquer cette augmentation.
\end{enumerate}
\tcblower
\textbf{Solution détaillée.}
\begin{enumerate}
\item On multiplie par 2 le nombre de pulsations comptées en 30 s :
\begin{align*}
FC_{\text{Aïcha}} &= 32 \times 2 = 64 \text{ battements/min}\\
FC_{\text{Bernard}} &= 47 \times 2 = 94 \text{ battements/min}\\
FC_{\text{Carine}} &= 38 \times 2 = 76 \text{ battements/min}
\end{align*}
\item Les trois valeurs sont comprises entre 60 et 100 battements/min : elles sont toutes \textbf{normales}. Celle de Bernard (94) est proche de la limite supérieure ; elle reste acceptable mais pourrait inviter à vérifier qu'il était bien au repos.
\item Pendant l'effort, les muscles consomment beaucoup plus de dioxygène et de nutriments et produisent davantage de $CO_2$. Pour répondre à ces besoins accrus, le cœur accélère : la fréquence cardiaque augmente afin d'augmenter le débit sanguin vers les muscles. Une FC de 145 battements/min à l'effort est donc une réponse physiologique normale.
\end{enumerate}
\end{exoresolu}

\begin{aretenir}[L'essentiel de la section]
\begin{itemize}
\item Le système cardiovasculaire = \textbf{cœur + vaisseaux + sang} ; il transporte l'$O_2$, les nutriments et les déchets.
\item Le cœur est une \textbf{double pompe} à 4 cavités (2 oreillettes, 2 ventricules) ; les \textbf{valvules} empêchent le reflux du sang.
\item \textbf{Artères} : le sang quitte le cœur ; \textbf{veines} : il y revient ; \textbf{capillaires} : lieu des échanges.
\item \textbf{Petite circulation} : cœur droit $\rightarrow$ poumons $\rightarrow$ cœur gauche (oxygénation). \textbf{Grande circulation} : cœur gauche $\rightarrow$ organes $\rightarrow$ cœur droit (distribution et collecte).
\item Le cycle cardiaque alterne \textbf{systole} (éjection) et \textbf{diastole} (remplissage).
\item Fréquence cardiaque normale au repos : \textbf{60 à 100 battements/min} ; elle se mesure au poignet (artère radiale), en comptant 30 s et en multipliant par 2.
\end{itemize}
\end{aretenir}

Maintenant que la machinerie est connue, nous pouvons comprendre ce qui se passe quand la pression du sang devient trop forte dans les artères : c'est l'objet de la section suivante, consacrée à la \textbf{tension artérielle} — ce fameux « 17/11 » du supporter de Yaoundé.

\coursec{La tension artérielle et sa mesure}

Dans la section précédente, nous avons découvert l'architecture du système cardiovasculaire : le cœur, cette pompe musculaire, propulse le sang dans un réseau fermé de vaisseaux. Les artères, par leur calibre et leur élasticité, assurent la distribution du sang sous pression vers les organes. C'est précisément cette \textbf{pression exercée par le sang sur la paroi des artères} que l'on appelle la \textbf{tension artérielle (TA)}. Comprendre comment elle s'origine, comment on la mesure et comment l'interpréter est la clé pour dépister le premier facteur de risque des accidents cardiovasculaires : l'hypertension artérielle.

\courssub{Origine physique de la tension artérielle : le rôle de la pompe cardiaque}

\emph{Observation.} Lorsqu'on palpe l'artère radiale (au poignet) ou l'artère carotide (au cou), on perçoit un battement régulier : le pouls. Ce pouls correspond à l'onde de pression générée par l'éjection du sang dans l'aorte à chaque contraction ventriculaire (systole).

\emph{Problème.} La pression dans l'artère n'est pas constante : elle monte brutalement quand le cœur se contracte et redescend quand il se relâche. Comment quantifier ces variations ?

\emph{Hypothèse.} La pression maximale (systolique) correspond à l'éjection du sang ; la pression minimale (diastolique) correspond au relâchement cardiaque, maintenue par l'élasticité des grosses artères (effet « Windkessel » ou réservoir éolien).

\emph{Expérience de référence (historique).} En 1733, Stephen Hales canule l'artère crurale d'un cheval et relie le tuyau à un tube de verre vertical : le sang monte à \SI{2,7}{m} de haut, prouvant l'existence d'une pression artérielle considérable. Aujourd'hui, on utilise un manomètre à mercure (ou son équivalent électronique) qui équilibre la pression artérielle par une pression externe appliquée au bras.

\emph{Interprétation.} Lors de la systole ventriculaire, le ventricule gauche éjecte environ \SI{70}{mL} de sang dans l'aorte. La paroi artérielle se distend : la pression grimpe à son maximum → \textbf{tension artérielle systolique (TAS)} ou « maximale ». Pendant la diastole, la valve aortique se ferme ; le sang stocké dans l'aorte élastique continue de s'écouler vers la périphérie : la pression diminue jusqu'à son minimum → \textbf{tension artérielle diastolique (TAD)} ou « minimale ».

\emph{Conclusion.} La tension artérielle oscille à chaque cycle cardiaque entre une valeur maximale (systolique) et une valeur minimale (diastolique). Elle dépend du \textbf{débit cardiaque} (volume éjecté × fréquence cardiaque) et de la \textbf{résistance vasculaire périphérique} (calibre des artérioles, viscosité du sang, élasticité artérielle).

\begin{definition}[Tension artérielle]
La \cle{tension artérielle (TA)} est la pression exercée par le sang sur la paroi des artères. Elle s'exprime par deux valeurs, notées sous la forme d'une fraction \textbf{systolique / diastolique} (ex. \textbf{12/8}), en \textbf{centimètres de mercure (cmHg)} ou en millimètres de mercure (mmHg) : 1 cmHg = 10 mmHg.
\begin{itemize}
    \item \textbf{Tension systolique (TAS) :} pression maximale lors de la systole ventriculaire (cœur en contraction, éjection du sang).
    \item \textbf{Tension diastolique (TAD) :} pression minimale lors de la diastole ventriculaire (cœur au repos, remplissage), maintenue par le recul élastique des grosses artères.
\end{itemize}
\end{definition}

\begin{popfigure}
\begin{tikzpicture}[scale=0.9, every node/.style={font=\small}]
    % Axes
    \draw[->, thick] (0,0) -- (4.5,0) node[right] {Temps (s)};
    \draw[->, thick] (0,0) -- (0,5.5) node[above] {Pression (cmHg)};
    % Graduations axe Y (0.4 unité = 1 cmHg)
    \foreach \y/\lab in {0.8/0, 1.6/2, 2.4/4, 3.2/6, 4.0/8, 4.8/10, 5.6/12, 6.4/14} {
        \draw (0,\y) -- (-0.1,\y) node[left] {\lab};
    }
    % Courbe de pression artérielle : TAS=12 (y=4.8), TAD=8 (y=3.2), Période=0.8s (75 bpm)
    \draw[popblue, very thick, smooth, domain=0:4, samples=200] plot (\x, {4.0 + 0.8*sin(360*\x/0.8)});
    % Annotations pour un cycle (au 2ème cycle, x=1.2)
    \draw[poporange, thick, <->] (1.2, 3.2) -- (1.2, 4.8) node[midway, right=2mm, poporange] {Amplitude (Pouls) = 4 cmHg};
    \node[popblue] at (1.2, 5.1) {TAS = 12};
    \node[popblue] at (1.2, 2.9) {TAD = 8};
    % Ligne de base diastolique
    \draw[dashed, popmuted] (0,3.2) -- (4,3.2);
    \node[popmuted, left] at (0,3.2) {Diastole (8)};
    % Ligne systolique
    \draw[dashed, popmuted] (0,4.8) -- (4,4.8);
    \node[popmuted, left] at (0,4.8) {Systole (12)};
    % Légende cycle
    \draw[<->, popdark] (0.8, -0.3) -- (1.6, -0.3) node[midway, below] {1 cycle cardiaque ($\approx$ 0,8 s à 75 bpm)};
\end{tikzpicture}
\legende{Variation schématique de la pression artérielle au cours du temps. La courbe oscille entre la tension systolique (TAS = 12 cmHg) et la tension diastolique (TAD = 8 cmHg). L'amplitude (TAS - TAD) est la pression de pouls (4 cmHg). Période = 0,8 s (FC = 75 bpm).}
\end{popfigure}

\courssub{Notation, valeurs de référence et classification}

En pratique clinique, on note la tension sous la forme \textbf{TAS / TAD} (ex. \textbf{12/8}). L'unité officielle au Cameroun (et dans la francophonie médicale) est le \textbf{centimètre de mercure (cmHg)}. Attention : certains tensiomètres électroniques grand public affichent en \textbf{mmHg} (120/80 mmHg = 12/8 cmHg). Ne pas confondre les deux échelles !

\begin{propriete}[Valeurs de référence de la tension artérielle chez l'adulte (au repos)]
\begin{center}
\begin{tabularx}{\linewidth}{|l|X|X|X|}\hline
\rowcolor{popblueL}
\textbf{Catégorie} & \textbf{TAS (cmHg)} & \textbf{TAD (cmHg)} & \textbf{Conduite à tenir} \\ \hline
\textbf{Optimale} & $<$ 12 & $<$ 8 & Surveillance routine \\ \hline
\textbf{Normale} & 12 -- 12,9 & 8 -- 8,4 & Surveillance routine \\ \hline
\textbf{Normale haute} & 13 -- 13,9 & 8,5 -- 8,9 & Conseil hygiéno-diététique, contrôle à 1 an \\ \hline
\textbf{Hypertension artérielle (HTA) Grade 1} & 14 -- 15,9 & 9 -- 9,9 & Mesures répétées, modification mode de vie, avis médical \\ \hline
\textbf{Hypertension artérielle (HTA) Grade 2} & 16 -- 17,9 & 10 -- 10,9 & Prise en charge médicale + hygiène de vie \\ \hline
\textbf{Hypertension artérielle (HTA) Grade 3} & $\ge$ 18 & $\ge$ 11 & Prise en charge médicale urgente \\ \hline
\end{tabularx}
\end{center}
\emph{Seuil diagnostique de l'HTA :} \textbf{TAS $\ge$ 14 cmHg \textbf{et/ou} TAD $\ge$ 9 cmHg}, confirmé par des mesures répétées.
\end{propriete}

\begin{attention}[Piège : les unités et le « 12/8 »]
\begin{itemize}
    \item \textbf{cmHg vs mmHg :} 12/8 cmHg = 120/80 mmHg. Si l'appareil affiche 120/80, c'est en mmHg. Ne pas écrire « 120/80 cmHg » (ce serait une tension incompatible avec la vie !).
    \item \textbf{Seuil 14/9 :} C'est le seuil \textbf{au cabinet médical}. En automesure (domicile), le seuil est souvent abaissé à \textbf{13,5/8,5} car l'effet « blouse blanche » (stress du médecin) est absent.
    \item \textbf{Une seule mesure ne suffit JAMAIS} pour diagnostiquer une HTA (voir boîte \texttt{Attention} plus bas).
\end{itemize}
\end{attention}

\courssub{L'appareil de mesure : le tensiomètre}

Deux technologies coexistent dans les formations sanitaires et les foyers camerounais.

\begin{experience}[Le tensiomètre à mercure (auscultatoire) — Référence historique]
\textbf{Matériel :} Brassard (manchette) gonflable, poire de gonflage avec valve de dégonflage, colonne de mercure (manomètre), stéthoscope.
\textbf{Principe (Méthode de Korotkoff) :}
\begin{enumerate}
    \item On place le brassard au bras (niveau du cœur), stéthoscope sur l'artère brachiale (creux du coude).
    \item On gonfle le brassard jusqu'à \textbf{occlure totalement l'artère} (pression $>$ TAS, ex. 180 mmHg) : plus de bruit.
    \item On dégonfle \textbf{lentement} ($\approx$ 2--3 mmHg/s).
    \item \textbf{Premier bruit clair (ton Korotkoff I)} $\rightarrow$ la pression artérielle vainc la pression du brassard : \textbf{TAS}.
    \item Les bruits s'atténuent, deviennent sourds, puis \textbf{disparaissent complètement (ton Korotkoff V)} $\rightarrow$ \textbf{TAD}.
\end{enumerate}
\textbf{Avantage :} Étalon de référence, pas de pile. \textbf{Inconvénient :} Mercure toxique, fragile, nécessite formation (auscultation), effet « blouse blanche » majoré.
\end{experience}

\begin{experience}[Le tensiomètre électronique (oscillométrique) — Usage courant]
\textbf{Principe :} Le brassard se gonfle/dégonfle automatiquement. Un capteur détecte les \textbf{oscillations de pression} dans le brassard transmises par la paroi artérielle. L'amplitude maximale des oscillations correspond à la pression moyenne ; des algorithmes estiment TAS et TAD.
\textbf{Avantage :} Simple, autonome, élimine l'erreur d'auscultation, permet l'automésure à domicile.
\textbf{Inconvénient :} Moins fiable en cas d'arythmie (fibrillation auriculaire), de bras très conique, ou d'appareil non validé (norme ISO 81060-2).
\end{experience}

\begin{popfigure}
\begin{tikzpicture}[scale=0.85, every node/.style={font=\small}]
    % Brassard
    \draw[popblue, line width=3mm, rounded corners=3mm] (0,0) rectangle (8,2.5);
    \draw[popblue!50, line width=1mm] (0.5,0.5) rectangle (7.5,2.0);
    \node[popblue, rotate=90] at (-0.5,1.25) {BRAS};
    % Tuyau vers manomètre
    \draw[poporange, line width=2mm] (8,1.25) -- (9.5,1.25);
    % Manomètre à mercure (schématique)
    \draw[popdark, thick] (9.5, 0.2) rectangle (11.5, 5.5);
    % Mercure
    \fill[poporange!80!popdark] (9.7, 0.4) rectangle (11.3, 3.8);
    % Graduations
    \foreach \y/\lab in {0.4/0, 1.2/2, 2.0/4, 2.8/6, 3.6/8, 4.4/10, 5.2/12} {
        \draw[popdark] (11.3, \y) -- (11.8, \y) node[right] {\lab};
    }
    \node[popdark, above] at (10.5, 5.5) {cmHg};
    % Flèches de lecture
    \draw[popgreen, thick, <->] (10.5, 0.4) -- (10.5, 3.8) node[midway, right=2mm, popgreen] {Lecture : 12 / 8};
    % Poire
    \draw[poppurple, thick] (8,1.25) .. controls (7,3) and (5,3) .. (5,1.25);
    \node[poppurple, above] at (6, 3.2) {Poire de gonflage};
    % Stéthoscope
    \draw[poppink, thick] (4, -0.5) -- (4, -1.5) -- (2, -1.5);
    \draw[poppink, thick] (4, -0.5) -- (4, -1.5) -- (6, -1.5);
    \node[poppink, below] at (4, -1.8) {Stéthoscope (artère brachiale)};
    % Légende valeurs
    \node[align=center, popdark] at (4, 3.5) {Tensiomètre à mercure (auscultatoire)\\\textbf{TAS = 12 cmHg} --- \textbf{TAD = 8 cmHg}};
\end{tikzpicture}
\legende{Schéma de principe d'un tensiomètre à mercure (manuel) avec brassard, colonne de mercure, poire de gonflage et stéthoscope. La lecture se fait au niveau du ménisque du mercure au moment des bruits de Korotkoff I (systolique) et V (diastolique).}
\end{popfigure}

\courssub{Conditions d'une mesure fiable : le protocole standardisé}

Une mesure de tension n'est valide que si elle respecte des conditions strictes, sinon elle surestime (souvent) ou sous-estime la vraie pression de repos.

\begin{methode}[Protocole de mesure de la tension artérielle (recommandations OMS / Société Camerounaise de Cardiologie)]
\begin{enumerate}
    \item \textbf{Repos :} Le patient est assis, calme, depuis \textbf{au moins 5 minutes} (pas de marche, pas d'escalier, pas de stress immédiat).
    \item \textbf{Position :} Assis, dos appuyé, \textbf{bras nu} (pas de vêtement serré), avant-bras soutenu sur une table, \textbf{paume vers le haut}, \textbf{bras au niveau du cœur} (niveau du 4ème espace intercostal, sternum).
    \item \textbf{Abstentions (30 min avant) :} Pas de café, thé, cola, alcool, tabac, effort physique intense.
    \item \textbf{Vessie vide :} Une vessie pleine peut majorer la TA de 10--15 mmHg (1--1,5 cmHg).
    \item \textbf{Brassard adapté :} Largeur = 40\% de la circonférence du bras (ou 2/3 de la longueur du bras). Brassard trop petit $\rightarrow$ surestimation (fausse HTA) ; trop grand $\rightarrow$ sous-estimation.
    \item \textbf{Mesures répétées :} \textbf{Au moins 2 mesures} espacées de 1--2 minutes. Si écart $>$ 5 mmHg (0,5 cmHg) $\rightarrow$ 3ème mesure. On retient la \textbf{moyenne des 2 dernières}.
    \item \textbf{Les deux bras :} À la première consultation, mesurer aux deux bras. Si différence $>$ 10 mmHg (1 cmHg) $\rightarrow$ pathologie vasculaire possible ; on retient le bras le plus haut pour le suivi.
\end{enumerate}
\end{methode}

\begin{attention}[Erreurs fréquentes en TP / consultation]
\begin{itemize}
    \item \textbf{Bras pendouillant} (sous niveau du cœur) $\rightarrow$ surestimation par effet hydrostatique ($\approx$ +0,8 cmHg par 10 cm sous le cœur).
    \item \textbf{Parler ou bouger} pendant la mesure $\rightarrow$ élévation artefacts.
    \item \textbf{Brassard sur vêtement épais} $\rightarrow$ mesure faussée.
    \item \textbf{Dégonflage trop rapide} (manuel) $\rightarrow$ sous-estimation de la TAS, surestimation de la TAD.
    \item \textbf{Arrondir au cmHg supérieur} : on note la valeur lue (ex. 13,6 $\rightarrow$ 13,6 ou 13,5 selon graduation), pas 14.
\end{itemize}
\end{attention}

\courssub{Interprétation de mesures simples : démarche d'investigation guidée}

\emph{Situation concrète.} Au dispensaire de Bafia, l'infirmière mesure la TA de trois patients au repos, selon le protocole. Voici les moyennes des deux dernières mesures :

\begin{center}
\begin{tabular}{|c|c|c|}\hline
\textbf{Patient} & \textbf{TAS (cmHg)} & \textbf{TAD (cmHg)} \\ \hline
Mme Atangana & 11 & 7 \\ \hline
M. Essono & 16 & 10 \\ \hline
Mme Ngo & 14 & 9 \\ \hline
\end{tabular}
\end{center}

\emph{Question.} Classer chaque patient selon la classification OMS et proposer une conduite à tenir.

\emph{Analyse étape par étape.}
\begin{enumerate}
    \item \textbf{Mme Atangana (11/7) :} TAS $<$ 12 \textbf{et} TAD $<$ 8 $\rightarrow$ \textbf{Optimale}. \emph{Conduite :} Encourager le mode de vie sain, contrôle dans 2--3 ans.
    \item \textbf{M. Essono (16/10) :} TAS $\ge$ 14 \textbf{et} TAD $\ge$ 9 $\rightarrow$ \textbf{HTA Grade 2}. \emph{Conduite :} \textbf{Confirmation impérative} (automesure sur 3--7 jours ou MAPA), consultation médicale pour bilan (recherche cause, atteinte d'organes) + mesures hygiéno-diététiques immédiates.
    \item \textbf{Mme Ngo (14/9) :} TAS = 14 \textbf{et} TAD = 9 $\rightarrow$ \textbf{HTA Grade 1 (seuil)}. \emph{Conduite :} Ne pas traiter sur une seule visite ! Prescrire automesure domicile (règle des 3 : 3 matins, 3 soirs, 3 jours) ou MAPA (Mesure Ambulatoire de la Pression Artérielle). Si confirmée $\rightarrow$ hygiène de vie + réévaluation à 3 mois.
\end{enumerate}

\begin{exemplebox}[Interprétation clinique : le cas de M. Essono]
M. Essono, 52 ans, cultivateur à Ebolowa, se présente pour un malaise. TA mesurée trois fois à 5 min d'intervalle : \textbf{16/10}, \textbf{15/10}, \textbf{16/11}. Moyenne = \textbf{15,7/10,3} $\approx$ \textbf{16/10}.
\begin{itemize}
    \item \textbf{Classification :} HTA Grade 2 (TAS 16--17,9 et/ou TAD 10--10,9).
    \item \textbf{Conduite immédiate :} Repos, interdiction sel ajouté, activité physique modérée, rendez-vous médecin dans la semaine pour bilan (créatinine, ECBU, ECG, fond d'œil, glycémie, lipidogramme).
    \item \textbf{Message clé :} « Monsieur, votre tension est trop haute (16/10). Ce n'est pas une urgence vitale aujourd'hui, mais si elle reste ainsi, elle abîme votre cœur, vos reins et votre cerveau. Nous allons vérifier chez vous sur plusieurs jours et le médecin décidera d'un traitement si besoin. »
\end{itemize}
\end{exemplebox}

\begin{exoresolu}[Exercice d'interprétation de mesures]
\textbf{Énoncé.} Lors d'une campagne de dépistage au marché de Mfoundi, on obtient les mesures suivantes (moyenne de 2 mesures au repos, brassard adapté) :
\begin{itemize}
    \item A. 13,5 / 8,2
    \item B. 14,2 / 8,8
    \item C. 12,8 / 9,2
    \item D. 17,5 / 11,0
\end{itemize}
Pour chaque cas, donner la catégorie OMS et la conduite à tenir prioritaire.

\tcblower
\textbf{Solution détaillée.}
\begin{itemize}
    \item \textbf{A (13,5/8,2) :} TAS 13--13,9, TAD 8--8,4 $\rightarrow$ \textbf{Normale haute}. \emph{Conduite :} Conseil hygiéno-diététique (sel, poids, activité), contrôle à 1 an.
    \item \textbf{B (14,2/8,8) :} TAS $\ge$ 14, TAD $<$ 9 $\rightarrow$ \textbf{HTA Grade 1 (systolique isolée)}. \emph{Conduite :} Confirmation par automesure/MAPA. Si confirmée : hygiène de vie + réévaluation 3 mois.
    \item \textbf{C (12,8/9,2) :} TAS $<$ 14, TAD $\ge$ 9 $\rightarrow$ \textbf{HTA Grade 1 (diastolique isolée)}. \emph{Conduite :} Idem B. Fréquente chez le sujet jeune (hyperkinétique).
    \item \textbf{D (17,5/11,0) :} TAS $\ge$ 16, TAD $\ge$ 10 $\rightarrow$ \textbf{HTA Grade 2}. \emph{Conduite :} Confirmation rapide (automesure 3 jours), consultation médicale dans la semaine pour bilan et probable traitement médicamenteux + hygiène de vie.
\end{itemize}
\emph{Rappel :} Le diagnostic d'HTA ne se pose \textbf{jamais} sur une seule séance de mesure (sauf urgence hypertensive avec signes de gravité : encéphalopathie, œdème pulmonaire aigu, dissection aortique, TAS $\ge$ 18 \textbf{et} TAD $\ge$ 11 avec symptômes).
\end{exoresolu}

\courssub{Visualisation comparative : tensions normales et pathologiques}

Le graphique ci-dessous compare quatre profils tensionnels types. Les barres bleues représentent la TAS, les barres oranges la TAD. La zone rouge hachurée marque le seuil d'hypertension (14/9).

\begin{popfigure}
\begin{tikzpicture}
\begin{axis}[
    ybar,
    bar width=12pt,
    width=\linewidth,
    height=8cm,
    ylabel={Tension (cmHg)},
    ymin=0, ymax=20,
    xtick={1,2,3,4},
    xticklabels={11/7 (Optimale), 12/8 (Normale), 14/9 (Seuil HTA), 17/11 (HTA Gr2)},
    x tick label style={rotate=30, anchor=east, font=\small},
    ytick={0,2,4,6,8,10,12,14,16,18,20},
    legend style={at={(0.5,-0.25)}, anchor=north, legend columns=2, font=\small},
    grid=major, grid style={popline},
    enlargelimits=0.15,
]
% Seuil HTA : lignes horizontales
\addplot[popred!50, dashed, thick, forget plot] coordinates {(0,14) (5,14)};
\addplot[popred!50, dashed, thick, forget plot] coordinates {(0,9) (5,9)};
% Zone seuil (hachurée)
\addplot[popred!15, fill=popred!15, forget plot, mark=none] coordinates {(0,14) (5,14) (5,20) (0,20)} \closedcycle;
\addplot[popred!15, fill=popred!15, pattern=north east lines, pattern color=popred!50, forget plot, mark=none] coordinates {(0,9) (5,9) (5,14) (0,14)} \closedcycle;
% Données TAS (décalées légèrement à gauche de chaque catégorie)
\addplot[popblue, fill=popblue!70] coordinates {(0.85,11) (1.85,12) (2.85,14) (3.85,17)};
% Données TAD (décalées légèrement à droite de chaque catégorie)
\addplot[poporange, fill=poporange!70] coordinates {(1.15,7) (2.15,8) (3.15,9) (4.15,11)};
\legend{TAS (Systolique), TAD (Diastolique)}
\end{axis}
\end{tikzpicture}
\legende{Comparaison de profils tensionnels. Les barres bleues = TAS, oranges = TAD. La zone hachurée rouge délimite la zone d'hypertension artérielle (TAS $\ge$ 14 \textbf{et/ou} TAD $\ge$ 9). Seuls les deux premiers profils sont normaux.}
\end{popfigure}

\begin{savaistu}[L'hypertension au Cameroun : le « tueur silencieux »]
Au Cameroun, les enquêtes STEPS (OMS) et les études locales (ex. \textit{Kingue et al., 2015}) estiment la prévalence de l'HTA à \textbf{25--35\%} des adultes (environ \textbf{1 adulte sur 3} après 25 ans).
\begin{itemize}
    \item \textbf{Moins de 50\%} des hypertendus se savent malades (asymptomatique longtemps).
    \item \textbf{Moins de 30\%} des traités sont contrôlés (TA $<$ 14/9).
    \item L'HTA est la 1ère cause de consultation en cardiologie à l'Hôpital Général de Yaoundé et à Laquintinie (Douala).
    \item \textbf{Conséquences majeures :} AVC (1ère cause de décès neurologique), insuffisance cardiaque, insuffisance rénale terminale (dialyse), rétinopathie hypertensive.
\end{itemize}
D'où l'importance cruciale du \textbf{dépistage systématique} (toute occasion : vaccination, CPN, consultation banale) et de l'\textbf{éducation thérapeutique} (observance, automesure, hygiène de vie).
\end{savaistu}

\begin{aretenir}[Essentiel à retenir sur la tension artérielle]
\begin{itemize}
    \item La TA = pression du sang sur les artères. Notée \textbf{TAS / TAD} en \textbf{cmHg} (ex. 12/8).
    \item \textbf{TAS} = pression à la systole (éjection) ; \textbf{TAD} = pression à la diastole (recul élastique artères).
    \item \textbf{Seuil HTA :} $\ge$ \textbf{14/9} (cabinet), confirmé sur \textbf{mesures répétées} (au moins 2 visites ou automesure 3j).
    \item \textbf{Mesure fiable :} repos 5 min, assis, bras nu au niveau du cœur, pas de café/tabac 30 min avant, brassard adapté, moyenne de 2 mesures.
    \item \textbf{Interprétation :} comparer TAS et TAD aux seuils (Optimale $<$12/8 ; Normale 12--12,9/8--8,4 ; Normale haute 13--13,9/8,5--8,9 ; HTA Grade 1 14--15,9/9--9,9 ; Grade 2 16--17,9/10--10,9 ; Grade 3 $\ge$18/$\ge$11).
    \item \textbf{Au Cameroun :} 1 adulte sur 3 hypertendu, souvent ignoré $\rightarrow$ dépistage et prévention prioritaires.
\end{itemize}
\end{aretenir}

\begin{exercice}[Entraînement : Mesure et interprétation]
\begin{enumerate}
    \item Un tensiomètre électronique affiche \textbf{138/86 mmHg}. Convertir en cmHg et classer selon l'OMS.
    \item Lors d'une mesure manuelle, on entend le premier bruit de Korotkoff à 15 cmHg et le dernier à 9 cmHg. Quelle est la TA ? Quelle catégorie ?
    \item Pourquoi est-il impératif d'avoir le bras \textbf{au niveau du cœur} ? Que se passe-t-il si le bras est 20 cm plus bas ? (Indication : $\Delta P = \rho g h$).
    \item M. K., 45 ans, présente 14,5/9,2 à la consultation. Quelle conduite proposer selon les recommandations ?
\end{enumerate}
\end{exercice}

\begin{corrige}[Exercice : Mesure et interprétation]
\begin{enumerate}
    \item 138/86 mmHg = \textbf{13,8/8,6 cmHg}. TAS 13--13,9, TAD 8,5--8,9 $\rightarrow$ \textbf{Normale haute}. Conseil hygiéno-diététique, contrôle 1 an.
    \item TA = \textbf{15/9 cmHg}. TAS $\ge$ 14, TAD $\ge$ 9 $\rightarrow$ \textbf{HTA Grade 1}. Nécessite confirmation (automesure/MAPA) avant tout diagnostic définitif.
    \item Pression hydrostatique : $\Delta P = \rho g h$. Sang : $\rho \approx 1060$ kg/m$^3$, $g=9,81$ m/s$^2$, $h=0,20$ m $\rightarrow$ $\Delta P \approx 2080$ Pa $\approx$ 15,6 mmHg $\approx$ \textbf{1,6 cmHg}. Si le bras est 20 cm \textbf{sous} le cœur, la mesure surestime la vraie TA d'environ \textbf{1,6 cmHg} (risque de fausse HTA).
    \item TA = 14,5/9,2 $\rightarrow$ \textbf{HTA Grade 1}. Conduite : \textbf{Pas de traitement immédiat}. Prescrire automesure domicile (règle des 3 : 3 matins, 3 soirs, 3 jours) ou MAPA. Si moyenne domicile $\ge$ 13,5/8,5 confirmée $\rightarrow$ mesures hygiéno-diététiques (sel $<$ 5g/j, DASH, activité, poids, alcool) + réévaluation à 3 mois. Si non contrôlée $\rightarrow$ traitement médicamenteux.
\end{enumerate}
\end{corrige}

\coursec{Les principaux accidents cardiovasculaires}

\noindent Dans la section précédente, nous avons appris à mesurer la tension artérielle et à comprendre son sens physiologique : elle reflète la force exercée par le sang sur les parois artérielles. Une tension normale assure une perfusion efficace de tous les organes. Mais que se passe-t-il lorsque ce système de distribution se grippe ? Lorsque les artères se bouchent ou se rompent, ou lorsque la pression devient durablement trop forte ? C'est l'objet de cette section : comprendre les \textbf{accidents cardiovasculaires}, ces événements brutaux ou chroniques qui menacent le cœur et le cerveau, et savoir réagir face à une victime.

\courssub{L'infarctus du myocarde : quand le cœur manque d'oxygène}

\begin{experience}[Observation : un cas d'infarctus à l'hôpital de district de Bafoussam]
\textbf{Situation :} M. K., 58 ans, cultivateur, est admis aux urgences pour une \textbf{douleur thoracique intense, constrictive, survenant au repos} depuis 45 minutes. La douleur irradie vers le \textbf{bras gauche} (face interne) et la \textbf{mâchoire}. Il est pâle, en sueur froide (sudation profuse), anxieux, avec une \textbf{fréquence cardiaque} à \SI{110}{bpm} et une \textbf{tension artérielle} à \SI{90}{mmHg} (systolique) / \SI{60}{mmHg} (hypotension). L'ECG montre un \textbf{sus-décalage du segment ST} dans les dérivations antérieures (V1--V4). Le dosage des \textbf{troponines} (protéines contractiles cardiaques spécifiques : TnT ou TnI) est très élevé.
\textbf{Question :} Quel est le mécanisme lésant le muscle cardiaque ? Pourquoi la douleur irradie-t-elle au bras et à la mâchoire ?
\end{experience}

\begin{definition}[Infarctus du myocarde]
L'\textbf{infarctus du myocarde} (ou crise cardiaque) est la \textbf{nécrose (mort) irréversible} d'une partie du muscle cardiaque (myocarde) due à l'\textbf{obstruction brutale et complète} d'une artère coronaire (le plus souvent par un thrombus sur une plaque d'athérome), entraînant l'arrêt de l'apport en oxygène et en nutriments.
\end{definition}

\noindent \textbf{Mécanisme physiopathologique (démarche d'investigation) :}
\begin{enumerate}
    \item \textbf{Observation :} L'angio-coronarographie (radiographie des artères du cœur après injection de produit de contraste iodé) montre une \textbf{sténose (rétrécissement) critique} ou une \textbf{occlusion totale} d'une artère coronaire (souvent l'interventriculaire antérieure, branche de la coronaire gauche).
    \item \textbf{Hypothèse :} Une plaque d'\textbf{athérome} (dépôt de cholestérol LDL oxydé, cellules inflammatoires -- macrophages/mousseux, tissu fibreux) s'est fissurée ou érodée, exposant son contenu thrombogène (facteur tissulaire, collagène) au sang. Cela déclenche l'activation plaquettaire et la coagulation : formation d'un \textbf{thrombus} (caillot blanc plaquettaire / rouge fibrinique) qui bouche l'artère.
    \item \textbf{Expérience / Données :} En l'absence de reperfusion (réouverture de l'artère par thrombolyse ou angioplastie primaire) dans les \textbf{premières 3 à 6 heures} (fenêtre thérapeutique optimale : \emph{« golden hours »}), les cardiomyocytes (cellules musculaires cardiaques) meurent par nécrose de coagulation (ischémie prolongée $\rightarrow$ défaillance mitochondriale $\rightarrow$ rupture membranaire). La zone nécrosée est remplacée par une \textbf{cicatrice fibreuse} (tissu conjonctif dense, non contractile) en 4 à 6 semaines (remodelage ventriculaire).
    \item \textbf{Interprétation :} Le myocarde en aval de l'obstruction passe par trois stades : \textbf{ischémie} (manque d'oxygène, réversible si bref) $\rightarrow$ \textbf{souffrance / lésion} (douleur, troubles du rythme, sus-décalage ST, élévation troponines) $\rightarrow$ \textbf{nécrose} (mort cellulaire définitive, onde Q pathologique à l'ECG ultérieur).
    \item \textbf{Conclusion :} L'infarctus est une \textbf{urgence vitale absolue} : « \emph{Time is muscle} » (le temps, c'est du muscle sauvé). La reperfusion rapide (thrombolyse si délai $<$ 3h et pas de centre d'angioplastie, sinon angioplastie primaire avec pose de stent) limite l'étendue de la nécrose et préserve la fonction ventriculaire gauche.
\end{enumerate}

\begin{popfigure}
\begin{tikzpicture}[scale=1.1, every node/.style={font=\small}]
    % Arterie saine
    \begin{scope}[xshift=-5.5cm]
        \draw[popblue, line width=2pt] (0,0) -- (3,0); % Lumen
        \draw[popblue!50, line width=1pt] (0,-0.5) -- (3,-0.5); % Intima
        \draw[popblue!50, line width=1pt] (0,0.5) -- (3,0.5);
        \draw[popdark, line width=1.5pt] (0,-0.8) -- (3,-0.8); % Media
        \draw[popdark, line width=1.5pt] (0,0.8) -- (3,0.8);
        \draw[popmuted, line width=1pt] (0,-1.1) -- (3,-1.1); % Adventice
        \draw[popmuted, line width=1pt] (0,1.1) -- (3,1.1);
        
        \node[popblue, align=center] at (1.5, 1.6) {\textbf{Artère coronaire saine}};
        \draw[<-, popblue] (1.5, 1.3) -- (1.5, 1.1);
        \node[popblue, align=center, anchor=west] at (3.2, 0) {Lumen large\\(flux sanguin optimal)};
        \draw[<-, popblue] (3.1, 0.1) -- (3, 0.5);
        \node[popdark, align=center, anchor=west] at (3.2, -0.5) {Média épaisse\\(muscle lisse, élastine)};
        \draw[<-, popdark] (3.1, -0.4) -- (3, -0.2);
    \end{scope}

    % Arterie obstruee
    \begin{scope}[xshift=5.5cm]
        \draw[poporange, line width=2pt] (0,0) -- (3,0); % Lumen residual
        \draw[poporange!50, line width=1pt] (0,-0.5) -- (3,-0.5);
        \draw[poporange!50, line width=1pt] (0,0.5) -- (3,0.5);
        \draw[popdark, line width=1.5pt] (0,-0.8) -- (3,-0.8);
        \draw[popdark, line width=1.5pt] (0,0.8) -- (3,0.8);
        \draw[popmuted, line width=1pt] (0,-1.1) -- (3,-1.1);
        \draw[popmuted, line width=1pt] (0,1.1) -- (3,1.1);
        
        % Plaque d'atherome
        \fill[poporange!80!popdark] (1.0, -0.5) .. controls (1.5, -0.2) and (2.0, -0.2) .. (2.5, -0.5) -- (2.5, 0) .. controls (2.0, 0.2) and (1.5, 0.2) .. (1.0, 0) -- cycle;
        % Thrombus
        \fill[popred!80!popdark] (1.2, 0) .. controls (1.6, 0.3) and (2.1, 0.3) .. (2.3, 0) -- (2.3, -0.1) .. controls (2.0, 0.1) and (1.5, 0.1) .. (1.2, -0.1) -- cycle;
        
        \node[poporange, align=center] at (1.5, 1.6) {\textbf{Artère coronaire obstruée (Infarctus)}};
        \draw[<-, poporange] (1.5, 1.3) -- (1.5, 1.1);
        \node[popred, align=center, anchor=west] at (3.2, 0.3) {Thrombus (caillot)\\(obstruction brutale)};
        \draw[<-, popred] (3.1, 0.25) -- (2.5, 0.15);
        \node[poporange!80!popdark, align=center, anchor=west] at (3.2, -0.3) {Plaque d'athérome\\(fissurée, lipidique)};
        \draw[<-, poporange!80!popdark] (3.1, -0.25) -- (2.5, -0.4);
        \node[poporange, align=center, anchor=west] at (3.2, -0.8) {Lumen rétréci/occlus\\(ischémie / nécrose)};
        \draw[<-, poporange] (3.1, -0.75) -- (3, -0.4);
    \end{scope}
\end{tikzpicture}
\legende{Comparaison histologique schématique : à gauche, une artère coronaire saine à paroi régulière (intima fine, média musculaire épaisse, adventice) et lumen large ; à droite, une artère obstruée par une plaque d'athérome fissurée (jaune/orange, noyau lipidique, cap fibreux rompu) surmontée d'un thrombus occlusif (rouge, riche en plaquettes/fibrine), responsable de l'infarctus du myocarde en aval.}
\end{popfigure}

\noindent \textbf{Signes cliniques d'alerte (à connaître par cœur pour le Probatoire) :}
\begin{itemize}
    \item \textbf{Douleur thoracique} : rétrosternale, de type \textbf{constrictif, en étau, en barre}, survenant le plus souvent \textbf{au repos} (ou à l'effort mais non cédée par l'arrêt), d'une durée \textbf{$>$ 20 min}, \textbf{non cédée par le repos} ni la trinitrine (différence majeure avec l'angine de poitrine stable).
    \item \textbf{Irradiation} : vers le \textbf{bras gauche} (face interne, jusqu'au petit doigt), la \textbf{mâchoire} (dents), le dos (inter-scapulaire), l'épaule droite, l'épigastre (peut simuler une pathologie digestive).
    \item \textbf{Signes d'accompagnement} : \textbf{sueurs froides} (sudation profuse), \textbf{pâleur}, \textbf{angoisse} (sentiment de mort imminente), \textbf{essoufflement} (dyspnée), nausées/vomissements (surtout infarctus inférieur).
    \item \textbf{Signes de gravité} : hypotension (choc cardiogénique : TA systolique $<$ 90 mmHg), troubles du rythme graves (fibrillation ventriculaire = arrêt cardiaque soudain, blocs AV), insuffisance cardiaque aiguë (œdème pulmonaire aigu : râles crépitants, expectorations mousseuses rosées).
\end{itemize}

\begin{exemplebox}[Cas concret : M. Essomba, 62 ans, Yaoundé]
M. Essomba, ancien fonctionnaire, diabétique de type 2 non équilibré (HbA1c > 9\%) et fumeur (20 PA), ressent vers 3h du matin une « barre sur la poitrine » qui lui « prend la gorge » et lui « serre le bras gauche ». Il attend 2h en pensant à une indigestion. Sa femme appelle le \textbf{119 (SAMU Cameroun)}. À l'arrivée des secours : FC \SI{105}{bpm}, TA \SI{100}{mmHg}/\SI{70}{mmHg}, sueurs froides. ECG : sus-décalage ST V1-V4 (infarctus antérieur). Prise en charge pré-hospitalière : \textbf{aspirine \SI{300}{mg} à croquer} (antiagrégant plaquettaire), oxygène si SpO2 $<$ 94\%, morphine (antalgique), monitoring continu, transfert en unité de soins intensifs cardiologiques (USIC) pour \textbf{angioplastie primaire} (dilatation + stent).
\emph{Leçon :} Ne jamais attendre face à une douleur thoracique atypique ou prolongée chez un sujet à risque. \textbf{Appeler le 119 immédiatement.} L'automédication ou l'attente retarde la reperfusion.
\end{exemplebox}

\courssub{L'accident vasculaire cérébral (AVC) : urgence cérébrale}

\begin{experience}[Observation : deux types d'AVC au CHU de Yaoundé]
\textbf{Cas 1 (Mme A., 72 ans, HTA non traitée) :} Brutalement, pendant le repas, \textbf{paralysie du bras et de la jambe droits} (hémiplégie droite), \textbf{bouche déviée à gauche} (atteinte faciale controlatérale), \textbf{impossibilité de parler} (aphasie de Broca : expression altérée, compréhension conservée). Scanner cérébral (TDM sans injection) : \textbf{hypodensité} (zone noire, substance blanche) dans le territoire de l'artère cérébrale moyenne (ACM) gauche, \textbf{sans hypersignal hémorragique}. $\rightarrow$ \textbf{AVC ischémique} (infarctus cérébral par thrombose sur athérome carotidien ou embolie cardiaque).
\textbf{Cas 2 (M. B., 55 ans, HTA sévère non observée, effort physique) :} \textbf{Céphalée brutale violente} (« coup de massue », céphalée tonnerre), \textbf{vomissements en jet} (signe d'hypertension intracrânienne), \textbf{perte de connaissance} rapide (score de Glasgow bas), \textbf{paralysie gauche} (hémiplégie controlatérale à l'hématome). Scanner : \textbf{hypersignal blanc} (sang hyperdense) dans les noyaux de la base droits (putamen/thalamus, siège des artères lenticulostriées), avec \textbf{effet de masse} (déviation de la ligne médiane, compression des ventricules). $\rightarrow$ \textbf{AVC hémorragique} (hématome intra-parenchymateux hypertensif).
\textbf{Question :} Quels sont les deux mécanismes radicalement opposés ? Pourquoi le scanner (TDM) est-il indispensable avant tout traitement spécifique ?
\end{experience}

\begin{definition}[Accident Vasculaire Cérébral (AVC)]
L'\textbf{AVC} (ou « attaque ») est une \textbf{atteinte brutale de la fonction neurologique} (installation en quelques minutes) due à une \textbf{perturbation soudaine de la circulation sanguine cérébrale}. On distingue deux types majeurs aux pronostics et traitements opposés :
\begin{itemize}
    \item \textbf{AVC ischémique (environ 80\% des cas au Cameroun comme ailleurs) :} \textbf{Obstruction} d'une artère cérébrale (thrombose locale sur plaque d'athérome carotidien ou intracrânien, ou embolie d'origine cardiaque -- fibrillation atrielle, valvulopathie, cardiopathie ischémique) $\rightarrow$ infarctus cérébral (nécrose du tissu neuronal par anoxie/aglycie).
    \item \textbf{AVC hémorragique (environ 20\% des cas, mais plus grave) :} \textbf{Rupture} d'une artère cérébrale (souvent sur HTA mal contrôlée $\rightarrow$ micro-anévrismes de Charcot-Bouchard sur les artères perforantes profondes ; ou rupture d'anévrisme sacciforme sur le cercle de Willis $\rightarrow$ hémorragie méningée) $\rightarrow$ hématome intra-parenchymateux (dans la substance cérébrale) ou hémorragie méningée (dans l'espace sous-arachnoïdien).
\end{itemize}
\end{definition}

\begin{popfigure}
\begin{tikzpicture}[scale=0.9, every node/.style={font=\small}]
    % Cerveau gauche : AVC Ischémique
    \begin{scope}[xshift=-5.5cm]
        % Hémisphère gauche
        \draw[popblue, line width=1.5pt] (0,0) .. controls (0,2.5) and (2,3.5) .. (3,3) .. controls (4,2.5) and (4,1) .. (3.5,0) .. controls (2.5,-0.5) and (1,-0.5) .. (0,0);
        % Artère cérébrale moyenne (ACM) - trunk et branches
        \draw[popblue, line width=1.2pt] (1.5,0.5) .. controls (2,1.5) and (2.5,2) .. (2.8,2.5);
        \draw[popblue, line width=0.8pt] (2,1.5) .. controls (2.2,1.8) and (2.5,2.2) .. (2.8,2.5); % branche sup
        \draw[popblue, line width=0.8pt] (2,1.5) .. controls (1.8,1.2) and (1.5,0.8) .. (1.2,0.5); % branche inf
        % Thrombus sur tronc ACM
        \fill[popred!80!popdark] (2.2, 1.7) circle (0.15);
        \draw[<-, popred, thick] (2.35, 1.85) -- (3.5, 2.2) node[right, align=left, popred] {Thrombus\\obstruant l'ACM};
        % Zone d'infarctus (ischémie/nécrose)
        \fill[popblue!20!white, opacity=0.7] (1.8,1.2) .. controls (2.5,1.8) and (3,2.2) .. (3.2,2.5) .. controls (3,2) and (2.5,1.5) .. (1.8,1.2);
        \node[popblue, align=center] at (1.5, -0.7) {\textbf{AVC Ischémique}};
        \node[popblue, align=center, font=\tiny] at (1.5, -1.1) {Obstruction artérielle $\rightarrow$ Ischémie $\rightarrow$ Nécrose};
    \end{scope}

    % Cerveau droit : AVC Hémorragique
    \begin{scope}[xshift=5.5cm]
        % Hémisphère droit
        \draw[poporange, line width=1.5pt] (0,0) .. controls (0,2.5) and (-2,3.5) .. (-3,3) .. controls (-4,2.5) and (-4,1) .. (-3.5,0) .. controls (-2.5,-0.5) and (-1,-0.5) .. (0,0);
        % Artère lenticulostriée (profonde, perforante)
        \draw[poporange, line width=1pt] (-1.5,0.5) -- (-1.5,2.2);
        % Rupture / Hématome
        \fill[popred!90!popdark] (-1.5, 1.5) circle (0.4);
        \draw[popred, thick] (-1.5, 1.5) circle (0.55); % effet de masse (contour élargi)
        \draw[<-, popred, thick] (-1.0, 1.9) -- (-0.2, 2.4) node[right, align=left, popred] {Hématome intra-\\parenchymateux};
        \draw[<-, popred, thick] (-2.0, 1.1) -- (-3.0, 0.6) node[left, align=right, popred] {Effet de masse\\(compression, déviation ligne médiane)};
        \node[poporange, align=center] at (-1.5, -0.7) {\textbf{AVC Hémorragique}};
        \node[poporange, align=center, font=\tiny] at (-1.5, -1.1) {Rupture artérielle $\rightarrow$ Hématome $\rightarrow$ Compression};
    \end{scope}
\end{tikzpicture}
\legende{Schéma comparatif des deux types d'AVC. À gauche : AVC ischémique par obstruction de l'artère cérébrale moyenne (thrombus rouge sur le tronc) entraînant une zone d'infarctus ischémique (bleu clair) dans son territoire (cortex frontal/pariétal/temporal controlatéral). À droite : AVC hémorragique par rupture d'une artère perforante profonde (artère lenticulostriée) formant un hématome (rouge foncé) dans les noyaux de la base, comprimant le tissu cérébral sain et déviant la ligne médiane. Le scanner (TDM) sans injection est l'examen de première intention (disponible, rapide) pour les différencier : hypodensité (ischémie) vs hypersignal blanc (sang).}
\end{popfigure}

\noindent \textbf{Signes cliniques d'alerte : la règle \cle{VITE} (adaptée de FAST internationale) :}
\begin{center}
\begin{tabularx}{\linewidth}{|l|X|}\hline
\textbf{Lettre} & \textbf{Signe à rechercher (Face à la victime -- test simple)} \\ \hline
\textbf{V} - \textbf{Visage} & Demander de \textbf{sourire} / montrer les dents : \textbf{déviation de la commissure des lèvres} (bouche tordue vers le côté sain), aplatissement du sillon naso-génien du côté paralysé. \\ \hline
\textbf{I} - \textbf{Incapacité motrice} & Demander de \textbf{lever les deux bras} (paumes vers le haut, yeux fermés) pendant 10s : \textbf{dérive} ou \textbf{paralysie} (hémiplégie/hémiparésie) d'un côté (bras retombe, paume pronation). Idem pour les jambes (décollement talon 5s). \\ \hline
\textbf{T} - \textbf{Trouble de la parole} & Demander de \textbf{répéter une phrase simple} (ex: « Le ciel est bleu ») : \textbf{aphasie} (ne trouve pas les mots, phrase incohérente, ne comprend pas), \textbf{dysarthrie} (parole pâteuse, « langue épaisse », articulation difficile), mutisme. \\ \hline
\textbf{E} - \textbf{Extrême urgence / Appeler le 119} & \textbf{Noter l'heure exacte de début des signes} (ou heure « dernier vu normal » si sommeil/début non vu). \textbf{Ne pas attendre.} Allonger la victime, tête légèrement surélevée (30°). Ne rien donner à manger/boire. Rassurer. \\ \hline
\end{tabularx}
\end{center}

\begin{propriete}[Rôle majeur de l'hypertension artérielle (admise au Probatoire)]
L'\textbf{hypertension artérielle (HTA)} est le \textbf{principal facteur de risque modifiable} de l'AVC (ischémique et hémorragique) et de l'infarctus du myocarde. Elle fragilise la paroi artérielle (athérosclérose accélérée des gros troncs, hyalinose des artérioles, micro-anévrismes de Charcot-Bouchard sur les petites artères cérébrales perforantes) et favorise la formation de thrombus (stase, endothélium lésé) ou la rupture vasculaire. \textbf{Contrôler sa TA (cible $<$ 140/90 mmHg, voire $<$ 130/80 mmHg selon profil) divise par 2 à 4 le risque d'AVC.}
\end{propriete}

\courssub{L'hypertension artérielle (HTA) : le « tueur silencieux »}

\noindent Contrairement à l'infarctus ou à l'AVC qui sont des \textbf{événements aigus} (brutaux, symptômes dramatiques), l'HTA est une \textbf{maladie chronique}, le plus souvent \textbf{asymptomatique} pendant des années, d'où son surnom.

\begin{definition}[Hypertension Artérielle (HTA) -- Critères diagnostiques OMS/ESC simplifiés]
On parle d'\textbf{HTA} lorsque la tension artérielle mesurée à plusieurs reprises (au moins \textbf{3 mesures sur 3 consultations distinctes} espacées de quelques semaines, hors contexte d'urgence) au cabinet médical est \textbf{durablement supérieure ou égale à \SI{140}{mmHg} (systolique) et/ou \SI{90}{mmHg} (diastolique)} (soit \SI{14}{cmHg}/\SI{9}{cmHg}).
\begin{itemize}
    \item \textbf{HTA Grade 1 (modérée) :} \SIrange{140}{159}{mmHg} / \SIrange{90}{99}{mmHg}.
    \item \textbf{HTA Grade 2 (sévère) :} \SIrange{160}{179}{mmHg} / \SIrange{100}{109}{mmHg}.
    \item \textbf{HTA Grade 3 (très sévère) :} $\ge$ \SI{180}{mmHg} / $\ge$ \SI{110}{mmHg}.
\end{itemize}
\emph{Nota :} En automesure validée (domicile, protocole règle des 3 : 3 matins, 3 soirs, 3 jours), les seuils diagnostiques sont abaissés à $\ge$ \SI{135}{mmHg} / $\ge$ \SI{85}{mmHg}. L'Hypertension Artérielle Masquée (normale au cabinet, élevée à domicile) et l'Effet Blouse Blanche (élevée au cabinet, normale à domicile) justifient l'automesure ou la MAPA (Mesure Ambulatoire de la Pression Artérielle sur 24h).
\end{definition}

\noindent \textbf{Conséquences de l'HTA non dépistée / non traitée (lésions d'organes cibles -- Preuve de l'atteinte) :}
\begin{itemize}
    \item \textbf{Cœur :} \textbf{Hypertrophie ventriculaire gauche (HVG)} (cœur gros, parois épaissies, rigide, ischémie sous-endocardique) $\rightarrow$ \textbf{insuffisance cardiaque} (diastolique puis systolique), \textbf{arythmies} (fibrillation atrielle $\rightarrow$ risque d'embolie systémique/AVC), \textbf{angine de poitrine / infarctus du myocarde} (coronaires atteintes).
    \item \textbf{Cerveau :} \textbf{AVC ischémique} (athérosclérose carotidienne/intracrânienne, infarctus lacunaires) ou \textbf{AVC hémorragique} (rupture micro-anévrismes), \textbf{démence vasculaire} (atteinte substance blanche, troubles de la marche/mémoire).
    \item \textbf{Reins :} \textbf{Néphro-angiosclérose} (artériolosclérose hyaline, ischémie glomérulaire) $\rightarrow$ \textbf{insuffisance rénale chronique} (protéinurie, élévation créatinine, DFG diminuée), stade terminal $\rightarrow$ dialyse/greffe.
    \item \textbf{Yeux :} \textbf{Rétinopathie hypertensive} (stades Keith-Wagener-Barker : rétrécissement artérioles $\rightarrow$ croisements artério-veineux $\rightarrow$ hémorragies/exsudats $\rightarrow$ papillite/œdème papillaire) $\rightarrow$ baisse de vision, cécité.
    \item \textbf{Artères périphériques :} \textbf{Artériopathie Oblitérante des Membres Inférieurs (AOMI)} (sténoses iléo-fémorales) $\rightarrow$ claudication (douleur à la marche), ischémie critique, risque d'amputation.
\end{itemize}

\begin{attention}[Piège majeur : l'absence de symptômes $\neq$ absence de maladie]
\textbf{Erreur fréquente au Probatoire et dans la vie courante :} « Je me sens bien, je n'ai pas mal à la tête, donc ma tension est bonne. » \textbf{FAUX.} L'HTA est \textbf{silencieuse (asymptomatique)} dans plus de 90\% des cas aux stades précoces et modérés. Les céphalées matinales (réveil), bourdonnements d'oreilles (acouphènes), vertiges ou saignements de nez (épistaxis) sont \textbf{inconstants, non spécifiques et tardifs}, signes d'une HTA très sévère (Grade 3) ou mal contrôlée, voire d'une urgence hypertensive. \textbf{Seule la mesure répétée et rigoureuse de la TA (cabinet, domicile, MAPA) permet le diagnostic.} Attendre des symptômes pour consulter, c'est découvrir l'HTA au stade de complication irréversible (AVC, insuffisance cardiaque, insuffisance rénale terminale).
\end{attention}

\courssub{Conduite à tenir face à une victime : chaque minute compte}

\noindent Devant un malaise brutal évoquant un infarctus ou un AVC, la conduite à tenir est \textbf{standardisée, simple et vitale}. Elle ne demande pas de connaissances médicales complexes, mais du sang-froid, de la rapidité et le respect de l'ordre des actions.

\begin{methode}[Conduite à tenir devant une douleur thoracique suspecte d'infarctus ou signes d'AVC (VITE)]
\begin{enumerate}
    \item \textbf{ALERTER IMMÉDIATEMENT LES SECOURS :} Composer le \textbf{119 (SAMU Cameroun)} ou le numéro d'urgence local (ex: 112). \textbf{Ne pas appeler le médecin traitant, ne pas emmener soi-même la victime à l'hôpital en taxi/moto/voiture particulière.} Le SAMU apporte le plateau technique (ECG, défibrillateur DAE, thrombolyse, monitoring) \textbf{sur place} et oriente vers l'unité adaptée (USIC pour infarctus, Unité Neuro-Vasculaire UNV pour AVC). \emph{Le transport par les secours médicalisés gagne du temps médical.}
    \item \textbf{INSTALLER LA VICTIME AU REPOS ABSOLU (Position adaptée) :}
        \begin{itemize}
            \item \textbf{Infarctus suspecté (douleur thoracique typique) :} Position \textbf{demi-assise} (buste relevé 30--45°, jambes pliées si possible) pour réduire le retour veineux (précharge) et le travail cardiaque, soulager l'essoufflement. Desserer vêtements (cravate, ceinture, col, soutien-gorge).
            \item \textbf{AVC suspecté (signes VITE positifs) :}
                \begin{itemize}
                    \item Victime \textbf{consciente} : Allongée, \textbf{tête légèrement surélevée (30°)} pour réduire la pression intracrânienne (PIC) et favoriser le drainage veineux cérébral. Tête tournée du côté paralysé si salivation.
                    \item Victime \textbf{inconsciente / troubles de la vigilance majeurs} : \textbf{Position Latérale de Sécurité (PLS) côté paralysé} (côté de la lésion) pour libérer les voies aériennes (écoulement salives/vomissements), prévenir l'inhalation, tout en surveillant la respiration.
                \end{itemize}
        \end{itemize}
    \item \textbf{RASSURER, SURVEILLER, NE RIEN DONNER :} Parler calmement, expliquer qu'on a appelé les secours. \textbf{Ne rien donner à boire ni à manger} (risque de fausse route pulmonaire si trouble de la déglutition/inconscience ; nécessité d'estomac vide pour une éventuelle anesthésie générale / intubation / angioplastie). Surveiller en continu : conscience (appel verbal, stimulation douloureuse), respiration (fréquence, aspect), pouls (fréquence, régularité). Si \textbf{arrêt cardiaque} (inconscient, ne respire pas ou respiration agonique) : \textbf{DÉFIBRILLER (DAE si dispo) et MASSER (RCP : 30 compressions / 2 insufflations)} immédiatement sans attendre le SAMU.
    \item \textbf{TRANSMETTRE LES INFOS ESSENTIELLES AU 119 (Mémo I.S.T.) :} \textbf{I}dentité/Âge, \textbf{S}ymptômes (douleur, paralysie, heure début), \textbf{T}raitements/Antécédents (HTA, diabète, cardiaque, anticoagulants). Lieu précis (repères, code porte, étage).
\end{enumerate}
\end{methode}

\begin{popfigure}
\begin{tikzpicture}[scale=0.85, every node/.style={font=\small, align=center},
    decision/.style={diamond, draw=popblue, thick, fill=popblue!10, aspect=2, inner sep=4pt},
    action/.style={rectangle, draw=popgreen, thick, fill=popgreen!10, rounded corners, inner sep=6pt},
    sign/.style={rectangle, draw=poporange, thick, fill=poporange!10, rounded corners, inner sep=6pt},
    arrow/.style={-Stealth, thick, popdark},
    start/.style={ellipse, draw=poppurple, thick, fill=poppurple!10, inner sep=6pt},
    urgent/.style={rectangle, draw=popred, thick, fill=popred!10, rounded corners, inner sep=6pt}]
    
    % Nœuds
    \node[start, align=center] (start){Malaise brutal\\Témoin / Victime};
    \node[sign, below=1.2cm of start] (signs) {Signes observés ?};
    
    % Branche Infarctus
    \node[decision, below left=1.5cm and 2.8cm of signs, align=center] (dec_inf){Douleur thoracique\\intense $\ge$ 20 min,\\irradiation bras/gauche/mâchoire,\\sueurs, angoisse ?};
    \node[action, below left=1.5cm and 1.5cm of dec_inf, align=center] (act_inf){\textbf{INFARCTUS SUSPECTÉ}\\\textit{Position demi-assise (30-45°) + Aspirine 300 mg à croquer*}};
    
    % Branche AVC
    \node[decision, below right=1.5cm and 2.8cm of signs, align=center] (dec_avc){Signes \textbf{VITE} positifs :\\(Visage tordu, Incapacité\\bras/jambe, Trouble parole) ?};
    \node[action, below right=1.5cm and 1.5cm of dec_avc, align=center] (act_avc){\textbf{AVC SUSPECTÉ}\\\textit{Conscient : Allongé tête 30° ; Inconscient : PLS côté paralysé}};
    
    % Branche Autre / Inconnu
    \node[decision, below=2.5cm of signs, align=center] (dec_other){Autre malaise\\ (inconscient sans signes VITE,\\dyspnée isolée, douleur abdominale...) ?};
    \node[action, below=1.5cm of dec_other, align=center] (act_other){Appeler 119\\Adapter position\\Surveiller constants};
    
    % Convergence : Appel 119
    \node[urgent, below=2.5cm of act_inf, xshift=3.5cm, align=center] (call){\textbf{APPELER LE 119 (SAMU) IMMÉDIATEMENT}\\\textit{Heure début - Lieu précis - Symptômes - Antécédents - Traitements}};
    
    % Post-appel
    \node[action, below=1.5cm of call, align=center] (post){Ne rien donner à boire/manger\\Rassurer - Surveiller respiration/conscience\\Préparer ordonnances / carte groupe sanguin / DAE};
    
    % Flèches
    \draw[arrow] (start) -- (signs);
    \draw[arrow] (signs) -- node[left, pos=0.5, popmuted, font=\tiny] {OUI} (dec_inf);
    \draw[arrow] (signs) -- node[right, pos=0.5, popmuted, font=\tiny] {OUI} (dec_avc);
    \draw[arrow] (signs) -- node[above, pos=0.5, popmuted, font=\tiny] {NON / AUTRE} (dec_other);
    
    \draw[arrow] (dec_inf) -- node[left, pos=0.5, popmuted, font=\tiny] {OUI} (act_inf);
    \draw[arrow] (dec_avc) -- node[right, pos=0.5, popmuted, font=\tiny] {OUI} (act_avc);
    \draw[arrow] (dec_other) -- (act_other);
    
    \draw[arrow] (act_inf) -- (call);
    \draw[arrow] (act_avc) -- (call);
    \draw[arrow] (act_other) -- (call);
    \draw[arrow] (call) -- (post);
    
    % Non pour les décisions (vers suivant)
    \draw[arrow, popmuted, dashed] (dec_inf) -- ++(-1.5,0) node[left, popmuted, font=\tiny] {NON} |- (dec_avc);
    \draw[arrow, popmuted, dashed] (dec_avc) -- ++(1.5,0) node[right, popmuted, font=\tiny] {NON} |- (dec_other);
    
    % Note aspirine
    \node[font=\tiny, popmuted, anchor=north] at (act_inf.south) {*Si pas d'allergie, pas d'hémorragie active, pas d'ulcère gastrique évolutif};
\end{tikzpicture}
\legende{Arbre de décision pour la conduite à tenir face à un malaise brutal. L'identification rapide des signes cardinaux (douleur thoracique typique d'infarctus ou signes VITE d'AVC) oriente vers la position adaptée (demi-assis vs PLS/tête 30°), mais \textbf{l'appel au 119 est l'action commune immédiate et prioritaire dans tous les cas}. L'aspirine 300 mg (mâchée) est recommandée dès le soupçon d'infarctus si pas de contre-indication.}
\end{popfigure}

\begin{savaistu}[Le temps, c'est du cerveau : la fenêtre thérapeutique de l'AVC]
Au Cameroun comme ailleurs, la prise en charge de l'AVC ischémique a été révolutionnée par la \textbf{thrombolyse intraveineuse par rt-PA (activase, altéplase)} (injection d'un activateur du plasminogène qui dissout le caillot : fibrinolyse). Mais elle ne peut être réalisée que dans une \textbf{fenêtre de 4h30 maximum après le début des signes} (parfois étendue à 9h en imagerie de perfusion sélectionnant le « penumbra » viable). Au-delà, le risque d'hémorragie de transformation (hémorragie parenchymateuse symptomatique) dépasse le bénéfice clinique. Pour l'AVC hémorragique, la neurochirurgie (évacuation de l'hématome si volumineux, accessible, détérioration clinique) ou le traitement médical intensif (contrôle strict de la TA : cible 130-140 mmHg systolique en 1h, contrôle de la PIC, prévention complications) doivent débuter \textbf{le plus précocement possible}. \textbf{Chaque minute de retard = 1,9 millions de neurones perdus, 14 milliards de synapses, 12 km de fibres myéliniques.} D'où l'importance vitale de noter l'heure de début (ou l'heure « dernier vu normal » si réveil avec déficit) et d'appeler le 119 sans délai. Les Unités Neuro-Vasculaires (UNV) se développent dans les grands centres (CHU Yaoundé, Douala, Hôpital Central, Hôpital Général) pour cette prise en charge spécialisée multidisciplinaire (neurologue, réanimateur, radiologue interventionnel, infirmiers formés, kiné, orthophoniste).
\end{savaistu}

\begin{exoresolu}[Analyse d'une situation d'urgence cardiovasculaire]
\textbf{Énoncé :} Lors d'un match de football inter-quartiers à Douala-Bonabéri, un spectateur de 52 ans, M. N., s'effondre soudainement. Il est inconscient, ne répond pas aux stimulations verbales ni douloureuses (score de Glasgow bas). Sa respiration est bruyante, irrégulière, laborieuse (respiration agonique / « gasping »). Un témoin signale qu'il s'est plaint d'un « mal de tête violent » il y a 10 minutes avant de s'effondrer. Il est connu hypertendu, mais « oublie souvent ses médicaments ».
\begin{enumerate}
    \item Quel type d'accident cardiovasculaire suspectez-vous en priorité ? Justifiez par deux arguments cliniques.
    \item Quelle est la conduite à tenir immédiate du premier témoin ? Détaillez les 3 premières actions dans l'ordre chronologique.
    \item Pourquoi ne doit-on \textbf{pas} lui donner de l'eau, du sucre, ou son traitement antihypertenseur habituel ?
\end{enumerate}
\tcblower
\textbf{Solution détaillée :}
\begin{enumerate}
    \item \textbf{Suspicion diagnostique prioritaire : AVC hémorragique (hémorragie intra-cérébrale hypertensive, siège probable putamen/thalame ou hémorragie méningée anévrismale).}
        \textbf{Arguments cliniques majeurs :}
        \begin{itemize}
            \item \textbf{Céphalée brutale violente (« coup de massue », céphalée tonnerre)} survenant 10 min avant l'effondrement : signe d'alerte cardinal de rupture vasculaire intracrânienne (augmentation brutale de la pression intravasculaire sur paroi fragilisée par l'HTA chronique).
            \item \textbf{Inconscience brutale + respiration agonique (bruyante, irrégulière, « gasping »)} : traduction d'une atteinte tronco-cérébrale (compression du tronc cérébral / bulbe rachidien par l'hématome expansif + œdème péritumoral + effet de masse avec herniation tentorielle/amygdalienne). La respiration agonique est un signe de détresse vitale extrême, souvent confondu à tort avec une respiration normale.
        \end{itemize}
        \emph{Note différentielle :} Un infarctus du myocarde compliqué d'arrêt cardiaque peut donner un tableau d'inconscience + respiration agonique, mais le contexte de céphalée initiale violente chez un hypertendu non observant oriente massivement vers l'AVC hémorragique. L'ECG (si fait) montrerait un rythme (souvent rythme sinusal ou bradycardie par compression bulbaire) et non une fibrillation ventriculaire primaire.
    \item \textbf{Conduite à tenir immédiate (3 premières actions dans l'ordre) :}
        \begin{enumerate}
            \item \textbf{Appeler le 119 (SAMU)} en précisant clairement : « Homme 52 ans, effondrement brutal sur terrain de sport, \textbf{inconscient}, \textbf{respiration anormale (agonique)}, antécédent \textbf{HTA non observante}, \textbf{céphalée violente précoce} (il y a 10 min). Lieu précis : terrain Bonabéri, tribune X. » Demander si un \textbf{DAE (Défibrillateur Automatisé Externe)} est disponible sur site (stade, mairie, pharmacie voisine).
            \item \textbf{Installer en Position Latérale de Sécurité (PLS) côté droit (préférentiellement) :} Victime inconsciente \textbf{respirant} (même agonique) $\rightarrow$ risque majeur d'inhalation (fausse route) par perte des réflexes de protection des voies aériennes (toux, déglutition) et relaxation du sphincter œsophagien. La PLS maintient les voies aériennes pérennes et permet l'écoulement des fluides (salive, sang, vomissements). \textbf{Ne pas laisser sur le dos (décubitus dorsal).} Si traumatisme crânien suspecté (chute), PLS modifiée (stabilisation tête-cou).
            \item \textbf{Surveiller la respiration en continu (monitoring visuel/manuel) :} Si la respiration s'arrête complètement (arrêt cardiaque confirmé : absence de respiration normale, absence de pouls carotidien en $<$ 10s) $\rightarrow$ \textbf{Démarrer immédiatement la Réanimation Cardio-Pulmonaire (RCP) : 30 compressions thoraciques (centre thorax, 5-6 cm, 100-120/min, relâchement complet) / 2 insufflations} (bouche-à-bouche ou masque poche si formé) \textbf{+ utilisation du DAE dès disponibilité} (coller électrodes, suivre instructions vocales). Ne pas interrompre les compressions pour autre chose que le choc ou le contrôle rythme DAE.
        \end{enumerate}
    \item \textbf{Contre-indications formelles (erreurs fréquentes à ne pas commettre) :}
        \begin{itemize}
            \item \textbf{Eau / Sucre / Aliment / Médicament par voie orale :} Risque majeur de \textbf{fausse route pulmonaire} (inhalation) car les réflexes de déglutition et de protection des voies aériennes (toux) sont abolis par l'inconscience (atteinte bulbaire). L'estomac doit rester vide pour une éventuelle intubation à séquence rapide / anesthésie d'urgence en réanimation/UNV.
            \item \textbf{Traitement antihypertenseur habituel (per os) :} \textbf{Interdit formellement en phase ultra-aiguë.} Une baisse brutale, non contrôlée et imprévisible de la TA (par médicament oral mal absorbé, effet de premier passage hépatique variable, pic d'action retardé) aggrave l'ischémie cérébrale périlésionnelle (pénombre ischémique) qui est \textbf{perfusion-dépendante} (autoregulation perdue). Cela peut provoquer un arrêt cardiaque par hypoperfusion coronaire/cérébrale. La TA doit être gérée en milieu hospitalier spécialisé (UNV/Réanimation) par voie veineuse (perfusion électrique) avec des molécules titrables à demi-vie courte (nicardipine, labétalol, urapidil, clevidipine) sous monitoring continu invasif (artériel) et neurologique.
        \end{itemize}
\end{enumerate}
\end{exoresolu}

\begin{aretenir}[Synthèse des 3 accidents majeurs -- Fiche de révision Probatoire]
\begin{center}
\begin{tabularx}{\linewidth}{|l|X|X|X|}\hline
\textbf{Accident} & \textbf{Mécanisme principal} & \textbf{Signes cardinaux (Sémiologie)} & \textbf{Urgence / Action 1 immédiate} \\ \hline
\textbf{Infarctus du myocarde} & Obstruction coronaire aiguë (thrombus sur plaque d'athérome fissurée) $\rightarrow$ Ischémie $\rightarrow$ Nécrose myocarde (troponines $\uparrow$) & Douleur thoracique intense $\ge$ 20 min, type étau/barre, au repos, irradiation bras G / mâchoire / dos, sueurs froides, angoisse, dyspnée. ECG : Sus-décalage ST (STEMI) ouonde Q / onde T négative (NSTEMI). & \textbf{1. Appeler 119} \\ \textbf{2. Position demi-assise} \\ \textbf{3. Aspirine 300 mg mâchée} (si pas CI) \\ \textbf{4. O2 si SpO2 $<$ 94\%} \\ \hline
\textbf{AVC Ischémique} & Obstruction artère cérébrale (thrombose athéromateuse / embolie cardiaque FA) $\rightarrow$ Infarctus cérébral (nécrose neuronale) & Règle \textbf{VITE} : \textbf{V}isage tordu (sourire asymétrique), \textbf{I}ncapacité motrice (bras/jambe dérivent/tombent), \textbf{T}rouble parole (aphasie/dysarthrie), \textbf{E}xtrême urgence. Déficit controlatéral à la lésion. & \textbf{1. Appeler 119} \\ \textbf{2. Allongé tête 30°} (conscient) ou \textbf{PLS côté paralysé} (inconscient) \\ \textbf{3. Noter heure début} (fenêtre thrombolyse 4h30) \\ \hline
\textbf{AVC Hémorragique} & Rupture artère cérébrale (HTA $\rightarrow$ micro-anévrismes Charcot-Bouchard / anévrisme sacciforme) $\rightarrow$ Hématome + Compression (PIC $\uparrow$) & Céphalée brutale violente (tonnerre), vomissements en jet (PIC), inconscience rapide, signes VITE (souvent majorés), signes de focalisation (hémiparésie, déviation regard conjugué vers l'hématome). Scanner : hypersignal blanc. & \textbf{1. Appeler 119} \\ \textbf{2. PLS} (souvent inconscient) \\ \textbf{3. PAS d'aspirine / anticoagulant} \\ \textbf{4. Contrôle TA hôpital uniquement} \\ \hline
\textbf{HTA (Maladie chronique)} & Pression artérielle durable $\ge$ 140/90 mmHg $\rightarrow$ Remodelage vasculaire (athérosclérose, hyalinose) + Remodelage cardiaque (HVG) $\rightarrow$ Lésions organes cibles & \textbf{Souvent AUCUN symptôme} (tueur silencieux). Céphalées matinales, acouphènes, vertiges, épistaxis = signes tardifs / HTA sévère. Dépistage : automesure (règle des 3) ou cabinet (3 mesures / 3 consultations). & \textbf{1. Dépistage régulier} (annuel après 40 ans, ou si facteur de risque) \\ \textbf{2. Hygiène de vie} (sel, poids, sport, alcool, tabac) \\ \textbf{3. Observance traitement} (vie entière, ne pas arrêter seul) \\ \hline
\end{tabularx}
\end{center}
\end{aretenir}

\noindent \textbf{En résumé :} L'infarctus du myocarde, l'AVC (ischémique et hémorragique) et l'hypertension artérielle forment un continuum pathophysiologique où l'athérosclérose et l'hypertension sont les artisans silencieux de la rupture brutale. Connaître les signes d'alerte (douleur thoracique typique, règle VITE) et le réflexe unique \textbf{« J'appelle le 119 »} sauve des vies. La section suivante abordera les \textbf{facteurs de risque cardiovasculaires} sur lesquels nous pouvons agir pour prévenir ces catastrophes.

\coursec{Les facteurs de risque cardiovasculaire}

\courssub{Transition depuis les accidents cardiovasculaires}
Dans la section précédente, nous avons identifié les trois accidents cardiovasculaires majeurs — infarctus du myocarde, accident vasculaire cérébral (AVC) et hypertension artérielle (HTA) — et décrit leurs manifestations cliniques. Nous nous posons maintenant la question fondamentale : \emph{qu'est‑ce qui favorise l'apparition de ces accidents ?} La réponse réside dans la notion de \cle{facteur de risque}. Comprendre ces facteurs, les classer et savoir les repérer dans une situation concrète est la clé de la prévention primaire.

\begin{definition}[Facteur de risque cardiovasculaire]
Un \textbf{facteur de risque cardiovasculaire} est une caractéristique individuelle, un comportement ou une condition pathologique qui \textbf{augmente la probabilité} de survenue d'un accident cardiovasculaire (infarctus, AVC, HTA sévère, etc.). Il ne s'agit pas d'une cause unique et suffisante, mais d'un élément qui \emph{pèse} sur la balance du risque.
\end{definition}

\courssub{Classification des facteurs de risque}
Les facteurs de risque se répartissent en deux grandes familles : ceux sur lesquels nous ne pouvons pas agir (\emph{non modifiables}) et ceux sur lesquels une intervention est possible (\emph{modifiables}). Le tableau ci‑dessous synthétise cette classification.

\begin{popfigure}
\begin{tikzpicture}[
  cell/.style={rectangle,draw=popline,minimum height=1cm,align=center,text width=7cm},
  header/.style={cell,fill=popblueL,font=\bfseries,text width=7cm},
  leftcol/.style={cell,fill=popgreenL,text width=5cm},
  rightcol/.style={cell,fill=poporangeL,text width=5cm}
]
% En-têtes
\node[header] (h1) at (0,0) {Facteurs \textbf{non modifiables}};
\node[header] (h2) at (14,0) {Facteurs \textbf{modifiables}};
% Lignes
\node[leftcol] (L1) at (0,-1) {Âge (risque croissant après 50 ans chez l'homme, 60 ans chez la femme)};
\node[rightcol] (R1) at (14,-1) {Alimentation déséquilibrée (excès de sel, graisses saturées, sucres ; déficit en fruits/légumes)};
\node[leftcol] (L2) at (0,-2) {Sexe (homme \textgreater femme avant la ménopause)};
\node[rightcol] (R2) at (14,-2) {Tabagisme (cigarettes, tabac à chiquer, narguilé)};
\node[leftcol] (L3) at (0,-3) {Hérédité (antécédents familiaux d'infarctus ou d'AVC avant 55 ans chez le père/frère, avant 65 ans chez la mère/sœur)};
\node[rightcol] (R3) at (14,-3) {Consommation excessive d'alcool};
\node[leftcol] (L4) at (0,-4) {};
\node[rightcol] (R4) at (14,-4) {Sédentarité (activité physique \textless 150 min par semaine)};
\node[leftcol] (L5) at (0,-5) {};
\node[rightcol] (R5) at (14,-5) {Stress chronique mal géré};
\node[leftcol] (L6) at (0,-6) {};
\node[rightcol] (R6) at (14,-6) {Obésité (IMC \textgreater 30 kg/m²) et surpoids abdominal};
\node[leftcol] (L7) at (0,-7) {};
\node[rightcol] (R7) at (14,-7) {Diabète de type 2 mal équilibré};
\node[leftcol] (L8) at (0,-8) {};
\node[rightcol] (R8) at (14,-8) {Dyslipidémie (LDL‑cholestérol élevé, HDL bas)};
\node[leftcol] (L9) at (0,-9) {};
\node[rightcol] (R9) at (14,-9) {Hypertension artérielle non traitée ou mal contrôlée};
% Lignes de séparation
\foreach \i in {1,...,9}{
  \draw[popline] (L\i.north west) -- (R\i.north east);
}
\draw[popline,thick] (h1.south west) -- (h2.south east);
\end{tikzpicture}
\legende{Tableau récapitulatif des facteurs de risque cardiovasculaires classés en non modifiables (à gauche) et modifiables (à droite).}
\end{popfigure}

\begin{attention}[Piège fréquent]
Ne pas croire que \emph{seules les personnes âgées} sont concernées. Les habitudes néfastes (tabac, alimentation riche en sel et graisses, inactivité) adoptées dès l'adolescence initient le processus athéromateux qui se manifestera des décennies plus tard. La prévention commence \emph{dès le plus jeune âge}.
\end{attention}

\courssub{Facteurs non modifiables : âge, sexe, hérédité}
\begin{itemize}
  \item \textbf{Âge} : le vieillissement artériel (perte d'élasticité, épaississement de l'intima) augmente mécaniquement le risque. Après 50 ans chez l'homme et 60 ans chez la femme, l'incidence des accidents cardiovasculaires double environ chaque décennie.
  \item \textbf{Sexe} : avant la ménopause, les femmes bénéficient d'un effet protecteur des œstrogènes (vasodilatation, profil lipidique favorable). Après la ménopause, le risque rejoint celui des hommes.
  \item \textbf{Hérédité} : un antécédent d'infarctus ou d'AVC survenu avant 55 ans chez un parent du premier degré de sexe masculin (père, frère) ou avant 65 ans chez un parent de sexe féminin (mère, sœur) multiplie le risque par 2 à 3. Ce facteur traduit une prédisposition génétique (polymorphismes du gène de l'apolipoprotéine E, du récepteur LDL, etc.) mais aussi un partage d'habitudes de vie familiales.
\end{itemize}

\courssub{Facteurs modifiables liés à l'alimentation}
L'alimentation est le levier le plus accessible. Les mécanismes physiopathologiques sont bien établis :
\begin{itemize}
  \item \textbf{Excès de sel (NaCl)} : augmente la volémie et la résistance vasculaire périphérique $\rightarrow$ élévation de la pression artérielle. L'OMS recommande \textless 5 g/jour (≈ 2 g de sodium).
  \item \textbf{Excès de graisses saturées} (viandes grasses, beurre, huile de palme, fromages) : élèvent le LDL‑cholestérol, substrat principal de la plaque d'athérome.
  \item \textbf{Excès de sucres rapides} (boissons sucrées, pâtisseries) : favorisent la prise de poids, l'insulinorésistance et la dyslipidémie (hypertriglycéridémie, HDL bas).
  \item \textbf{Pauvre en fruits et légumes} : déficit en potassium (effet antihypertenseur), en fibres (régulation du cholestérol) et en antioxydants (protection de l'endothélium).
\end{itemize}

\begin{exemplebox}[Repas typique à risque au Cameroun]
Un déjeuner composé de \emph{ndolé} très salé, de \emph{plantains frits} dans l'huile de palme, accompagné d'une boisson gazeuse sucrée, apporte en un seul repas : \textgreater 3 g de sel, \textgreater 30 g de graisses saturées, \textgreater 40 g de sucres ajoutés et presque aucun légume cru. Répété quotidiennement, ce profil alimentaire construit le terrain de l'athérosclérose.
\end{exemplebox}

\courssub{Facteurs modifiables liés au mode de vie}
\begin{itemize}
  \item \textbf{Tabagisme} : le monoxyde de carbone (CO) fixe l'hémoglobine $\rightarrow$ hypoxie tissulaire ; la nicotine stimule le système sympathique $\rightarrow$ vasoconstriction et tachycardie ; les radicaux libres oxydent le LDL‑cholestérol $\rightarrow$ plaque d'athérome plus instable. Même une consommation faible (1–4 cigarettes/jour) double le risque d'infarctus.
  \item \textbf{Alcool} : une consommation modérée (≤ 1 verre/jour chez la femme, ≤ 2 chez l'homme) peut avoir un effet neutre ou légèrement protecteur sur le HDL, mais au‑delà l'effet est délétère : HTA, cardiomyopathie dilatée, arythmies, prise de poids.
  \item \textbf{Sédentarité} : l'inactivité physique réduit la sensibilité à l'insuline, abaisse le HDL, favorise l'obésité abdominale et rigidifie les artères. L'OMS préconise ≥ 150 min par semaine d'activité d'intensité modérée (marche rapide, vélo, danse).
  \item \textbf{Stress chronique} : activation permanente de l'axe hypothalamo‑hypophyso‑surrénalien $\rightarrow$ cortisol et catécholamines élevés $\rightarrow$ HTA, inflammation bas bruit, comportements compensatoires (grignotage, tabac, alcool).
\end{itemize}

\courssub{Facteurs modifiables liés à l'état de santé}
Ces facteurs sont souvent la conséquence des précédents, mais ils deviennent à leur tour des leviers d'action thérapeutique :
\begin{itemize}
  \item \textbf{Obésité (IMC ≥ 30)} et \textbf{surpoids abdominal} (tour de taille > 94 cm homme, > 80 cm femme) : tissu adipeux viscéral sécrète des adipokines pro‑inflammatoires (TNF‑α, IL‑6) et favorise l'insulinorésistance.
  \item \textbf{Diabète de type 2} : l'hyperglycémie chronique glyque les protéines vasculaires (produits de glycation avancée) $\rightarrow$ rigidité artérielle, dysfonction endothéliale, micro‑ et macro‑angiopathie.
  \item \textbf{Dyslipidémie} : LDL‑c élevé (cible \textless 1,6 g/L selon le niveau de risque) et HDL‑c bas (\textless 0,40 g/L homme, \textless 0,50 g/L femme).
  \item \textbf{Hypertension artérielle non traitée} : pression systolique ≥ 140 mmHg et/ou diastolique ≥ 90 mmHg (confirmée sur 3 mesures). L'HTA est à la fois cause et conséquence de la rigidité artérielle.
\end{itemize}

\courssub{L'athérosclérose : mécanisme central}
L'athérosclérose est la lésion anatomique sous‑jacente à la majorité des accidents cardiovasculaires. Elle résulte d'un dépôt progressif de lipides, de cellules inflammatoires et de tissu fibreux dans l'intima des artères de moyen et gros calibre. Le schéma ci‑dessous illustre les trois stades clés.

\begin{popfigure}
\begin{tikzpicture}[scale=0.9,>=Stealth]
% Artère saine
\begin{scope}[xshift=0cm]
  \draw[popblue,thick] (0,0) ellipse (2 and 1);
  \draw[popblue!50,thick] (0,0) ellipse (1.6 and 0.6);
  \draw[popblue!30,thick] (0,0) ellipse (1.2 and 0.3);
  \node[popblue,font=\bfseries] at (0,-1.5) {Artère saine};
  \draw[->,popline] (1.2,0.3) -- ++(0.5,0.5) node[right,font=\small] {Endothélium intact};
  \draw[->,popline] (1.6,0.6) -- ++(0.5,0.5) node[right,font=\small] {Média musculaire};
  \draw[->,popline] (2,1) -- ++(0.5,0.5) node[right,font=\small] {Adventice};
\end{scope}
% Stade précoce : strie graisseuse
\begin{scope}[xshift=6cm]
  \draw[poporange,thick] (0,0) ellipse (2 and 1);
  \draw[poporange!50,thick] (0,0) ellipse (1.6 and 0.6);
  \fill[poporange!30] (-1.2,0) .. controls (-1.2,0.3) and (-0.5,0.4) .. (0,0.3) .. controls (0.5,0.4) and (1.2,0.3) .. (1.2,0) -- cycle;
  \node[poporange,font=\bfseries] at (0,-1.5) {Strie graisseuse};
  \draw[->,popline] (0,0.3) -- ++(0,0.8) node[above,font=\small] {LDL oxydé + macrophages (cellules spumeuses)};
\end{scope}
% Stade avancé : plaque d'athérome
\begin{scope}[xshift=12cm]
  \draw[popred,thick] (0,0) ellipse (2 and 1);
  \draw[popred!50,thick] (0,0) ellipse (1.6 and 0.6);
  \fill[popred!30] (-1.5,0) .. controls (-1.5,0.5) and (-0.5,0.7) .. (0,0.7) .. controls (0.5,0.7) and (1.5,0.5) .. (1.5,0) -- cycle;
  \draw[popred!80,thick] (-0.2,0.7) -- (0.2,0.7);
  \node[popred,font=\bfseries] at (0,-1.5) {Plaque d'athérome};
  \draw[->,popline] (0,0.7) -- ++(0,0.8) node[above,font=\small] {Nécrose lipidique, cap fibreux, calcification};
  \draw[->,popline] (-1.6,0) -- ++(-0.6,0) node[left,font=\small] {Lumière rétrécie};
\end{scope}
% Stade complication : thrombose
\begin{scope}[xshift=18cm]
  \draw[poppurple,thick] (0,0) ellipse (2 and 1);
  \draw[poppurple!50,thick] (0,0) ellipse (1.6 and 0.6);
  \fill[poppurple!30] (-1.5,0) .. controls (-1.5,0.5) and (-0.5,0.7) .. (0,0.7) .. controls (0.5,0.7) and (1.5,0.5) .. (1.5,0) -- cycle;
  \fill[poppurple!80] (0,0.7) circle (0.3);
  \node[poppurple,font=\bfseries] at (0,-1.5) {Thrombose occlusive};
  \draw[->,popline] (0,1) -- ++(0,0.8) node[above,font=\small] {Rupture de la plaque + caillot};
  \draw[->,popline] (-1.6,0) -- ++(-0.6,0) node[left,font=\small] {Lumière obstruée};
\end{scope}
% Flèches d'évolution
\foreach \x in {4,10,16}{
  \draw[popline,->,thick] (\x,-1.5) -- ++(2,0);
}
\end{tikzpicture}
\legende{Évolution de la paroi artérielle : artère saine → strie graisseuse (dépôt de LDL oxydé et macrophages) → plaque d'athérome fibreuse (rétrécissement) → rupture de la plaque et thrombose occlusive (infarctus ou AVC).}
\end{popfigure}

\begin{propriete}[Potentialisation des facteurs de risque (admise)]
Les facteurs de risque ne s'additionnent pas simplement, ils \textbf{se potentialisent} (effet synergique). Un sujet hypertendu, fumeur et sédentaire présente un risque \textbf{bien supérieur} à la somme arithmétique des risques isolés. Cette interaction justifie la prise en charge globale plutôt que le traitement d'un seul facteur.
\end{propriete}

\courssub{Effet cumulatif : plus de facteurs, plus de danger}
Pour visualiser cette potentialisation, le diagramme en barres ci‑dessous montre l'évolution du niveau de risque global (faible, modéré, élevé) en fonction du nombre de facteurs de risque majeurs présents (selon l'échelle de risque SCORE adaptée).

\begin{popfigure}
\begin{tikzpicture}
\begin{axis}[
  ybar,
  bar width=1.2cm,
  width=\linewidth,
  height=8cm,
  ymin=0,ymax=100,
  ylabel={Risque cardiovasculaire à 10 ans (\% )},
  xtick={1,2,3,4,5},
  xticklabels={0 facteur,1 facteur,2 facteurs,3 facteurs,≥4 facteurs},
  xlabel={Nombre de facteurs de risque majeurs},
  ymajorgrids=true,
  grid style=dashed,
  nodes near coords,
  nodes near coords align={vertical},
  every node near coord/.append style={font=\small,popdark},
  legend style={at={(0.5,-0.15)},anchor=north,legend columns=-1},
  enlarge x limits=0.15
]
\addplot[fill=popgreenL,draw=popgreen] coordinates {(1,2) (2,5) (3,12) (4,25) (5,45)};
\addlegendentry{Risque faible (<10\%)}
\addplot[fill=poporangeL,draw=poporange] coordinates {(1,0) (2,0) (3,8) (4,20) (5,30)};
\addlegendentry{Risque modéré (10–20\%)}
\addplot[fill=popred!30,draw=popred] coordinates {(1,0) (2,0) (3,0) (4,5) (5,25)};
\addlegendentry{Risque élevé (>20\%)}
\end{axis}
\end{tikzpicture}
\legende{Risque cardiovasculaire à 10 ans selon le nombre de facteurs de risque majeurs (HTA, tabagisme, diabète, dyslipidémie, obésité abdominale). La courbe illustre la potentialisation : à partir de 3 facteurs, le risque bascule en zone élevée.}
\end{popfigure}

\courssub{Méthode : identifier les facteurs de risque dans une situation}
\begin{methode}[Analyse d'une situation de vie]
\begin{enumerate}
  \item \textbf{Recueillir les données} : âge, sexe, antécédents familiaux (père, mère, frères/sœurs), habitudes alimentaires (sel, graisses, sucres, fruits/légumes), consommation de tabac (nombre de cigarettes/jour), d'alcool (verres/semaine), niveau d'activité physique (minutes par semaine), stress perçu.
  \item \textbf{Noter les paramètres médicaux} : tension artérielle (moyenne de 3 mesures), IMC et tour de taille, glycémie à jeun ou HbA1c, profil lipidique (LDL‑c, HDL‑c, triglycérides).
  \item \textbf{Classer chaque élément} : \emph{non modifiable} (âge, sexe, hérédité) ou \emph{modifiable} (tous les autres).
  \item \textbf{Compter les facteurs majeurs} : HTA, tabagisme, diabète, dyslipidémie, obésité abdominale, sédentarité.
  \item \textbf{Évaluer le niveau de risque global} (faible / modéré / élevé) en s'appuyant sur une grille de risque validée (ex. SCORE, Framingham) ou, à défaut, sur le principe de potentialisation.
  \item \textbf{Proposer un plan d'action} : prioriser les facteurs modifiables les plus impactants (arrêt tabac, contrôle HTA, activité physique, rééquilibrage alimentaire).
\end{enumerate}
\end{methode}

\begin{exoresolu}[Analyse du cas de M. Ngo Bell]
\textbf{Énoncé.} M. Ngo Bell, 48 ans, agent de sécurité, fume 10 cigarettes/jour depuis 20 ans. Sa tension artérielle mesurée à trois reprises est de 150/95 mmHg. Il ne pratique aucune activité physique régulière. Son père est décédé d'un infarctus à 52 ans. Son IMC est de 29 kg/m² (tour de taille 98 cm). Glycémie à jeun 1,10 g/L, LDL‑c 1,70 g/L, HDL‑c 0,38 g/L.

\textbf{Solution détaillée.}
\tcblower
\begin{enumerate}
  \item \textbf{Facteurs non modifiables} : âge 48 ans (homme, risque croissant), sexe masculin, hérédité (père infarctus < 55 ans) → \textbf{3 facteurs non modifiables dont 1 majeur (hérédité)}.
  \item \textbf{Facteurs modifiables majeurs} :
  \begin{itemize}
    \item Tabagisme actif (10 cig/jour) → \textbf{1}.
    \item Hypertension artérielle non traitée (150/95) → \textbf{2}.
    \item Sédentarité (0 min par semaine) → \textbf{3}.
    \item Surpoids abdominal (tour de taille 98 cm > 94 cm) → \textbf{4}.
    \item Dyslipidémie (LDL‑c 1,70 > 1,60 ; HDL‑c 0,38 < 0,40) → \textbf{5}.
  \end{itemize}
  \item \textbf{Facteur additionnel} : Prédiabète (glycémie 1,10 g/L) → risque accru.
  \item \textbf{Bilan} : \textbf{6 facteurs de risque majeurs} (1 non modifiable + 5 modifiables) + prédiabète. Selon la grille de risque, M. Ngo Bell se situe en \textbf{risque élevé (>20\% à 10 ans)}.
  \item \textbf{Priorités d'intervention} : arrêt du tabac (gain immédiat), mise sous antihypertenseur, marche rapide 30 min/jour, rééquilibrage alimentaire (réduction sel, graisses saturées, augmentation fruits/légumes), suivi glycémique et lipidique à 3 mois.
\end{enumerate}
\end{exoresolu}

\begin{savaistu}[Le tabac, même à faible dose, est un poison vasculaire]
Dès la première cigarette, le monoxyde de carbone diminue le transport d'O₂ et la nicotine provoque une vasoconstriction coronarienne. Mais l'effet le plus insidieux est l'oxydation du LDL‑cholestérol, qui le rend athérogène. \textbf{Bonne nouvelle :} l'arrêt du tabac réduit le risque d'infarctus de 50 \% dès la première année, et après 10 à 15 ans le risque rejoint celui d'un non‑fumeur. Au Cameroun, le programme national de lutte contre le tabagisme (PNLT) propose des consultations d'aide à l'arrêt dans la plupart des hôpitaux de référence (Yaoundé, Douala, Garoua, etc.).
\end{savaistu}

\courssub{Synthèse des points clés}
\begin{aretenir}[À retenir pour le Probatoire]
\begin{itemize}
  \item Un facteur de risque cardiovasculaire augmente la probabilité d'accident (infarctus, AVC, HTA sévère).
  \item \textbf{Non modifiables} : âge, sexe, hérédité (infarctus/AVC avant 55 ans père/frère, avant 65 ans mère/sœur).
  \item \textbf{Modifiables} : alimentation (sel, graisses saturées, sucres, déficit fruits/légumes), tabac, alcool, sédentarité, stress, obésité, diabète, dyslipidémie, HTA non traitée.
  \item L'athérosclérose (plaques d'athérome) est la lésion commune ; elle est favorisée par le LDL‑c, le tabac, l'HTA, le diabète.
  \item Les facteurs se \textbf{potentialisent} : l'effet combiné est supérieur à la somme des effets isolés.
  \item La méthode d'analyse : recenser habitudes + données médicales + antécédents → classer modifiable/non modifiable → compter facteurs majeurs → estimer risque global → proposer plan de prévention ciblé.
\end{itemize}
\end{aretenir}

\begin{exercice}[Application : cas de Mme Atangana]
Mme Atangana, 55 ans, enseignante à la retraite, ne fume pas, boit un verre de vin rouge par jour. Elle marche 30 min trois fois par semaine. Sa tension est de 138/88 mmHg (non traitée). IMC 27 kg/m², tour de taille 88 cm. Glycémie à jeun 0,95 g/L. LDL‑c 1,55 g/L, HDL‑c 0,55 g/L. Sa mère a fait un AVC à 68 ans.
\begin{enumerate}
  \item Lister ses facteurs de risque et les classer.
  \item Estimer son niveau de risque global (faible / modéré / élevé) en justifiant.
  \item Proposer deux mesures de prévention prioritaires.
\end{enumerate}
\end{exercice}

\begin{corrige}[Exercice n°1]
\begin{enumerate}
  \item \textbf{Non modifiables} : âge 55 ans (femme post‑ménopausée). Hérédité : mère AVC à 68 ans (seuil de précocité 65 ans non atteint, donc \textbf{non comptabilisé} comme facteur majeur). \textbf{Modifiables} : HTA limite / normale élevée (138/88), surpoids (IMC 27), surpoids abdominal (tour de taille 88 cm > 80 cm), LDL‑c borderline (1,55 g/L). Tabac = 0, alcool modéré, activité physique correcte (90 min/sem, proche objectif), glycémie normale, HDL‑c bon.
  \item \textbf{Risque faible à modéré} : 1 facteur majeur (surpoids abdominal) + HTA limite + 1 non modifiable (âge). Selon SCORE femme 55 ans non fumeuse, SBP 138, cholestérol total ~2,2 g/L, risque 10 ans ≈ 3–5 \% (zone faible à modérée).
  \item \textbf{Priorités} : (1) Surveillance stricte de la TA et régime DASH (réduction sel < 5 g/j) pour éviter l'installation d'une HTA avérée. (2) Perte de 5 à 10 \% du poids (rééquilibrage alimentaire + augmenter marche à 150 min par semaine) pour faire passer le tour de taille < 80 cm.
\end{enumerate}
\end{corrige}

\coursec{Le système cardiovasculaire : structure et fonctionnement simplifié}

\begin{definition}[Système cardiovasculaire]
Le \textbf{système cardiovasculaire} (ou appareil circulatoire) assure le transport du sang dans l'organisme. Il comprend :
\begin{itemize}
\item le \textbf{cœur} (pompe musculaire à quatre cavités : deux oreillettes, deux ventricules) ;
\item les \textbf{vaisseaux sanguins} : artères (sang riche en O$_2$ vers les organes, sauf artère pulmonaire), veines (sang pauvre en O$_2$ vers le cœur, sauf veines pulmonaires), capillaires (lieu des échanges) ;
\item le \textbf{sang} (milieu de transport : globules rouges, plasma, plaquettes, globules blancs).
\end{itemize}
\end{definition}

\begin{experience}[Observation de la circulation sanguine (schéma animé ou dissection cœur de bœuf)]
\textbf{Matériel :} Cœur de bœuf frais, gants, scalpel, plateau, eau colorée (rouge/bleue), seringue ; ou vidéo animation circulation double.
\textbf{Protocole :}
\begin{enumerate}
\item Identifier les 4 cavités, les gros vaisseaux (aorte, veines caves, artère/veines pulmonaires), les valves.
\item Injecter eau rouge dans veine cave sup $\rightarrow$ suit : oreillette D $\rightarrow$ ventricule D $\rightarrow$ artère pulmonaire.
\item Injecter eau bleue dans veines pulmonaires $\rightarrow$ suit : oreillette G $\rightarrow$ ventricule G $\rightarrow$ aorte.
\end{enumerate}
\textbf{Interprétation :} Le cœur fonctionne comme une \textbf{double pompe} en série : \textbf{petite circulation} (cœur $\rightarrow$ poumons $\rightarrow$ cœur : oxygénation) et \textbf{grande circulation} (cœur $\rightarrow$ organes $\rightarrow$ cœur : nutrition). Les valves assurent l'unidirectionnalité du flux.
\textbf{Conclusion :} Le sang parcourt un circuit fermé, pulsatile, sous pression. La \textbf{fréquence cardiaque} (FC) et le \textbf{volume d'éjection systolique} (VES) déterminent le \textbf{débit cardiaque} (DC = FC $\times$ VES).
\end{experience}

\begin{propriete}[Double circulation et pressions (admis)]
\begin{itemize}
\item \textbf{Petite circulation (pulmonaire) :} Basse pression (PAS artère pulmonaire $\approx$ 25 mmHg). Permet l'échange gazeux sans œdème alvéolaire.
\item \textbf{Grande circulation (systémique) :} Haute pression (PAS aorte $\approx$ 120 mmHg au repos). Assure la perfusion de tous les organes, y compris le cerveau (autoregulation) et les reins (filtration).
\item \textbf{Coronaires :} Artères nourrissant le myocarde. Se remplissent en \textbf{diastole} (relaxation ventriculaire). Toute tachycardie excessive raccourcit la diastole $\rightarrow$ risque d'ischémie coronaire.
\end{itemize}
\end{propriete}

\coursec{La tension artérielle : définition, mesure et valeurs de référence}

\begin{definition}[Tension artérielle (TA)]
La \textbf{tension artérielle} est la pression exercée par le sang sur la paroi des artères. On distingue :
\begin{itemize}
\item \textbf{Pression artérielle systolique (PAS)} : pression maximale lors de la systole ventriculaire (éjection).
\item \textbf{Pression artérielle diastolique (PAD)} : pression minimale lors de la diastole (remplissage).
\item \textbf{Pression artérielle moyenne (PAM)} $\approx$ PAD + (PAS -- PAD)/3. Pression de perfusion effective des organes.
\end{itemize}
Unité : \textbf{millimètre de mercure (mmHg)}. Notation : PAS/PAD (ex. 120/80 mmHg).
\end{definition}

\begin{methode}[Mesure de la TA au brassard (méthode auscultatoire de Korotkoff)]
\begin{enumerate}
\item \textbf{Conditions :} Repos 5--10 min, assis, dos soutenu, bras nu au niveau du cœur, pas de café/tabac/exercice 30 min avant.
\item \textbf{Matériel :} Tensiomètre à mercure (étalon) ou électronique validé (brassard, pas poignet), brassard adapté (largeur $\approx$ 40\% circonférence bras, longueur $\approx$ 80\%).
\item \textbf{Technique :} Brassard 2--3 cm au-dessus du pli du coude. Stéthoscope sur artère brachiale. Gonfler 20--30 mmHg au-dessus de la disparition du pouls radial. Dégonfler lentement (2--3 mmHg/s).
\item \textbf{Lecture :} \textbf{1er bruit de Korotkoff} = PAS. \textbf{Disparition des bruits (5e phase)} = PAD.
\item \textbf{Répétition :} 2 mesures à 1--2 min d'intervalle. Moyenne retenue. Si écart $>$ 5 mmHg $\rightarrow$ 3e mesure.
\end{enumerate}
\end{methode}

\begin{propriete}[Classification de la TA (consultatoire, ESC/ISH 2023 -- admis)]
\begin{center}
\begin{tabular}{|l|c|c|}\hline
\textbf{Catégorie} & \textbf{PAS (mmHg)} & \textbf{PAD (mmHg)} \\ \hline
TA optimale & $<$ 120 & $<$ 80 \\ \hline
TA normale & 120--129 & 80--84 \\ \hline
TA normale élevée & 130--139 & 85--89 \\ \hline
\textbf{HTA Stade 1} & \textbf{140--159} & \textbf{90--99} \\ \hline
\textbf{HTA Stade 2} & \textbf{160--179} & \textbf{100--109} \\ \hline
\textbf{HTA Stade 3} & $\ge$ 180 & $\ge$ 110 \\ \hline
HTA systolique isolée & $\ge$ 140 & $<$ 90 \\ \hline
\end{tabular}
\end{center}
\textbf{Seuils automesure (domicile) :} HTA si moyenne $\ge$ \textbf{135/85 mmHg}. \textbf{MAPA (Mesure Ambulatoire Pression Artérielle) :} 24h, seuil jour 135/85, nuit 120/70, 24h 130/80.
\end{propriete}

\begin{attention}[Erreurs fréquentes de mesure]
\begin{itemize}
\item Brassard trop petit (surestimation) ou trop grand (sous-estimation).
\item Bras non soutenu $\rightarrow$ contraction isométrique $\rightarrow$ surestimation.
\item Parler / croiser les jambes / vessie pleine $\rightarrow$ élévation artefactuelle.
\item \textbf{Effet « blouse blanche » :} TA consultatoire $\ge$ 140/90 mais automesure $<$ 135/85. Fréquent (15--30\%). $\rightarrow$ Confirmer par automesure ou MAPA avant tout traitement.
\end{itemize}
\end{attention}

\coursec{Les principaux accidents cardiovasculaires : infarctus, AVC, hypertension artérielle}

\begin{definition}[Accident cardiovasculaire]
Survenue brutale d'un dysfonctionnement du système cardiovasculaire entraînant une ischémie (manque d'O$_2$) ou une hémorragie d'un organe vital (cœur, cerveau, aorte, membres). Urgence vitale.
\end{definition}

\begin{propriete}[Infarctus du myocarde (IDM) -- admis]
\begin{itemize}
\item \textbf{Mécanisme :} Rupture de plaque d'athérome coronaire $\rightarrow$ thrombose occlusive $\rightarrow$ ischémie $\rightarrow$ nécrose myocardique (irréversible $>$ 20--40 min).
\item \textbf{Signes d'alerte :} Douleur thoracique rétrosternale $\ge$ 20 min, irradiant bras G, mâchoire, dos, oppressante, non calmée par repos/nitrates. Sueurs froides, nausées, angoisse. \textbf{Femme/diabétique :} formes atypiques (dyspnée, fatigue, épigastralgie).
\item \textbf{Conduite :} \textbf{Appeler SAMU (112 / 119) immédiatement}. Aspirine 250--300 mg à croquer si pas contre-indication. Réperfusion urgente (thrombolyse ou angioplastie primaire $<$ 90--120 min).
\end{itemize}
\end{propriete}

\begin{propriete}[Accident Vasculaire Cérébral (AVC) -- admis]
\begin{itemize}
\item \textbf{AVC ischémique (80~\%) :} Thrombose/embolie artère cérébrale (souvent carotide/interne/moyenne). \textbf{AVC hémorragique (20~\%) :} Rupture anévrisme / micro-anévrisme (HTA mal contrôlée) / angiome.
\item \textbf{Signes (mnémotechnique \textbf{VITE} / \textbf{FAST}) :}
    \begin{itemize}
    \item \textbf{V} / \textbf{F} (Face) : Déviation bouche, asymétrie sourire.
    \item \textbf{I} / \textbf{A} (Arms/Bras) : Hémiparésie (un bras tombe).
    \item \textbf{T} / \textbf{S} (Speech/Parole) : Trouble langage (aphasie, dysarthrie).
    \item \textbf{E} / \textbf{T} (Time/Temps) : \textbf{URGENCE ABSOLUE}. Chaque minute = 1,9 millions de neurones perdus.
    \end{itemize}
\item \textbf{Fenêtre thérapeutique :} Thrombolyse IV $<$ 4h30 (idéal $<$ 3h), thrombectomie mécanique $<$ 6--24h selon imagerie. \textbf{Ne jamais donner d'aspirine avant scanner (risque hémorragie).}
\end{itemize}
\end{propriete}

\begin{propriete}[Hypertension Artérielle (HTA) : « tueur silencieux » -- admis]
\begin{itemize}
\item \textbf{HTA essentielle (90--95~\%) :} Multifactorielle (génétique + environnement). Pas de cause unique identifiable.
\item \textbf{HTA secondaire (5--10~\%) :} Rénale (sténose artère rénale, néphropathie), endocrinienne (hyperaldostéronisme, phéochromocytome, thyroïde), médicamenteuse (AINS, corticoïdes, pilule œstro-progestative, réglisse), apnée du sommeil.
\item \textbf{Complications (cibles) :} Cœur (hypertrophie VG, insuffisance cardiaque, IDM), Cerveau (AVC, démence vasculaire), Reins (néphro-angiosclérose, IRC), Yeux (rétinopathie hypertensive), Artères (anévrisme aorte, artériopathie membres inférieurs).
\end{itemize}
\end{propriete}

\coursec{Les facteurs de risque cardiovasculaire}

\begin{definition}[Facteur de risque cardiovasculaire (FRCV)]
Caractéristique (biologique, comportementale, environnementale) associée statistiquement à une augmentation de la probabilité de survenue d'un accident cardiovasculaire. La relation est souvent \textbf{dose-dépendante} et \textbf{multifactorielle} (effet synergique).
\end{definition}

\begin{propriete}[Classification des FRCV (admis -- SCORE2 / Framingham / OMS)]
\begin{center}
\begin{tabularx}{\linewidth}{|l|X|X|}\hline
\textbf{Type} & \textbf{Facteurs} & \textbf{Action possible} \\ \hline
\textbf{Immuables} & Âge ($\ge$ 50 ans H, $\ge$ 60 ans F / ménopause), Sexe (H $>$ F avant 60 ans), Hérédité (antécédent prématuré : IDM/AVC père $<$ 55 ans, mère $<$ 65 ans), Ethnie (antilles, Afrique subsaharienne : HTA plus fréquente/sévère) & \textbf{Non modifiables} $\rightarrow$ Identifient les sujets à dépister \textbf{plus tôt} et \textbf{plus souvent}. \\ \hline
\textbf{Modifiables \textbf{Majeurs}} & \textbf{HTA} (1er FRCV mondial), \textbf{Tabagisme} (actif/passif), \textbf{Diabète} (type 1 \& 2), \textbf{Dyslipidémie} (LDL-c $\uparrow$, HDL-c $\downarrow$, TG $\uparrow$), \textbf{Sédentarité}, \textbf{Surpoids/Obésité} (abdominale), \textbf{Alcool} excessif, \textbf{Stress} chronique, \textbf{Alimentation} déséquilibrée (sel, graisses saturées, sucres, pauvre fibres) & \textbf{Cibles prioritaires} de la prévention. Agir dessus r\'eduit le risque de 80~\% (OMS). \\ \hline
\textbf{Emergents / Marqueurs} & Protéine C-réactive (inflammation), Microalbuminurie, Homocystéine, Fibrinogène, Apnée du sommeil, Pollution air, Bas poids de naissance & En recherche / affinage score risque. \\ \hline
\end{tabularx}
\end{center}
\end{propriete}

\begin{experience}[Calcul de score de risque cardiovasculaire (SCORE2 simplifié -- atelier)]
\textbf{Objectif :} Estimer le risque à 10 ans d'événement CV fatal/non fatal chez un sujet asymptomatique 40--69 ans.
\textbf{Matériel :} Tableaux SCORE2 (région à risque modéré = Cameroun zone Afrique subsaharienne $\approx$ risque modéré/élevé) ou application \textit{HeartScore} / \textit{SCORE2-OP}.
\textbf{Données patient :} H, 55 ans, fumeur, PAS 145 mmHg, Chol total 2,4 g/L, HDL 0,40 g/L, non diabétique.
\textbf{Résultat :} Risque 10 ans $\approx$ 10--15\% (risque \textbf{élevé} $\ge$ 10\% pour 55 ans en zone modérée). $\rightarrow$ Indication de prise en charge intensive (hygiéno-diététique + médicamenteuse si HTA confirmée).
\textbf{Interprétation :} Le score guide la décision thérapeutique (seuils LDL-c, objectif TA). Il ne remplace pas le jugement clinique.
\end{experience}

\coursec{La prévention des accidents cardiovasculaires}

Dans les sections précédentes, nous avons identifié les \textbf{facteurs de risque cardiovasculaire} : certains sont immuables (âge, sexe, hérédité), mais la grande majorité sont \textbf{modifiables} (hypertension artérielle, tabagisme, diabète, dyslipidémie, surpoids, sédentarité, stress, alimentation déséquilibrée, consommation excessive d'alcool). C'est sur ces leviers que repose la prévention. \textbf{Agir sur les facteurs modifiables permet d'éviter jusqu'à 80~\% des infarctus du myocarde et des accidents vasculaires cérébraux (AVC)} (données OMS). Cette section détaille les mesures concrètes, accessibles à tous, pour protéger son cœur et ses vaisseaux tout au long de la vie.

\courssub{Principes généraux de la prévention cardiovasculaire}

\begin{definition}[Prévention cardiovasculaire]
La \textbf{prévention cardiovasculaire} regroupe l'ensemble des mesures — hygiène de vie (alimentation, activité physique, arrêt du tabac et de l'alcool), dépistage régulier des facteurs de risque et suivi médical des pathologies installées — visant à éviter la survenue, la récidive ou l'aggravation des accidents cardiovasculaires (infarctus du myocarde, AVC, insuffisance cardiaque, artériopathie des membres inférieurs).
\end{definition}

On distingue traditionnellement trois niveaux de prévention, illustrés par la pyramide ci-dessous :

\begin{popfigure}
\begin{tikzpicture}[scale=0.9, every node/.style={font=\small}]
% Couleurs
\definecolor{basecol}{HTML}{2E7D32}
\definecolor{midcol}{HTML}{F9A825}
\definecolor{topcol}{HTML}{C62828}
\definecolor{textcol}{HTML}{212121}

% Base (Niveau 1)
\node[draw=basecol, fill=basecol!15, rounded corners=5pt, minimum width=14cm, minimum height=3.5cm, align=center] (base) at (0,0) {
    \textbf{\Large NIVEAU 1 : PRÉVENTION PRIMAIRE}\\[0.4em]
    \textit{(Éviter l'apparition des facteurs de risque et des maladies)}\\[0.6em]
    \begin{minipage}{12.5cm}
    \textbf{Hygiène de vie quotidienne} :\\
    $\bullet$ \textbf{Alimentation équilibrée} : peu de sel, peu de graisses saturées, riches en fruits/légumes/fibres (ndolé, folon, gombo, banane plantain, avocat\ldots)\\
    $\bullet$ \textbf{Activité physique régulière} : $\ge$ 30 min marche rapide / jour, 5 fois/semaine\\
    $\bullet$ \textbf{Zéro tabac} (actif et passif) — \textbf{Alcool} : modération ou abstinence\\
    $\bullet$ \textbf{Gestion du stress} \& \textbf{Sommeil suffisant} (7--8 h/nuit)\\
    $\bullet$ \textbf{Poids santé} : IMC 18,5--25, tour de taille $<$ 94 cm (H) / 80 cm (F)
    \end{minipage}
};

% Milieu (Niveau 2)
\node[draw=midcol, fill=midcol!15, rounded corners=5pt, minimum width=11cm, minimum height=3cm, align=center] (mid) at (0,4) {
    \textbf{\Large NIVEAU 2 : DÉPISTAGE PRÉCOCE}\\[0.4em]
    \textit{(Détecter et traiter les facteurs de risque asymptomatiques)}\\[0.6em]
    \begin{minipage}{9.5cm}
    $\bullet$ \textbf{Tension artérielle} : $\ge$ 1 fois/an dès 18 ans (automesure validée)\\
    $\bullet$ \textbf{Glycémie à jeun} / \textbf{Profil lipidique} : selon facteurs de risque\\
    $\bullet$ \textbf{Poids / IMC / Tour de taille} : contrôle régulier\\
    $\bullet$ \textbf{Bilan médical} : ECG, recherche protéinurie si HTA/diabète
    \end{minipage}
};

% Sommet (Niveau 3)
\node[draw=topcol, fill=topcol!15, rounded corners=5pt, minimum width=8cm, minimum height=2.5cm, align=center] (top) at (0,7.3) {
    \textbf{\Large NIVEAU 3 : SUIVI MÉDICAL}\\[0.4em]
    \textit{(Prévenir récidives \& complications chez le malade connu)}\\[0.6em]
    \begin{minipage}{6.5cm}
    $\bullet$ \textbf{Observance stricte} des traitements (HTA, diabète, dyslipidémie)\\
    $\bullet$ \textbf{Automesure tensionnelle} (règle des 3 : 3 matins, 3 soirs, 3 jours)\\
    $\bullet$ \textbf{Consultations régulières} (médecin, cardiologue)\\
    $\bullet$ \textbf{Éducation thérapeutique} du patient
    \end{minipage}
};

% Flèches et étiquettes
\draw[->, thick, popdark] (base.north) -- (mid.south) node[midway, right=0.2cm, align=left, text width=4cm] {\small Si facteur de risque détecté $\rightarrow$ prise en charge};
\draw[->, thick, popdark] (mid.north) -- (top.south) node[midway, right=0.2cm, align=left, text width=4cm] {\small Si pathologie confirmée $\rightarrow$ suivi spécialisé};

% Titre général
\node[above=0.3cm of top.north, font=\bfseries\Large, text=popdark] {PYRAMIDE DE LA PRÉVENTION CARDIOVASCULAIRE};
\end{tikzpicture}
\legende{Pyramide des trois niveaux de prévention cardiovasculaire. La base large (prévention primaire) concerne toute la population ; le dépistage (niveau 2) cible les adultes et les sujets à risque ; le suivi médical (niveau 3) concerne les patients déjà diagnostiqués (hypertendus, diabétiques, coronariens).}
\end{popfigure}

\courssub{La prévention alimentaire : l'assiette santé camerounaise}

L'alimentation est le premier médicament cardiovasculaire. Au Cameroun, la transition nutritionnelle (urbanisation, \textit{fast-food}, cubes d'assaisonnement riches en sel, huiles de friture réutilisées) fait exploser l'hypertension et le diabète. Revenons à nos bases alimentaires, riches en fibres, potassium, antioxydants et acides gras insaturés.

\begin{propriete}[Effet de l'alimentation sur le risque cardiovasculaire (admis)]
Une alimentation de type \textbf{méditerranéen} ou \textbf{DASH} (Dietary Approaches to Stop Hypertension), adaptée au contexte local — riche en légumes, fruits, céréales complètes, légumineuses, poisson, pauvre en sel ($<$ 5 g/jour), graisses saturées et sucres ajoutés — réduit significativement la tension artérielle (baisse de 5 à 10 mmHg de la PAS), le LDL-cholestérol et l'incidence des infarctus et AVC.
\end{propriete}

\begin{methode}[Composer son assiette santé au quotidien]
\begin{enumerate}
\item \textbf{Remplir la moitié de l'assiette de légumes} (crus ou cuits peu gras) : ndolé, folon (feuilles de gombo), gombo, épinards, chou, carottes, tomates, concombre\ldots
\item \textbf{Un quart de féculents complets ou tubercules} : riz complet, maïs, mil, sorgho, igname, macabo, patate douce, banane plantain (bouillie ou vapeur, pas frite).
\item \textbf{Un quart de protéines maigres} : poisson (capitaine, carpe, maquereau — riches en oméga-3), poulet sans peau, œufs, légumineuses (haricots, niébé, soja), viande rouge maigre (2 fois/semaine max).
\item \textbf{Un fruit de saison} en dessert : papaye, ananas, pastèque, mangue, banane, avocat (modération : riche en graisses saines mais calorique).
\item \textbf{Assaisonner malin} : oignon, ail, gingembre, persil, céleri, piment, jus de citron \textbf{à la place des cubes} (Maggi, Jumbo\ldots) et du sel de table. Ne pas resaler à table.
\item \textbf{Boire de l'eau} (1,5--2 L/jour), limiter les boissons sucrées (jus industriels, sodas) et l'alcool.
\end{enumerate}
\end{methode}

\begin{popfigure}
\begin{tikzpicture}[scale=0.85, every node/.style={font=\small}]
% Assiette (cercle)
\draw[thick, popdark] (0,0) circle (6cm);
% Divisions
\draw[thick, popdark] (0,0) -- (0,6);
\draw[thick, popdark] (0,0) -- (-5.2,-3);
\draw[thick, popdark] (0,0) -- (5.2,-3);

% Étiquettes secteurs
\node[align=center, text width=5cm] at (0,2.5) {
    \textbf{\large 50\% LÉGUMES}\\
    \textit{(fibres, K, Mg, antioxydants)}\\
    Ndolé, folon, gombo, épinards, chou,\\
    carottes, tomates, concombre, aubergine\ldots
};
\node[align=center, text width=4.5cm] at (-3,-2) {
    \textbf{\large 25\% FÉCULENTS}\\
    \textit{(énergie, fibres si complets)}\\
    Riz complet, maïs, mil, sorgho,\\
    igname, macabo, patate douce,\\
    banane plantain (vapeur/bouillie)
};
\node[align=center, text width=4.5cm] at (3,-2) {
    \textbf{\large 25\% PROTÉINES}\\
    \textit{(construction, oméga-3 si poisson)}\\
    Poisson (capitaine, maquereau),\\
    poulet sans peau, œufs,\\
    haricots, niébé, soja,\\
    viande rouge maigre (2$\times$/sem)
};
% Fruit à côté
\node[draw=popgreen, fill=popgreen!10, rounded corners=5pt, minimum width=3.5cm, minimum height=2.5cm, align=center] at (8.5, 3) {
    \textbf{+ 1 FRUIT}\\
    \textit{(vitamines, fibres)}\\
    Papaye, ananas, pastèque,\\
    mangue, banane, avocat ($\frac{1}{2}$)
};
\draw[->, thick, popgreen] (6,3) -- (7.5,3);
% Eau
\node[draw=popblue, fill=popblue!10, rounded corners=5pt, minimum width=3.5cm, minimum height=1.5cm, align=center] at (8.5, -1) {
    \textbf{EAU}\\
    1,5 -- 2 L / jour
};
\draw[->, thick, popblue] (6,-1.5) -- (7.5,-1.5);
% Interdits
\node[draw=poppink, fill=poppink!10, rounded corners=5pt, minimum width=5cm, minimum height=2.5cm, align=center] at (-9, -4) {
    \textbf{À LIMITER / ÉVITER}\\
    $\times$ Sel ajouté \& cubes (Maggi, Jumbo)\\
    $\times$ Fritures (beignets, poisson frit, plantain frit)\\
    $\times$ Graisses animales (gras de bœuf, porc, peau poulet)\\
    $\times$ Charcuteries, abats\\
    $\times$ Boissons sucrées, alcool excessif
};
\draw[->, thick, poppink] (-6,-3.5) -- (-7.5,-3.5);

% Titre
\node[above=0.3cm of {(0,6)}, font=\bfseries\Large, text=popdark] {L'ASSIETTE SANTÉ CAMEROUNAISE};
\end{tikzpicture}
\legende{Repères visuels pour composer un repas protecteur du cœur et des vaisseaux. La moitié de l'assiette en légumes verts locaux (ndolé, folon, gombo\ldots), un quart en féculents complets ou tubercules, un quart en protéines maigres (poisson de préférence), un fruit de saison et de l'eau. On bannit le resalage et les cubes d'assaisonnement industriels.}
\end{popfigure}

\begin{exemplebox}[Réduire le sel : du diagnostic à l'action]
\textbf{Situation :} M. K., 52 ans, chef de service à Douala, consulte pour une tension à 165/100 mmHg. L'anamnèse alimentaire révèle : 3 cubes Maggi par repas, sauce tomate industrielle, resalage systématique à table, beignets haricots au petit-déjeuner, poisson frit 3 fois/semaine. Apport sodique estimé $\approx$ 12--14 g/jour (sel $\approx$ 30--35 g/jour).

\textbf{Objectif :} Ramener l'apport en sel $<$ 5 g/jour (soit Na $<$ 2 g/jour) sans médicament dans un premier temps (HTA stade 1, pas de complication).

\textbf{Mesures concrètes proposées :}
\begin{enumerate}
\item \textbf{Cuisine} : Arrêt total des cubes. Utilisation d'oignon, ail, gingembre, persil, céleri, piment, jus de citron, tomates fraîches.
\item \textbf{À table} : Retirer la salière. Ne pas resaler.
\item \textbf{Produits transformés} : Remplacer le pain blanc industriel par du pain maison peu salé ou de la banane plantain bouillie. Supprimer les beignets (riches en sel caché et graisses \textit{trans}).
\item \textbf{Cuisson} : Poisson grillé ou en sauce (sans cube) au lieu de frit. Huile rouge non chauffée (source de vitamine E et tocotriénols) en petite quantité.
\item \textbf{Suivi} : Automesure tensionnelle (règle des 3) + consultation à 1 mois.
\end{enumerate}

\textbf{Résultat attendu (données cliniques) :} Une réduction de 10 g à 5 g de sel/jour abaisse la PAS de \textbf{5 à 10 mmHg} (environ \textbf{0,7 à 1,3 cmHg}) chez un sujet hypertendu sensible au sel, effet comparable à un antihypertenseur en monothérapie.
\end{exemplebox}

\begin{attention}[Piège fréquent : les « faux amis » alimentaires]
\begin{itemize}
\item \textbf{« L'huile rouge, c'est naturel donc illimité. »} \textbf{Faux.} L'huile de palme non raffinée contient des antioxydants précieux, mais reste une graisse saturée à \textbf{50~\%} (acide palmitique). À consommer avec modération (1--2 c. à soupe/jour), non chauffée.
\item \textbf{« L'avocat, c'est du bon gras, j'en mange 2 par jour. »} \textbf{Attention.} L'avocat est riche en acides gras mono-insaturés (excellents), mais très calorique ($\approx$ 160 kcal/100 g). Un avocat moyen = 300--400 kcal. $\rightarrow$ \textbf{$\frac{1}{2}$ avocat/jour max} en prévention cardiovasculaire.
\item \textbf{« Je bois du jus de fruit naturel, c'est sain. »} \textbf{Piège.} Un verre de jus d'ananas/mangue = 3--4 fruits pressés, fibres éliminées, index glycémique élevé. $\rightarrow$ \textbf{Manger le fruit entier}, boire de l'eau.
\end{itemize}
\end{attention}

\courssub{L'activité physique : le médicament sous-estimé}

\begin{propriete}[Activité physique et risque cardiovasculaire (admis)]
Une activité physique d'endurance modérée (\textbf{marche rapide, vélo, natation, danse, travail agricole soutenu}) réalisée \textbf{au moins 30 minutes par jour, 5 jours par semaine} (soit 150 min/semaine), réduit la mortalité cardiovasculaire de \textbf{20 à 30~\%}, abaisse la PAS de 4 à 9 mmHg, augmente le HDL-cholestérol, améliore la sensibilité à l'insuline et favorise la perte de poids.
\end{propriete}

\begin{methode}[Prescrire (ou s'auto-prescrire) l'activité physique en toute sécurité]
\begin{enumerate}
\item \textbf{Évaluer le niveau actuel} : Questionnaires (ex. IPAQ court) ou test de marche de 6 minutes.
\item \textbf{Fixer un objectif progressif} : Si sédentaire total $\rightarrow$ commencer par 10 min marche/jour, augmenter de 5 min/semaine jusqu'à 30 min.
\item \textbf{Choisir une activité plaisante et accessible} : Marche aller-retour au travail/marché, danse traditionnelle, football léger, jardinage, vélo.
\item \textbf{Respecter la « règle du discours »} : En effort modéré, on peut parler mais pas chanter. Fréquence cardiaque cible $\approx$ 50--70\% FC max (FC max $\approx$ 220 -- âge).
\item \textbf{Intégrer le renforcement musculaire} : 2 séances/semaine (port de charges, pompes, squats, travail aux champs) pour le métabolisme de base.
\item \textbf{Lutter contre la sédentarité} : Se lever toutes les heures (bureau, moto-taxi, télévision). « Le mieux est l'ennemi du bien » : \textbf{bouger un peu vaut mieux que ne rien faire}.
\end{enumerate}
\end{methode}

\begin{savaistu}[Le cœur au travail]
Au Cameroun, les métiers physiques (agriculteurs, porteurs, maçons, pêcheurs) ont historiquement une faible prévalence d'HTA. Mais l'urbanisation crée une \textbf{« sédentarité active »} : on marche pour aller au travail, mais on reste assis 8--10 h/jour (bureau, moto-taxi, écrans). \textbf{La sédentarité (temps assis) est un facteur de risque indépendant de l'activité physique.} Même un sportif qui reste assis 10 h/jour a un risque accru. \textbf{Solution :} pause active de 2--3 min toutes les heures (étirements, marche, escaliers).
\end{savaistu}

\courssub{Lutte contre le tabac, l'alcool, le stress et le sommeil}

\begin{experience}[Effet aigu du tabac sur la tension artérielle et la fréquence cardiaque]
\textbf{Matériel :} Tensiomètre brassard validé, oxymètre de pouls, 2 volontaires fumeurs (acceptant de fumer une cigarette), 2 volontaires non-fumeurs (témoins).
\textbf{Protocole :}
\begin{enumerate}
\item Repos 10 min assis. Mesure TA et FC basales (moyenne 3 mesures).
\item Volontaires fumeurs : fumer 1 cigarette en 5 min. Témoins : rester assis.
\item Mesure TA et FC immédiates, puis à 10, 20, 30 min post-exposition.
\end{enumerate}
\textbf{Observations types :}
\begin{center}
\begin{tabular}{|l|c|c|c|c|}\hline
 & \textbf{Basal} & \textbf{Immédiat} & \textbf{10 min} & \textbf{30 min} \\ \hline
\textbf{Fumeur (TA sys)} & 125 mmHg & 145 mmHg & 138 mmHg & 130 mmHg \\ \hline
\textbf{Fumeur (FC)} & 72 bpm & 95 bpm & 85 bpm & 78 bpm \\ \hline
\textbf{Témoin (TA sys)} & 122 mmHg & 123 mmHg & 121 mmHg & 122 mmHg \\ \hline
\textbf{Témoin (FC)} & 70 bpm & 71 bpm & 70 bpm & 69 bpm \\ \hline
\end{tabular}
\end{center}
\textbf{Interprétation :} La nicotine provoque une \textbf{libération d'adrénaline} $\rightarrow$ vasoconstriction ($\uparrow$ résistances périphériques) + $\uparrow$ FC + $\uparrow$ contractilité $\rightarrow$ \textbf{pic tensionnel et tachycardie}. Répété 10--20 fois/jour, ce stress hémodynamique chronique rigidifie les artères et accélère l'athérosclérose.
\textbf{Conclusion :} Chaque cigarette est une agression vasculaire aiguë. L'arrêt du tabac fait baisser le risque cardiovasculaire de \textbf{50~\% dès 1 an}, et rejoint celui d'un non-fumeur à 5--10 ans.
\end{experience}

\begin{aretenir}[Tabac, alcool, stress, sommeil : 4 piliers à ne pas négliger]
\begin{itemize}
\item \textbf{Tabac} : \textbf{Zéro cigarette}. Pas de « cigarette light » (inhalation plus profonde). Éviter le tabagisme passif (foyer, bars, transports). Aide : ligne verte du MINSANTE (\textbf{1510}), substituts nicotiniques, thérapies cognitivo-comportementales.
\item \textbf{Alcool} : \textbf{Modération stricte} : $\le$ 2 verres/jour (H), $\le$ 1 verre/jour (F), pas tous les jours. 1 verre = 10 g alcool pur (25 cl bière 4\%, 10 cl vin 12\%, 3 cl whisky 40\%). \textbf{Binge drinking} (ivresse épisodique) = risque aigu d'AVC hémorragique et d'arythmie.
\item \textbf{Stress chronique} : Active l'axe corticosurrénalien (cortisol) et le système sympathique $\rightarrow$ HTA, prise de poids abdominale, troubles du sommeil. \textbf{Gestion} : respiration diaphragmatique (cohérence cardiaque 5-5-5), prière/méditation, loisirs, soutien social, aide psychologique si besoin.
\item \textbf{Sommeil} : $<$ 6 h ou $>$ 9 h/nuit $\rightarrow$ risque cardiovasculaire accru. \textbf{Objectif :} 7--8 h, rythme régulier, obscurité, pas d'écran 1 h avant coucher.
\end{itemize}
\end{aretenir}

\courssub{Le dépistage : détecter pour agir à temps}

L'hypertension artérielle est le « tueur silencieux » : \textbf{50~\% des hypertendus s'ignorent} (enquête STEPS Cameroun 2023). Le dépistage est l'arme majeure de la prévention secondaire.

\begin{methode}[Organiser son dépistage cardiovasculaire]
\begin{enumerate}
\item \textbf{Tension artérielle (TA)} : \textbf{Au moins 1 fois par an dès 18 ans.} Si TA $\ge$ 130/85 mmHg $\rightarrow$ contrôle à 3--6 mois. Si TA $\ge$ 140/90 mmHg $\rightarrow$ confirmation (automesure ou MAPA) et prise en charge.
\item \textbf{Automesure tensionnelle (règle des 3)} : 3 mesures le matin (avant café/traitement), 3 mesures le soir (avant coucher), pendant 3 jours consécutifs. Moyenne des 18 valeurs (en excluant la 1ère journée). Seuil normalité : $<$ 135/85 mmHg.
\item \textbf{Glycémie à jeun} : Dès 40 ans, ou avant si surpoids/obésité, antécédents familiaux diabète, HTA, dyslipidémie, syndrome des ovaires polykystiques, grossesse antérieure macrosomie. Seuil diabète : $\ge$ 1,26 g/L (2 fois) ou HbA1c $\ge$ 6,5~\%.
\item \textbf{Profil lipidique} : Cholestérol total, HDL, LDL, triglycérides. Selon facteurs de risque (âge, HTA, tabac, diabète, antécédents familiaux précoces). Objectif LDL : $<$ 1,6 g/L (risque modéré), $<$ 1,0 g/L (risque élevé), $<$ 0,7 g/L (très haut risque / secondaire).
\item \textbf{Poids / IMC / Tour de taille} : À chaque consultation. IMC $\ge$ 25 = surpoids, $\ge$ 30 = obésité. Tour de taille $\ge$ 94 cm (H) / 80 cm (F) = obésité abdominale (risque métabolique majoré).
\item \textbf{Bilan rénal (créatininémie, protéinurie)} et \textbf{ECG} : Si HTA ou diabète confirmés.
\end{enumerate}
\end{methode}

\begin{popfigure}
\begin{tikzpicture}
\begin{axis}[
    width=\linewidth, height=8cm,
    xlabel={Temps (mois)}, ylabel={Pression artérielle systolique (mmHg)},
    xmin=0, xmax=6, ymin=110, ymax=170,
    xtick={0,1,2,3,4,5,6}, ytick={110,120,130,140,150,160,170},
    grid=major, grid style={popline},
    legend style={at={(0.98,0.98)}, anchor=north east, fill=white, draw=popdark},
    axis line style={popdark}, tick style={popdark},
    label style={font=\small}, tick label style={font=\small},
    legend style={font=\small}
]
% Courbe sans intervention (témoin simulé)
\addplot[poppink, thick, mark=*, mark size=2] coordinates {
    (0,158) (1,160) (2,162) (3,165) (4,163) (5,166) (6,168)
};
\addlegendentry{Sans intervention (évolution naturelle)}

% Courbe avec mesures hygiéno-diététiques
\addplot[popgreen, thick, mark=square*, mark size=2] coordinates {
    (0,158) (1,148) (2,142) (3,138) (4,135) (5,132) (6,130)
};
\addlegendentry{Avec mesures hygiéno-diététiques}

% Ligne cible
\addplot[popblue, dashed, thick, domain=0:6] {135};
\addlegendentry{Cible automesure ($<$ 135 mmHg)}

% Ligne HTA
\addplot[poporange, dashed, thick, domain=0:6] {140};
\addlegendentry{Seuil HTA consultatoire (140 mmHg)}
\end{axis}
\end{tikzpicture}
\legende{Évolution simulée de la pression artérielle systolique (PAS) sur 6 mois chez un patient hypertendu stade 1 (PAS initiale 158 mmHg). \textbf{Courbe rose} : absence de modification du mode de vie (aggravation progressive). \textbf{Courbe verte} : adoption des mesures (sel $<$ 5 g/j, marche 30 min/j, arrêt tabac, perte 5 kg, DASH). La cible d'automesure ($<$ 135 mmHg) est atteinte vers le 3\textsuperscript{e} mois. Source : données synthétiques inspirées d'essais cliniques (PREMIER, DASH-Sodium).}
\end{popfigure}

\begin{exoresolu}[Interpréter une automesure et proposer un plan d'action]
\textbf{Énoncé :} Mme B., 48 ans, vendeuse au marché Mfoundi, réalise son automesure tensionnelle sur 3 jours (règle des 3). Elle n'a pas de traitement. Résultats (moyennes des 2 dernières mesures de chaque créneau) :
\begin{center}
\begin{tabular}{|c|c|c|c|}\hline
\textbf{Jour} & \textbf{Matin (M)} & \textbf{Soir (S)} & \textbf{Moyenne jour} \\ \hline
J1 (exclu) & 148/94 & 152/96 & -- \\ \hline
J2 & 142/90 & 146/92 & 144/91 \\ \hline
J3 & 138/88 & 140/90 & 139/89 \\ \hline
J4 & 136/86 & 138/88 & 137/87 \\ \hline
\end{tabular}
\end{center}
\textbf{Questions :}
\begin{enumerate}
\item Calculer la moyenne globale des 6 mesures valides (J2--J4). Interpréter au regard du seuil d'automesure (135/85 mmHg).
\item Identifier 3 facteurs de risque probables dans son contexte (marché, alimentation, stress, sédentarité\ldots).
\item Proposer un plan de prévention concret, réaliste et priorisé (5 actions max).
\end{enumerate}

\tcblower
\textbf{Solution détaillée :}
\begin{enumerate}
\item \textbf{Moyenne globale} : PAS = (142+146+138+140+136+138)/6 = \textbf{140 mmHg}. PAD = (90+92+88+90+86+88)/6 = \textbf{89 mmHg}.
\textbf{Interprétation} : Moyenne 140/89 mmHg $>$ seuil automesure 135/85 mmHg $\rightarrow$ \textbf{HTA confirmée (stade 1)}. Nécessite consultation médicale pour bilan (recherche complications, bilan biologique) et discussion prise en charge (hygiéno-diététique d'abord 3 mois si pas de complication/antécédent).
\item \textbf{Facteurs de risque probables} :
    \begin{itemize}
    \item \textbf{Sel excessif} : cuisine au cube (Maggi/Jumbo), sauces industrielles, resalage, beignets, poisson fumé/salé.
    \item \textbf{Sédentarité professionnelle} : assise 8--10 h/jour derrière son étal, déplacements en moto-taxi.
    \item \textbf{Stress chronique} : instabilité revenus, bruit, chaleur, charge mentale (famille, crédit tontine).
    \item \textbf{Sommeil court/fragmenté} : lever 4 h 30, coucher tardif.
    \end{itemize}
\item \textbf{Plan d'action priorisé (négocié avec la patiente)} :
    \begin{enumerate}
    \item \textbf{Sel} : Arrêt cubes $\rightarrow$ épices locales (penja, djansang, ail, gingembre, citron). Ne pas resaler. Rincer poisson fumé. Objectif $<$ 5 g/j.
    \item \textbf{Activité} : Marche 20 min aller-retour domicile-marché (ou 10 min pause midi + 10 min soir) $\rightarrow$ viser 30 min/j. Monter escaliers au lieu d'ascenseur si bâtiment.
    \item \textbf{Poids} : Perte 3--5 kg sur 3 mois (réduction portions féculents, supprimer beignets/sodas, fruit au lieu de gâteau).
    \item \textbf{Stress/Sommeil} : Respiration 5 min à l'ouverture de l'étal (cohérence cardiaque). Coucher 22 h max. Sieste 20 min si possible.
    \item \textbf{Suivi} : Automesure mensuelle (3 jours). RDV médecin à 3 mois pour réévaluation. Si PAS $>$ 140 persistante $\rightarrow$ discussion traitement.
    \end{enumerate}
\end{enumerate}
\end{exoresolu}

\courssub{Suivi médical et observance : la chronicité ne prend pas de vacances}

\begin{attention}[Erreur majeure : « Je me sens bien, j'arrête mon traitement. »]
\textbf{C'est l'erreur n°1 en cardiologie.} L'hypertension, le diabète, la dyslipidémie sont des \textbf{maladies chroniques silencieuses}. L'absence de symptômes ne signifie pas l'absence de maladie.
\begin{itemize}
\item \textbf{Arrêter l'anti-hypertenseur} $\rightarrow$ « effet rebond » : TA remonte souvent \textbf{au-dessus} des valeurs initiales en quelques jours $\rightarrow$ risque d'AVC, d'œdème aigu du poumon, d'infarctus.
\item \textbf{Le traitement se prend à vie}, sous contrôle médical (adaptation posologie, changement molécule si effets secondaires).
\item \textbf{Observance} = prendre son traitement \textbf{tous les jours, à la même heure, sans oublier}. Astuces : pilulier hebdomadaire, alarme téléphone, associer à un rituel (brossage dents), emporter réserve en voyage.
\item \textbf{Automesure} = boussole du traitement. Noter les valeurs dans un carnet ou application (ex. \textit{MyHTA}, \textit{Cardiocheck}) et l'apporter à chaque consultation.
\end{itemize}
\end{attention}

\begin{propriete}[Rôle de l'élève : relais de santé communautaire (admis)]
L'élève de Première, formé aux gestes de prévention, devient un \textbf{acteur de santé publique} dans son entourage :
\begin{itemize}
\item \textbf{Famille} : Inciter parents/grands-parents à mesurer leur TA (offrir un tensiomètre brassard validé), cuisiner moins salé, marcher ensemble le week-end.
\item \textbf{École} : Animer un club « Cœur en santé », organiser une journée « Zéro sel / Zéro tabac » le 29 septembre (Journée mondiale du cœur).
\item \textbf{Communauté} : Relayer les messages de la Campagne nationale de lutte contre les MNT (Maladies Non Transmissibles) du MINSANTE, signaler les signes d'alerte (douleur thoracique, paralysie brutale, trouble de la parole $\rightarrow$ appeler les urgences / SAMU 112 ou 119 ou se rendre au centre de santé le plus proche).
\end{itemize}
\end{propriete}

\begin{savaistu}[Journée Mondiale du Cœur — 29 Septembre]
Instituée par la Fédération Mondiale du Cœur (World Heart Federation) et l'OMS, la \textbf{Journée Mondiale du Cœur} (29 septembre) vise à sensibiliser sur les maladies cardiovasculaires, première cause de mortalité mondiale (17,9 millions de décès/an, dont 75~\% dans les pays à revenu faible ou intermédiaire). Au Cameroun, le MINSANTE et la Société Camerounaise de Cardiologie organisent des dépistages gratuits, marches sportives, conférences dans les lycées et marchés. Thème récent : \textbf{« Utilisez votre cœur pour agir »} — chaque individu a le pouvoir de changer son mode de vie et d'influencer son entourage. \textbf{Toi aussi, tu as un rôle à jouer !}
\end{savaistu}

\coursec{Synthèse et évaluation des acquis}

\courssub{Synthèse : ton plan d'action personnel « Cœur en santé »}

\begin{exercice}[Bilan de compétences — Prévention cardiovasculaire]
\textbf{Consigne :} À partir de ton mode de vie actuel, remplis le tableau ci-dessous en toute honnêteté. Pour chaque item, coche ta situation et écris \textbf{une action concrète} que tu t'engages à mettre en place cette semaine.

\begin{center}
\begin{tabularx}{\linewidth}{|l|X|X|}\hline
\textbf{Domaine} & \textbf{Mes habitudes actuelles} & \textbf{Mon engagement concret (SMART : Spécifique, Mesurable, Atteignable, Réaliste, Temporel)} \\ \hline
\textbf{Sel / Cubes} & \textit{Ex: 2 cubes/repas, resalage table} & \textit{Ex: 0 cube dès demain ; citron/ail/gingembre ; salière retirée} \\ \hline
\textbf{Légumes / Fruits} & & \\ \hline
\textbf{Activité physique} & & \\ \hline
\textbf{Tabac / Alcool} & & \\ \hline
\textbf{Stress / Sommeil} & & \\ \hline
\textbf{Dépistage (TA, glycémie)} & & \\ \hline
\textbf{Relais familial} & & \\ \hline
\end{tabularx}
\end{center}
\textbf{Partage ton engagement avec un camarade (binôme) : vous vous relancez mutuellement dans 1 semaine.}
\end{exercice}

\begin{corrige}[Exercice n°1 — Bilan de compétences]
\textit{Corrigé indicatif (chaque élève a son propre bilan). L'important est la formulation SMART.}
\begin{itemize}
\item \textbf{Mal formulé :} « Je vais manger mieux. » $\rightarrow$ \textbf{Bien formulé :} « Je mange 200 g de légumes (ndolé/folon) au déjeuner et 1 fruit (papaye) au dîner, 5 soirs sur 7 cette semaine. »
\item \textbf{Mal formulé :} « Je vais faire du sport. » $\rightarrow$ \textbf{Bien formulé :} « Je marche 30 min (aller-retour lycée-domicile) les lundis, mercredis, vendredis, samedis, dimanches. »
\item \textbf{Mal formulé :} « Je vais arrêter le sel. » $\rightarrow$ \textbf{Bien formulé :} « Je ne rachète pas de cubes Maggi ; j'achète du persil, céleri, piment, citron pour cuisiner ; je ne pose pas la salière sur la table. »
\end{itemize}
\textbf{Critère de réussite :} Au moins 3 engagements tenus sur 5 à J+7. Si échec $\rightarrow$ analyser l'obstacle (goût, temps, fatigue, entourage) et adapter l'action (plus petite, plus plaisante, soutien familial).
\end{corrige}

\begin{aretenir}[Les 5 commandements de la prévention cardiovasculaire au quotidien]
\begin{enumerate}
\item \textbf{Assiette} : $\frac{1}{2}$ légumes locaux, $\frac{1}{4}$ féculents complets, $\frac{1}{4}$ protéines maigres (poisson $>$ viande), 1 fruit, eau. \textbf{0 cube, 0 resalage.}
\item \textbf{Mouvement} : 30 min marche rapide / jour (ou équivalent). Bouger toutes les heures.
\item \textbf{Substances} : \textbf{Zéro tabac} (actif/passif). Alcool : 0 ou modéré (max 1--2 verres/j, pas quotidien).
\item \textbf{Veille} : TA 1 fois/an (automesure validée). Poids/IMC/tour de taille réguliers. Bilan sanguin si facteur de risque.
\item \textbf{Chronicité} : Si traitement $\rightarrow$ \textbf{à vie, tous les jours, automesure mensuelle, RDV médecin 2--4 fois/an}. Ne jamais arrêter seul.
\end{enumerate}
\end{aretenir}

\textbf{Tu as maintenant toutes les clés pour protéger ton cœur et celui de tes proches.} La prévention cardiovasculaire n'est pas une contrainte, c'est un \textbf{investissement quotidien pour une vie longue, active et libre}. Le cœur qui bat en toi est le moteur de tes rêves — prends-en soin dès aujourd'hui.

\coursec{Synthèse et évaluation des acquis}

Après avoir exploré les piliers de la prévention — alimentation équilibrée, activité physique régulière et dépistage systématique — nous arrivons au terme de cette leçon. Cette dernière section a pour but de \textbf{consolider l'ensemble des connaissances} acquises, de \textbf{vérifier votre capacité à les mobiliser} face à des situations concrètes (mesures tensionnelles, profils de risque, cas cliniques) et de vous préparer aux exigences du Probatoire, notamment la rédaction de messages de sensibilisation argumentés. Comme le dirait un clinicien expérimenté : « Connaître les facteurs de risque ne suffit pas ; il faut savoir les repérer chez un patient et agir avant l'accident. »

\courssub{Bilan structuré : les quatre piliers de la leçon}

Pour organiser votre révision, nous vous proposons une carte mentale synthétique qui articule les quatre grands blocs de cette leçon : le système cardiovasculaire (substrat anatomique et physiologique), les accidents majeurs (pathologies), les facteurs de risque (étiologie) et la prévention (prise en charge).

\begin{popfigure}
\begin{tikzpicture}[
    mindmap,
    grow cyclic,
    every node/.style={concept, execute at begin node=\hskip0pt},
    concept color=popblue!80,
    level 1/.append style={level distance=4.5cm, sibling angle=90},
    level 2/.append style={level distance=3.2cm, sibling angle=45},
    level 3/.append style={level distance=2.8cm, sibling angle=30},
    popblue/.style={concept color=popblue},
    poporange/.style={concept color=poporange},
    popgreen/.style={concept color=popgreen},
    poppurple/.style={concept color=poppurple},
    poppink/.style={concept color=poppink},
    popgold/.style={concept color=popgold},
]
\node[concept, text width=3.2cm, align=center, font=\bfseries\large] {Lutte contre les\\ accidents cardiovasculaires}
    child[popblue] {
        node[concept, text width=2.8cm, align=center] {1. Système cardiovasculaire}
        child { node[concept, text width=2.2cm, align=center] {Cœur : pompe double (cavités, valves)} }
        child { node[concept, text width=2.2cm, align=center] {Vaisseaux : artères, veines, capillaires} }
        child { node[concept, text width=2.2cm, align=center] {Double circulation : petite \& grande} }
        child { node[concept, text width=2.2cm, align=center] {Tension artérielle : $P_{sys}$ / $P_{dias}$ (mmHg)} }
    }
    child[poporange] {
        node[concept, text width=2.8cm, align=center] {2. Accidents majeurs}
        child { node[concept, text width=2.2cm, align=center] {Infarctus du myocarde (IDM) : nécrose, douleur thoracique} }
        child { node[concept, text width=2.2cm, align=center] {AVC (ischémique/hémorragique) : signes FAST} }
        child { node[concept, text width=2.2cm, align=center] {HTA : tueur silencieux ($>140/90$)} }
        child { node[concept, text width=2.2cm, align=center] {Conduite à tenir : 112 / 113} }
    }
    child[popgreen] {
        node[concept, text width=2.8cm, align=center] {3. Facteurs de risque}
        child { node[concept, text width=2.2cm, align=center] {Non modifiables : âge, sexe, hérédité} }
        child { node[concept, text width=2.2cm, align=center] {Modifiables : tabac, alimentation, sédentarité, alcool, stress} }
        child { node[concept, text width=2.2cm, align=center] {Facteurs métaboliques : HTA, diabète, dyslipidémie} }
        child { node[concept, text width=2.2cm, align=center] {Effet cumulatif : score de risque (Framingham/SCORE2)} }
    }
    child[poppurple] {
        node[concept, text width=2.8cm, align=center] {4. Prévention (3 piliers)}
        child { node[concept, text width=2.2cm, align=center] {Alimentation : type méditerranéen, sel $<$5g/j} }
        child { node[concept, text width=2.2cm, align=center] {Activité physique : 150 min/sem (modérée) ou 75 min (intense)} }
        child { node[concept, text width=2.2cm, align=center] {Dépistage : TA, glycémie, lipidogramme, IMC, tour de taille} }
        child { node[concept, text width=2.2cm, align=center] {Éducation : message de sensibilisation argumenté} }
    };
\end{tikzpicture}
\legende{Carte mentale récapitulative de la leçon NCT03 : les quatre blocs interconnectés (système, accidents, facteurs de risque, prévention).}
\end{popfigure}

\courssub{Bilan 1 : Le système cardiovasculaire et la tension artérielle}

\begin{aretenir}[Rappel essentiel : le circuit du sang et la pression artérielle]
Le cœur est une pompe double à quatre cavités (2 oreillettes, 2 ventricules) assurant une \textbf{double circulation} :
\begin{itemize}
    \item \textbf{Petite circulation (pulmonaire)} : Ventricule droit $\rightarrow$ Artère pulmonaire $\rightarrow$ Poumons (hématose) $\rightarrow$ Veines pulmonaires $\rightarrow$ Oreillette gauche. Le sang se charge en \ce{O2}.
    \item \textbf{Grande circulation (systémique)} : Ventricule gauche $\rightarrow$ Aorte $\rightarrow$ Artères $\rightarrow$ Capillaires (échanges) $\rightarrow$ Veines caves $\rightarrow$ Oreillette droite. Le sang distribue \ce{O2} et nutriments.
\end{itemize}
La \textbf{tension artérielle (TA)} est la pression exercée par le sang sur la paroi des artères. Elle s'exprime en millimètres de mercure (mmHg) avec deux valeurs :
\begin{itemize}
    \item \textbf{Pression systolique (PAS)} : pression maximale lors de la systole ventriculaire (éjection). C'est \textbf{le premier chiffre} noté.
    \item \textbf{Pression diastolique (PAD)} : pression minimale lors de la diastole (remplissage). C'est \textbf{le deuxième chiffre}.
\end{itemize}
\textbf{Valeurs de référence chez l'adulte au repos} : $\text{PAS} < 140$ mmHg et $\text{PAD} < 90$ mmHg (seuil d'HTA). Une TA optimale est $< 120/80$ mmHg.
\end{aretenir}

\begin{attention}[Piège classique : inversion systolique/diastolique]
\emph{Erreur fréquente au Probatoire :} écrire « TA = 80/140 mmHg » ou dire « la diastolique est le premier chiffre ».
\textbf{Règle d'or} : Le \textbf{premier chiffre noté est toujours la systolique} (la plus élevée). On note toujours $\text{Systolique} / \text{Diastolique}$.
\emph{Astuce mnémotechnique} : \textbf{S}ystolique = \textbf{S}upérieur = \textbf{S}econd (dans l'ordre physiologique) mais \textbf{Premier} dans l'écriture. Exemple correct : $130/85$ mmHg.
\end{attention}

\courssub{Bilan 2 : Les trois accidents majeurs — Signes et conduite à tenir}

\begin{definition}[Infarctus du myocarde (IDM)]
Nécrose (mort) d'une partie du muscle cardiaque (myocarde) due à l'obstruction brutale d'une artère coronaire (thrombus sur plaque d'athérome). Le territoire vascularisé en aval n'est plus perfusé $\rightarrow$ ischémie $\rightarrow$ nécrose si non revascularisé rapidement (\emph{time is muscle}).
\end{definition}

\begin{definition}[Accident Vasculaire Cérébral (AVC)]
Atteinte brutale de la fonction neurologique d'origine vasculaire.
\begin{itemize}
    \item \textbf{AVC ischémique (80\%)} : obstruction artère cérébrale (thrombose/embolie) $\rightarrow$ ischémie.
    \item \textbf{AVC hémorragique (20\%)} : rupture anévrisme ou HTA mal contrôlée $\rightarrow$ hématome intracérébral.
\end{itemize}
\end{definition}

\begin{definition}[Hypertension Artérielle (HTA)]
Élévation permanente et anormale de la pression artérielle au repos : $\text{PAS} \ge 140$ mmHg \textbf{et/ou} $\text{PAD} \ge 90$ mmHg, confirmée sur 3 mesures distinctes sur 3 à 6 mois. Maladie chronique asymptomatique au début (\emph{tueur silencieux}), facteur de risque majeur d'IDM, AVC, insuffisance cardiaque, rénale, rétinopathie.
\end{definition}

\begin{propriete}[Signes d'alerte et Conduite à tenir (CAT) — Exigible au Probatoire]
\begin{center}
\begin{tabularx}{\linewidth}{|l|X|X|}\hline
\textbf{Accident} & \textbf{Signes d'alerte principaux} & \textbf{Conduite à tenir immédiate} \\\hline
\textbf{IDM} & Douleur thoracique rétrosternale, constrictive, irradiant bras gauche/mâchoire/dos, $>20$ min, non cédée au repos/nitroglycérine. Sueurs froides, angoisse, dyspnée. & \textbf{1. Appeler le 112 (urgence universelle) ou 113 (sapeurs-pompiers) immédiatement.} \\
& & 2. Mettre au repos absolu (demi-assis). \\
& & 3. Ne rien faire manger/boire. \\
& & 4. Surveiller FC/TA/conscience. \\\hline
\textbf{AVC} & \textbf{Méthode FAST (VITE)} : \textbf{F}ace (bouche déviée), \textbf{A}rm (faiblesse bras), \textbf{S}peech (troubles parole), \textbf{T}ime (temps = urgence). Céphalée brutale intense, troubles vision/équilibre. & \textbf{1. Appeler le 112 / 113 immédiatement.} \\
& & 2. Allonger tête légèrement surélevée (30°). \\
& & 3. Noter l'heure de début des signes. \\
& & 4. Ne pas donner d'aspirine (risque hémorragie). \\\hline
\textbf{HTA} & Souvent \textbf{asymptomatique}. Parfois : céphalées matinales, vertiges, bourdonnements d'oreilles, saignements de nez, fatigue visuelle. & \textbf{1. Confirmer par 3 mesures au repos (JOUR 1, JOUR 2, JOUR 3).} \\
& & 2. Consulter médecin traitant / cardiologue. \\
& & 3. Bilan : ECG, créatininémie, lipidogramme, glycémie, ECBU, fond d'œil. \\
& & 4. Hygiène de vie + Traitement au long cours si indiqué. \\\hline
\end{tabularx}
\end{center}
\end{propriete}

\courssub{Bilan 3 : Facteurs de risque — Classification et effet cumulatif}

\begin{propriete}[Classification des facteurs de risque cardiovasculaire (FRCV)]
\begin{center}
\begin{tabularx}{\linewidth}{|l|X|X|}\hline
\textbf{Catégorie} & \textbf{Facteurs} & \textbf{Action possible} \\\hline
\textbf{Non modifiables} & Âge ($>$ 50 ans H, $>$ 60 ans F), Sexe (H $>$ F avant ménopause), Hérédité (antécédent IDM/AVC prématuré parent 1er degré $<$ 55 ans H, $<$ 65 ans F), Ethnie. & \textbf{Surveillance renforcée} (dépistage précoce, targets TA/Lipides plus stricts). \\\hline
\textbf{Modifiables (Comportementaux)} & Tabagisme (actif/passif), Alimentation déséquilibrée (sel, graisses saturées, sucre, bas fibres), Sédentarité ($<$ 150 min/sem), Alcool (excès), Stress chronique, Sommeil insuffisant. & \textbf{Modification du mode de vie} (sevrage tabagique, régime méditerranéen, activité physique, gestion stress). \\\hline
\textbf{Modifiables (Métaboliques / Cliniques)} & HTA, Diabète de type 2, Dyslipidémie (LDL-c élevé, HDL-c bas, Triglycérides élevés), Surpoids/Obésité (IMC $\ge 25$, Tour de taille $>$ 94 cm H / $>$ 80 cm F), Hyperuricémie. & \textbf{Prise en charge thérapeutique} (médicamenteuse si besoin) \textbf{+} Hygiène de vie. \\\hline
\end{tabularx}
\end{center}
\end{propriete}

\begin{aretenir}[Effet cumulatif et Score de risque global]
Les FRCV \textbf{ne s'additionnent pas simplement, ils se potentialisent} (effet multiplicatif). Un homme de 55 ans, fumeur, HTA, diabétique a un risque \emph{beaucoup plus grand} que la somme des risques individuels.
$\rightarrow$ En pratique clinique, on utilise des \textbf{scores de risque global} (ex: \textbf{SCORE2} en Europe, \textbf{Framingham} aux USA) qui intègrent âge, sexe, TA, tabagisme, cholestérol total/HDL pour estimer le risque à 10 ans d'accident CV fatal/non fatal. Cela guide l'intensité de la prévention (cible LDL-c, décision traitement HTA).
\end{aretenir}

\courssub{Bilan 4 : Les trois piliers de la prévention}

\begin{propriete}[Les 3 piliers de la prévention cardiovasculaire (OMS / Société Camerounaise de Cardiologie)]
\begin{enumerate}
    \item \textbf{Alimentation de type méditerranéen / DASH} :
    \begin{itemize}
        \item Fruits/légumes : $\ge 5$ portions/jour (fibres, potassium, antioxydants).
        \item Céréales complètes, légumineuses, poisson (2-3/sem, oméga-3), huile d'olive/colza.
        \item \textbf{Limiter} : Sel $<$ 5 g/j (1 c. à café), graisses saturées (viande rouge, charcuterie, fritures, pâtisseries), sucres ajoutés, alcool (max 2 verres/j, pas tous les jours).
        \item \emph{Contexte Cameroun} : Attention au \emph{ndolé} trop huileux/salé, \emph{koki} (haricots + huile rouge), \emph{poisson braisé} (sauce pimentée salée), boissons sucrées industrielles. Privilégier \emph{légumes vapeur}, \emph{poisson grillé sans excès de sel}, fruits de saison (papaye, ananas, banane plantain bouillie).
    \end{itemize}
    \item \textbf{Activité Physique (AP) régulière} :
    \begin{itemize}
        \item Adulte : $\ge 150$ min/sem d'intensité modérée (marche rapide, vélo, danse, travaux champs) \textbf{OU} $\ge 75$ min/sem intense (football, course, HIIT). + Renforcement musculaire 2/sem.
        \item Enfant/Adolescent : 60 min/jour modérée à vigoureuse.
        \item Lutte contre la sédentarité : lever toutes les heures, marche active transport.
    \end{itemize}
    \item \textbf{Dépistage et suivi médical} :
    \begin{itemize}
        \item TA : $\ge 18$ ans, tous les ans (ou 6 mois si HTA).
        \item Glycémie à jeun / HbA1c : $\ge 45$ ans (ou avant si surpoids/FRCV), tous les 3 ans.
        \item Lipidogramme (CT, LDL, HDL, TG) : $\ge 35$ ans H / $\ge 45$ ans F (ou avant si FRCV).
        \item IMC + Tour de taille : annuel.
        \item ECG repos : si symptômes ou HTA/diabète.
    \end{itemize}
\end{enumerate}
\end{propriete}

\courssub{Entraînement 1 : Interprétation de mesures tensionnelles et de fréquence cardiaque}

La compétence « Interpréter une mesure simple de tension artérielle ou de fréquence cardiaque » est centrale. Voici la méthode rigoureuse à appliquer systématiquement.

\begin{methode}[Réflexe d'examen : Interpréter une mesure de TA et FC]
\begin{enumerate}
    \item \textbf{Lire la situation} : Identifier le contexte (âge, sexe, repos/effort, traitement en cours, symptômes).
    \item \textbf{Souligner les données chiffrées} : Noter explicitement $\text{PAS}$, $\text{PAD}$, $\text{FC}$ (bpm). Vérifier l'unité (mmHg pour TA, bpm pour FC).
    \item \textbf{Comparer aux valeurs de référence} (adulte au repos) :
    \begin{itemize}
        \item TA optimale : $< 120/80$ ; Normale : $120-129 / 80-84$ ; Normale haute : $130-139 / 85-89$.
        \item HTA Grade 1 : $140-159 / 90-99$ ; Grade 2 : $160-179 / 100-109$ ; Grade 3 : $\ge 180 / \ge 110$.
        \item FC normale au repos : $60 - 100$ bpm. Bradycardie $< 60$ ; Tachycardie $> 100$.
    \end{itemize}
    \item \textbf{Conclure} : « TA normale / HTA Grade X / Hypotension » et « FC normale / Bradycardie / Tachycardie ». Préciser si mesure unique (nécessite confirmation) ou moyenne validée.
    \item \textbf{Proposer des mesures} : Hygiène de vie (si normale haute ou Grade 1 sans FRCV), Consultation médicale rapide (Grade 2/3, symptômes, FRCV multiples), Urgence 112/113 (Urgence hypertensive : $>180/110$ avec signes atteinte d'organe : céphalée sévère, troubles visuels, douleur thoracique, dyspnée, confusion).
\end{enumerate}
\end{methode}

\begin{exoresolu}[Interprétation de cas cliniques variés]
\textbf{Cas 1 :} Monsieur M., 52 ans, comptable à Yaoundé, sédentaire, fumeur (10 cig/j). Mesure au cabinet médical au repos : $\text{TA} = 158/96$ mmHg, $\text{FC} = 88$ bpm. IMC $= 29,2$ kg/m$^2$. Tour de taille $= 102$ cm.
\textbf{Solution :}
\begin{itemize}
    \item \textbf{TA} : $158/96$ mmHg $\rightarrow$ \textbf{HTA Grade 1} (PAS $\in [140;159]$ et PAD $\in [90;99]$). \emph{Rappel} : la classification retient le grade le plus élevé ; ici les deux valeurs sont en Grade 1.
    \item \textbf{FC} : $88$ bpm $\rightarrow$ \textbf{Normale} ($60-100$).
    \item \textbf{FRCV identifiés} : Âge ($>50$ H), Tabac, Surpoids (IMC $>25$), Obésité abdominale (Tour taille $>94$ H), Sédentarité. \textbf{Risque global TRÈS ÉLEVÉ}.
    \item \textbf{Conduite} : Confirmer HTA par automesure (règle des 3 : 3 matins/3 soirs sur 3 jours). Bilan complet (créatinine, K$^+$, Na$^+$, glycémie, lipidogramme, ECG, fond d'œil). \textbf{Indication formelle de traitement médicamenteux immédiat + hygiène de vie stricte} (score risque élevé). Arrêt tabac prioritaire.
\end{itemize}

\textbf{Cas 2 :} Étudiante K., 21 ans, stress examen. Mesure infirmerie : $\text{TA} = 135/85$ mmHg, $\text{FC} = 108$ bpm.
\textbf{Solution :}
\begin{itemize}
    \item \textbf{TA} : $135/85$ $\rightarrow$ \textbf{Normale haute} (seuil $130-139 / 85-89$). Pas d'HTA sur une mesure unique en contexte stress.
    \item \textbf{FC} : $108$ bpm $\rightarrow$ \textbf{Tachycardie sinusale} (réaction au stress/aigu).
    \item \textbf{Conduite} : Repos 15 min, réévaluer. Si normalise $\rightarrow$ simple réaction émotionnelle. Conseil gestion stress (respiration, sport). Pas de traitement.
\end{itemize}

\textbf{Cas 3 :} Papa N., 68 ans, traité HTA (Amlodipine 5mg + Hydrochlorothiazide 12,5mg). Automesure matin : $\text{TA} = 102/58$ mmHg, $\text{FC} = 54$ bpm. Il signale fatigue, vertiges à la station debout.
\textbf{Solution :}
\begin{itemize}
    \item \textbf{TA} : $102/58$ $\rightarrow$ \textbf{Hypotension artérielle} (PAS $< 100$ ou PAD $< 60$ souvent symptomatique). \textbf{Hypotension orthostatique probable} (vertiges debout).
    \item \textbf{FC} : $54$ bpm $\rightarrow$ \textbf{Bradycardie} ($< 60$). Attention : l'Hydrochlorothiazide ne donne pas de bradycardie, l'Amlodipine (dihydropyridine) non plus. Vérifier s'il prend un \textbf{bêta-bloquant} ou un \textbf{calcium-antagoniste non dihydropyridine} (ex: vérapamil) non mentionné, ou pathologie propre (bloc AV). \textbf{Urgence relative} : Contacter médecin traitant \textbf{aujourd'hui} pour réajustement traitement (arrêt diurétique si déshydraté, baisse posologie). Ne pas arrêter brutalement seul.
\end{itemize}
\end{exoresolu}

\begin{exemplebox}[Exemple chiffré : Tachycardie au repos]
Un adulte de 35 ans, au calme depuis 10 minutes, présente une $\text{FC} = 110$ bpm.
\textbf{Interprétation} : Valeur \textbf{anormalement élevée} (Tachycardie, $> 100$ bpm au repos).
\textbf{Causes possibles} : Fièvre, anémie, hyperthyroïdie, déshydratation, anxiété/stress aigu, consommation café/thé/énergisants, effet médicamenteux (bronchodilatateurs, thyroxine surdosée), arythmie (flutter/fibrillation auriculaire).
\textbf{Conduite} : \textbf{Avis médical conseillé dans la journée} (recherche cause étiologique). Si signes de gravité (douleur thoracique, dyspnée au repos, lipothymie, FC $> 150$ irrégulière) $\rightarrow$ Appeler 112/113.
\end{exemplebox}

\courssub{Entraînement 2 : Analyse de situations de vie et proposition de prévention}

Au Probatoire, on vous présente souvent un « profil patient » ou une « situation familiale ». Il faut lister les FRCV, classer (modifiable/non), évaluer le risque global et proposer un plan d'action personnalisé, réaliste et priorisé.

\begin{methode}[Analyse d'une situation de vie : Profil de risque \& Plan prévention]
\begin{enumerate}
    \item \textbf{Lister exhaustivement} tous les FRCV présents dans le texte (données chiffrées : TA, IMC, Tour taille, Glycémie, Lipides, Tabac, Alcool, Activité physique, Antécédents familiaux, Âge, Sexe).
    \item \textbf{Classer} chaque facteur : \textbf{NM} (Non Modifiable) ou \textbf{M} (Modifiable : Comportemental ou Métabolique).
    \item \textbf{Identifier les cibles prioritaires} : Facteurs modifiables les plus impactants (Tabac = \#1, HTA non contrôlée, Diabète déséquilibré, Dyslipidémie majeure, Sédentarité, Obésité abdominale).
    \item \textbf{Rédiger le plan d'action} (SMART : Spécifique, Mesurable, Atteignable, Réaliste, Temporel) :
    \begin{itemize}
        \item \textbf{Médical} : Consultation spécialisée, bilan, adaptation traitement, vaccin grippe/COVID.
        \item \textbf{Alimentation} : Objectif concret (ex: « Remplacer huile rouge par huile arachide/colza pour cuisson ; 1 fruit à chaque repas ; Sel $<$ 1 c. à café/j »).
        \item \textbf{Activité physique} : Objectif progressif (ex: « Marche 30 min, 5 soirs/sem après le travail ; inscription club sport quartier »).
        \item \textbf{Comportement} : Sevrage tabagique (consultation tabacologue, substituts nicotiniques), modération alcool, sommeil 7-8h, gestion stress (cohérence cardiaque).
        \item \textbf{Suivi} : Prochaine TA (J+15), Prochain lipidogramme (3 mois), Prochaine consultation (1 mois).
    \end{itemize}
\end{enumerate}
\end{methode}

\begin{exoresolu}[Analyse de situation : Le cas de la famille B. à Douala]
\textbf{Énoncé :} Monsieur B., 54 ans, chauffeur de taxi (assise 12h/j), fumeur 20 cig/j depuis 30 ans. TA mesurée campagne dépistage : $162/102$ mmHg. IMC $= 31,5$. Tour taille $= 108$ cm. Glycémie à jeun $= 1,35$ g/L (2 mesures). Père décédé IDM à 58 ans. Maman vivante, HTA traitée.
Madame B., 48 ans, vendeuse marché (station debout), non fumeuse. TA $= 138/88$ mmHg. IMC $= 28$. Tour taille $= 92$ cm. Glycémie $= 1,05$ g/L. Mère diabétique.
Leur fils aîné, 22 ans, étudiant, fumeur 5 cig/j, IMC $= 22$, TA $= 118/72$, sédentaire (jeux vidéo).

\textbf{Solution détaillée :}

\begin{center}
\begin{tabularx}{\linewidth}{|l|X|X|X|}\hline
\textbf{Individu} & \textbf{FRCV Non Modifiables} & \textbf{FRCV Modifiables} & \textbf{Niveau de risque / Priorités} \\\hline
\textbf{M. B. (54 ans)} & Âge ($>50$), Sexe (H), Hérédité (Père IDM $<55$ ans) & \textbf{Majeurs} : Tabagisme lourd (30 PA), HTA Grade 2 ($162/102$), Obésité (IMC 31,5), Obésité abdominale (108 cm), Diabète type 2 (Glycémie $1,35$ g/L $\ge 1,26$), Sédentarité extrême (métier), Alcool probable (chauffeur). & \textbf{RISQUE TRÈS ÉLEVÉ (Score SCORE2 $\gg 10\%$)}. \\
& & & \textbf{Urgence} : Prise en charge médicale immédiate (Cardiologue/ Médecin traitant). Traitement HTA + Diabète + Statine (LDL-c cible $< 1,4$ ou $< 1,8$ g/L selon score). \textbf{Arrêt tabac IMPÉRATIF} (sevrage aidé). \\
& & & \textbf{Hygiène} : Régime hyposodé strict, fractionné, index glycémique bas. Activité : Début très progressif (marche 10 min $\times$ 3/j) $\rightarrow$ 30 min. Surveillance pieds (risque artérite). \\\hline
\textbf{Mme B. (48 ans)} & Sexe (F, pré-ménopause), Hérédité (Mère diabète) & HTA limite / Normale haute ($138/88$), Surpoids (IMC 28), Obésité abdominale (92 cm $> 80$), Prédiabète (Glycémie $1,05-1,25$), Sédentarité (station debout $\neq$ activité physique dynamique). & \textbf{RISQUE MODÉRÉ $\rightarrow$ ÉLEVÉ si non action}. \\
& & & \textbf{Priorité} : Perte de poids 5-10\% (cible IMC $< 25$, Tour taille $< 80$). Activité physique dynamique (marche rapide 30 min/j). Régime méditerranéen/DASH. Contrôle TA / Glycémie / Lipides dans 6 mois. \\\hline
\textbf{Fils (22 ans)} & Sexe (H), Hérédité (Père IDM prématuré + HTA + Diabète ; Mère Prédiabète/HTA) & Tabac (5 cig/j), Sédentarité extrême. & \textbf{RISQUE FAIBLE ACTUEL MAJORÉ PAR HÉRÉDITÉ FORTE}. \\
& & & \textbf{Priorité PRÉVENTION PRIMAIRE} : \textbf{Arrêt tabac maintenant} (avant dépendance forte). Intégration sport loisir (foot, salle, course) 3$\times$/sem. Éducation nutritionnelle. Dépistage lipidique/glycémie à 25 ans (antécédents). \\\hline
\end{tabularx}
\end{center}

\textbf{Message de sensibilisation familial (exercice transversal) :}
\begin{quote}
\textbf{« Famille B., votre cœur vous parle. Écoutez-le ! »} \\
Papa, ton cœur fatigue sous le poids du tabac, du sel, du sucre et de l'assise. 162/102 et 1,35 g/L, c'est l'alerte rouge. Le médecin t'attend \textbf{cette semaine} pour protéger ton cœur et ton cerveau. Arrêter le taxi 30 min pour marcher, c'est gagner des années de vie. \\
Maman, ton tour de taille (92 cm) et ta glycémie (1,05) sont des signaux d'alarme silencieux. Tu as le pouvoir d'inverser la tendance : marché à pied, légumes vapeur, moins d'huile rouge, 30 min de danse active par soir. Tu évites le diabète et l'HTA. \\
Fils, 5 cigarettes à 22 ans avec l'héritage cardiaque de tes parents, c'est jouer à la roulette russe. Arrête aujourd'hui. Le sport, c'est ton meilleur ami pour la vie. \\
\textbf{Ensemble :} Cuisinez sain (ndolé léger, poisson grillé, fruits), marchez ensemble le week-end, soutenez Papa dans son sevrage. \textbf{Le cœur de la famille bat plus fort quand on en prend soin.}
\end{quote}
\end{exoresolu}

\courssub{Complément Probatoire : Rédiger un message de sensibilisation argumenté}

L'épreuve d'Éducation à la Santé (souvent couplée à SVT ou Éducation à la Citoyenneté) évalue fréquemment la capacité à \textbf{produire un message de prévention court, clair, scientifiquement juste et persuasif}, adapté à une cible (jeunes, adultes, population spécifique). C'est une compétence transversale : mobiliser des savoirs (biologie, physiopathologie), des savoir-faire (communication, argumentation) et des savoir-être (empathie, responsabilité).

\begin{methode}[Structure type d'un message de sensibilisation efficace (Modèle « AIDA » adapté santé)]
\begin{enumerate}
    \item \textbf{Accroche (Attention)} : Titre court, percutant, image forte ou question. Ex: « Votre cœur a-t-il 20 ans de plus que vous ? », « 1 cigarette = 11 minutes de vie en moins ».
    \item \textbf{Intérêt (Donnée choc / Réalité locale)} : Chiffre clé, statistique Cameroun/Monde, situation vécue. Ex: « Au Cameroun, 1 adulte sur 3 est hypertendu, la moitié l'ignore. », « L'AVC frappe aussi à 30 ans. »
    \item \textbf{Désir / Conviction (Argument scientifique simple)} : Expliquer \textbf{POURQUOI} (mécanisme simplifié) et \textbf{QUEL BÉNÉFICE} concret. Ex: « Le sel durcit les artères $\rightarrow$ le cœur force $\rightarrow$ il s'épuise (insuffisance cardiaque). Baisser le sel de 3g/j = -5 mmHg TA = -20\% AVC. »
    \item \textbf{Action (Appel à l'acte concret, faisable, immédiat)} : Un geste précis, mesurable. Ex: « Aujourd'hui : je pose la salière, je marche 20 min, je prends RDV pour ma TA. » \textbf{Éviter} : « Faites du sport », « Mangez mieux » (trop vague).
    \item \textbf{Signature / Ressource} : Structure émettrice (Ministère Santé, ONG, École), Numéro vert (112, 113, ligne verte tabac), Lieu/Date prochaine campagne.
\end{enumerate}
\end{methode}

\begin{attention}[Pièges à éviter dans la rédaction (Probatoire)]
\begin{itemize}
    \item \textbf{Jargon médical incompris} : « Athérome », « LDL-c », « Systolique » sans définition simple. $\rightarrow$ Utiliser « graisse dans les artères », « mauvais cholestérol », « premier chiffre de la tension ».
    \item \textbf{Ton moralisateur / culpabilisant} : « Vous êtes irresponsables de fumer. » $\rightarrow$ Préférer : « Le tabac vous piège, mais vous pouvez vous en libérer. Aide dispo au 112. »
    \item \textbf{Consignes irréalistes} : « Arrêtez tout sel », « Courrez 1h/jour ». $\rightarrow$ « Réduisez le sel progressivement », « Marchez 10 min 3 fois/jour ».
    \item \textbf{Oublier le « Pourquoi » scientifique} : Un message sans argument biologique perd en crédibilité.
    \item \textbf{Mauvaise orthographe / syntaxe} : Pénalisé en communication.
\end{itemize}
\end{attention}

\begin{exoresolu}[Rédaction d'un message de sensibilisation — Sujet Probatoire type]
\textbf{Consigne :} « Vous êtes membre du club santé de votre lycée. Rédigez un message d'affichage (max 150 mots) pour inciter vos camarades à pratiquer une activité physique régulière. »

\textbf{Proposition de rédaction :}

\begin{center}
\fcolorbox{popblue}{popcream}{
\begin{minipage}{0.92\linewidth}
\textbf{\large \textcolor{popblue}{BOUGE TON CŒUR, GAGNE TA VIE !}} \\[4pt]
\textcolor{popdark}{Saviez-vous que \textbf{1 jeune sur 4} au Cameroun ne bouge pas assez (OMS) ? Rester assis 8h/j (cours, téléphone, TV) rigidifie vos artères \textbf{dès 15 ans}. Le cœur, c'est un muscle : s'il ne travaille pas, il s'atrophie. Résultat : essoufflement précoce, tension qui monte, risque d'infarctus à 40 ans.} \\[4pt]
\textbf{\textcolor{popgreen}{LE DÉFI 30 MIN / JOUR :}} \\
\checkmark Marche rapide vers l'école / arrêts de bus plus loin. \\
\checkmark Foot, basket, danse, corde à sauter \textbf{entre 2 révisions} (meilleure mémoire !). \\
\checkmark Zéro ascenseur, escaliers obligatoires. \\[4pt]
\textcolor{poporange}{\textbf{Cette semaine :}} Je choisis \textbf{UN} sport, je fixe \textbf{3 créneaux} dans mon agenda, j'invite \textbf{UN} ami. \\
\textcolor{poppurple}{Club Santé Lycée — Salle 12, Mardi 16h — Infirmerie pour test FC/TA gratuit. \#MonCœurBatFort}
\end{minipage}
}
\end{center}

\textbf{Analyse de la réponse :} Accroche positive (gain), Chiffre local (OMS Cameroun), Mécanisme simple (cœur muscle, artères rigides), Bénéfices immédiats (mémoire), Actions ultra-concrètes (arrêt bus, escaliers, créneaux agenda), Appel à l'acte daté (cette semaine), Ressource locale (Club, Infirmerie). \textbf{Note attendue : 5/5.}
\end{exoresolu}

\courssub{Évaluation formative : Testez vos acquis}

\begin{exercice}[Auto-évaluation — Niveau Probatoire]
\textbf{Exercice 1 (Interprétation) :} Un homme de 45 ans, au repos depuis 10 min, présente $\text{TA} = 142/92$ mmHg et $\text{FC} = 58$ bpm. Il prend un bêta-bloquant (Bisoprolol) pour HTA.
\begin{enumerate}
    \item Classer la TA et la FC.
    \item Expliquer la FC basse.
    \item Proposer la conduite à tenir immédiate.
\end{enumerate}

\textbf{Exercice 2 (Analyse de risque) :} Femme, 58 ans, ménopausée, employée bureau. TA $= 135/85$ mmHg (confirmée). Non fumeuse. IMC $= 26$. Tour taille $= 88$ cm. Glycémie $= 1,10$ g/L. Cholestérol total $= 2,40$ g/L (HDL $= 0,55$, LDL $= 1,60$). Mère : AVC 70 ans.
\begin{enumerate}
    \item Lister et classer ses FRCV.
    \item Estimer son niveau de risque global (Faible / Modéré / Élevé / Très élevé) et justifier.
    \item Définir 3 objectifs prioritaires de prévention (SMART).
\end{enumerate}

\textbf{Exercice 3 (Rédaction - Compétence transversale) :} Rédigez un message court (100-120 mots) pour une campagne radio locale (langue française + quelques mots pidgin/camfranglais acceptés si naturel) ciblant les conducteurs de moto-taxi (\emph{benskin}) sur le thème : « Hypertension et conduite : danger pour tous ». Structure : Accroche $\rightarrow$ Fait scientifique $\rightarrow$ Geste préventif $\rightarrow$ Ressource.
\end{exercice}

\begin{corrige}[Exercice 1]
\begin{enumerate}
    \item \textbf{TA} : $142/92$ mmHg $\rightarrow$ \textbf{HTA Grade 1} (seuils $140-159 / 90-99$). \textbf{FC} : $58$ bpm $\rightarrow$ \textbf{Bradycardie} ($< 60$ bpm).
    \item \textbf{Explication} : Le bêta-bloquant (Bisoprolol) bloque les récepteurs $\beta_1$ du nœud sinusal $\rightarrow$ diminue l'automatisme $\rightarrow$ baisse la FC (effet thérapeutique recherché pour protéger le myocarde). Une FC $58$ bpm \textbf{sous traitement} est souvent \textbf{bien tolérée si asymptomatique} (pas de vertiges, fatigue anormale, syncope). C'est un effet pharmacologique attendu, pas une pathologie primaire.
    \item \textbf{Conduite} : \textbf{Ne pas arrêter le traitement !} Vérifier tolérance clinique (symptômes). Si asymptomatique $\rightarrow$ Maintenir traitement, surveiller FC/TA à domicile, consulter médecin traitant à la prochaine visite programmée (ou plus tôt si symptômes apparaissent). Si symptomatique (asthénie majeure, vertiges, syncope) $\rightarrow$ Contacter médecin \textbf{dans les 24-48h} pour éventuelle réduction posologie.
\end{enumerate}
\end{corrige}

\begin{corrige}[Exercice 2]
\begin{enumerate}
    \item \textbf{FRCV :}
    \begin{itemize}
        \item NM : Âge ($>60$ F post-ménopause), Sexe (F), Hérédité (Mère AVC 70 ans $>65$ ans $\rightarrow$ pas « prématuré » strict mais antécédent familial).
        \item M Comportementaux : Sédentarité (bureau).
        \item M Métaboliques : \textbf{Normale haute} (TA $135/85$ confirmée $\rightarrow$ \textbf{PAS d'HTA} selon seuil $\ge 140/90$), Surpoids (IMC 26), Obésité abdominale (Tour taille 88 cm $> 80$ F), Prédiabète (Glycémie $1,10$ g/L $\in [1,00 ; 1,25]$), Dyslipidémie (LDL-c $1,60$ g/L $\ge 1,6$ cible risque modéré ? $\rightarrow$ Cible LDL dépend du risque global).
    \end{itemize}
    \item \textbf{Niveau de risque :} \textbf{MODÉRÉ}. Pas de diabète avéré, pas d'HTA avérée, pas de tabagisme, pas d'atteinte d'organe. Score SCORE2 (F, 58 ans, non fumeuse, PAS 135, CT 2,40, HDL 0,55) $\approx 3-5\%$ à 10 ans (risque modéré). \textbf{Mais} cumul facteurs (obésité abdominale, prédiabète, dyslipidémie, normale haute) $\rightarrow$ \textbf{Risque majoré, prévention active nécessaire}.
    \item \textbf{Objectifs SMART prioritaires :}
    \begin{enumerate}
        \item \textbf{Perte de poids :} Perdre 5 kg en 6 mois (cible IMC $< 24$, Tour taille $< 80$ cm) par régime méditerranéen hypocalorique modéré (-500 kcal/j) + Activité physique.
        \item \textbf{Activité physique :} Marche rapide 30 min, 5 soirs/sem (sortie bureau) + Renforcement 2$\times$/sem (gainage, squats). Objectif : 150 min/sem modérée.
        \item \textbf{Surveillance métabolique :} Contrôle TA automesure 1$\times$/mois (cible $< 130/80$), Glycémie/HbA1c + Lipidogramme dans 6 mois. Réévaluer risque et cible LDL-c.
    \end{enumerate}
\end{enumerate}
\end{corrige}

\begin{corrige}[Exercice 3 — Proposition de message radio]
\begin{quote}
\textbf{« Benskineur, ton cœur aussi porte casque ! »} \\[4pt]
Frère conducteur, rouler 12h/jour, stress, bruit, poussière, café-kola à gogo, ça use le moteur... et ton cœur ! \textbf{1 conducteur sur 3} a la tension haute sans le savoir. La tension qui monte, c'est comme une moto surrégimée : le moteur (cœur) pète, le frein (cerveau) lâche $\rightarrow$ AVC, malaise au volant = accident grave pour toi et ton client. \\[4pt]
\textbf{Le geste qui sauve :} \textbf{Tous les 1ers mardis du mois, arrête-toi au Centre de Santé du quartier (ou grand carrefour) : 5 min, gratuit, on te prend la tension.} Si $> 14/9$, on te guide. Prends ta vie en main, comme ton guidon. \\[4pt]
\textbf{Santé Publique Cameroun — Ligne verte 1515. Ton cœur, ta vie, ton volant.}
\end{quote}
\end{corrige}

\courssub{Synthèse finale et mots-clés pour le Probatoire}

\begin{aretenir}[Mots-clés et concepts « must-know » pour l'examen]
\begin{itemize}
    \item \textbf{Double circulation} : Petite (cœur $\leftrightarrow$ poumon) / Grande (cœur $\leftrightarrow$ corps). Cœur = pompe double.
    \item \textbf{TA} = $\text{PAS} / \text{PAD}$ (mmHg). \textbf{PAS = 1er chiffre}. Normale $< 140/90$. HTA $\ge 140/90$ confirmée.
    \item \textbf{IDM} : Douleur thoracique $>20$ min $\rightarrow$ \textbf{112 / 113}. \textbf{Time is muscle}.
    \item \textbf{AVC} : \textbf{FAST} (Visage, Bras, Parole, Temps) $\rightarrow$ \textbf{112 / 113}. Pas d'aspirine avant scanner.
    \item \textbf{FRCV} : Non modifiables (Âge, Sexe, Hérédité) vs Modifiables (Tabac, Nutrition, Sédentarité, Alcool, Stress, HTA, Diabète, Lipides, Poids). \textbf{Effet cumulatif}.
    \item \textbf{Prévention} : 3 piliers \textbf{Alimentation (Méditerranéen/DASH, Sel $<$5g)} — \textbf{Activité Physique (150 min/sem)} — \textbf{Dépistage (TA, Glycémie, Lipides, IMC, Tour taille)}.
    \item \textbf{Compétence transversale} : Message sensibilisation = Accroche + Fait scientifique local + Mécanisme simple + Action concrète (SMART) + Ressource.
\end{itemize}
\end{aretenir}

\begin{savaistu}[Chiffres clés pour vos dissertations et oraux]
\begin{itemize}
    \item \textbf{17,9 millions} de décès/an dans le monde par maladies cardiovasculaires (OMS 2019) = \textbf{32\% de la mortalité mondiale}. \textbf{Première cause de mortalité}.
    \item \textbf{80\%} des infarctus et AVC \textbf{prématurés} (avant 70 ans) sont \textbf{évitables} en agissant sur les facteurs de risque modifiables (Tabac, Alimentation, Activité physique, Alcool).
    \item Au \textbf{Cameroun} : Prévalence HTA $\approx$ \textbf{30\%} des adultes (1 sur 3). Seule \textbf{la moitié} le sait. Parmi les traités, \textbf{moins de 30\% sont contrôlés} ($< 140/90$). L'AVC est la 1ère cause de handicap moteur adulte.
    \item \textbf{Coût économique} : Les MNT (Maladies Non Transmissibles dont CV) coûtent aux pays à revenu faible/intermédiaire $\approx$ 7000 milliards \$ / an (pertes productivité, frais santé). La prévention rapporte \textbf{7 \$ pour 1 \$ investi} (OMS).
\end{itemize}
\end{savaistu}

\begin{experience}[TP / Activité pratique suggérée : « Journée Cœur en Santé » au lycée]
\textbf{Objectif} : Mettre en œuvre les savoir-faire « Mesurer TA/FC », « Identifier FRCV », « Conseiller prévention ».
\textbf{Matériel} : Tensiomètres brassard validés (électroniques), stéthoscopes, balances impédancemètres, mètre ruban, fiches anonymisées, affiches prévention.
\textbf{Protocole} :
\begin{enumerate}
    \item Atelier 1 : \textbf{Mesure TA/FC rigoureuse} (repos 5 min, bras au niveau cœur, 3 mesures à 1 min intervalle, moyenne des 2 dernières). Noter sur fiche.
    \item Atelier 2 : \textbf{Anthropométrie} : Taille, Poids $\rightarrow$ IMC. Tour de taille (milieu dernière côte - crête iliaque).
    \item Atelier 3 : \textbf{Questionnaire FRCV} (anonyme) : Tabac, Alcool, Activité physique (IPAQ court), Alimentation (fréquence fruits/légumes/sel/huile), Antécédents familiaux.
    \item Atelier 4 : \textbf{Conseil personnalisé} (binôme élève-infirmier/enseignant) : Interprétation résultats $\rightarrow$ Remise fiche « Mon Cœur » avec code couleur (Vert/Orange/Rouge) + 1 action concrète à faire cette semaine.
    \item Atelier 5 : \textbf{Production collective} : Création affiches / spots radio / vidéos TikTok (1 min) pour campagne lycée. Vote du meilleur message.
\end{enumerate}
\textbf{Exploitation} : Dépouillement statistique anonyme (prévalence HTA inconnue, surpoids, sédentarité dans l'établissement) $\rightarrow$ Restitution en assemblée générale + Plaidoyer auprès administration (cantine, temps sport, point d'eau).
\end{experience}

\courssub{Conclusion de la leçon}

Cette leçon NCT03 vous a permis de traverser tout le continuum : de la \textbf{physiologie normale} (cœur pompe, vaisseaux, pression) à la \textbf{pathologie aiguë} (IDM, AVC) et \textbf{chronique} (HTA), en passant par la \textbf{étiologie multifactorielle} (facteurs de risque) pour aboutir à l'\textbf{action concrète} (prévention primaire, secondaire, tertiaire). Vous disposez désormais des outils scientifiques pour \textbf{comprendre}, des grilles de lecture pour \textbf{évaluer} (mesures, scores) et des leviers pour \textbf{agir} (hygiène de vie, dépistage, message de santé publique).

Rappelez-vous : \textbf{La maladie cardiovasculaire n'est pas une fatalité génétique, c'est le plus souvent le résultat d'un mode de vie qui s'installe dans la durée.} Chaque choix alimentaire, chaque minute de marche, chaque cigarette non fumée, chaque contrôle de tension est un investissement sur votre capital santé. En tant que futurs bacheliers, citoyens, et peut-être professionnels de santé, vous êtes les premiers acteurs de cette prévention. \textbf{Prenez soin de votre cœur, il vous le rendra.}

\begin{center}
\textbf{— Fin de la leçon NCT03 —}
\end{center}

\newpage

\coursec{Activité d'intégration}

\courssub{Objectifs de l'activité}
Cette activité d'intégration vise à mobiliser l'ensemble des savoirs et savoir-faire acquis au cours de la leçon \emph{Lutte contre les accidents cardiovasculaires} dans une situation authentique et contextualisée. Elle permet à l'élève de :
\begin{itemize}
    \item \textbf{Interpréter} des mesures physiologiques réelles (tension artérielle, fréquence cardiaque) en les confrontant aux valeurs de référence ;
    \item \textbf{Identifier} et \textbf{classer} les facteurs de risque cardiovasculaire (modifiables et non modifiables) à partir d'un cas clinique ;
    \item \textbf{Proposer} des mesures de prévention concrètes, adaptées au contexte socio-culturel et économique camerounais ;
    \item \textbf{Communiquer} une information de santé claire, bienveillante et incitative sous forme d'un message de sensibilisation.
\end{itemize}
Cette démarche place l'élève en situation de \cle{résolution de problème} : il doit articuler des connaissances biologiques (fonctionnement cardiovasculaire, physiopathologie de l'hypertension, de l'athérosclérose) et des compétences transversales (analyse de données, argumentation, éducation à la santé).

\courssub{Situation : Enquête santé — Le profil cardiovasculaire de M.~Fotso}
\textbf{M.~Fotso}, 50~ans, est un commerçant dynamique opérant au marché central de Douala. Il consulte lors d'une journée de dépistage gratuit organisée par la délégation régionale de la Santé publique à l'occasion de la Journée mondiale de l'hypertension artérielle (17~mai). L'infirmier du poste de santé recueille les données suivantes :

\begin{center}
\begin{tabularx}{\linewidth}{|l|X|}\hline
\textbf{Paramètre} & \textbf{Valeur observée} \\ \hline
Âge & 50~ans \\ \hline
Profession & Commerçant (vente de pièces détachées, position assise prolongée, horaires étendus 7~h--20~h) \\ \hline
Tension artérielle (mesure 1, bras droit, position assise, après 5~min de repos) & \textbf{16 / 10} (soit \SI{160}{\mmHg} / \SI{100}{\mmHg}) \\ \hline
Tension artérielle (mesure 2, 10~min plus tard, mêmes conditions) & \textbf{15 / 9} (soit \SI{150}{\mmHg} / \SI{90}{\mmHg}) \\ \hline
Fréquence cardiaque au repos (radiale, 1~min) & \textbf{92~batt/min} \\ \hline
Tabagisme & Fumeur régulier depuis 20~ans (\textasciitilde 15~cigarettes/jour) \\ \hline
Alimentation & Repas quotidiens au restaurant du marché : \emph{ndolé} très salé, \emph{koki} riche en huile de palme, \emph{poisson braisé} sauce pimentée, bouillons de viande gras. Consommation d'alcool modérée (2 à 3 verres de bière locale le week-end). \\ \hline
Activité physique & Aucune activité sportive délibérée. Déplacements en moto-taxi. Station assise > 10~h/jour. \\ \hline
Antécédents familiaux & Père décédé d'un \textbf{accident vasculaire cérébral (AVC)} à 60~ans. Mère vivante, hypertendue traitée. \\ \hline
Autres & Pas de diabète connu (glycémie à jeun non réalisée), pas de plainte actuelle (asymptomatique). \\ \hline
\end{tabularx}
\end{center}

L'infirmier note \og HTA stade 1 -- Prise en charge urgente \fg sur la fiche de dépistage et oriente M.~Fotso vers le médecin du district pour confirmation diagnostique et prise en charge.

\courssub{Consignes}
Travaillant en groupes de 4 à 5 élèves, vous disposez de 45~minutes pour traiter les quatre tâches ci-dessous. Un rapporteur par groupe présentera les résultats à la classe entière (15~min de restitution).

\begin{enumerate}
    \item \textbf{Interprétation des mesures physiologiques} : Comparez les deux valeurs de tension artérielle (TA) et la fréquence cardiaque (FC) de M.~Fotso aux valeurs de référence pour un adulte de 50~ans au repos. Caractérisez le statut tensionnel et le rythme cardiaque. Justifiez votre conclusion en citant les seuils numériques.
    \item \textbf{Identification et classification des facteurs de risque} : Dressez la liste exhaustive des facteurs de risque cardiovasculaire présents chez M.~Fotso. Classez-les en deux colonnes : \textbf{Facteurs modifiables} (liés au mode de vie) et \textbf{Facteurs non modifiables} (liés à l'individu). Pour chaque facteur modifiable, précisez le mécanisme physiopathologique par lequel il favorise l'athérosclérose ou l'hypertension.
    \item \textbf{Construction d'un tableau de prévention personnalisé} : Élaborez un tableau à deux colonnes : \og Facteur de risque identifié \fg $\rightarrow$ \og Mesure de prévention concrète et adaptée au contexte de M.~Fotso (Douala, commerçant, 50~ans) \fg. Les mesures doivent être \textbf{spécifiques}, \textbf{réalisables} et \textbf{chiffrées} si possible (ex. : \og Réduire la consommation de sel à < 5~g/jour \fg plutôt que \og Manger moins salé \fg).
    \item \textbf{Rédaction d'un message de sensibilisation} : Rédigez un court message (5 lignes maximum, ton bienveillant, sans jargon médical complexe) que M.~Fotso pourrait recevoir par SMS ou affiche lors de la campagne. Ce message doit l'inciter à consulter et à modifier un comportement précis.
\end{enumerate}

\courssub{Ressources à mobiliser}
Pour réaliser cette activité, vous devez mobiliser les connaissances suivantes (rappel synthétique) :

\begin{itemize}
    \item \textbf{Valeurs de référence de la tension artérielle (OMS / Société Camerounaise de Cardiologie)} :
    \begin{center}
    \begin{tabular}{|l|c|c|}\hline
    \textbf{Catégorie} & \textbf{Systolique (mmHg)} & \textbf{Diastolique (mmHg)} \\ \hline
    Optimale & $< 120$ & $< 80$ \\ \hline
    Normale & 120--129 & 80--84 \\ \hline
    Normale haute & 130--139 & 85--89 \\ \hline
    \textbf{HTA Stade 1 (légère)} & \textbf{140--159} & \textbf{90--99} \\ \hline
    HTA Stade 2 (modérée) & 160--179 & 100--109 \\ \hline
    HTA Stade 3 (sévère) & $\ge 180$ & $\ge 110$ \\ \hline
    \end{tabular}
    \end{center}
    \textit{Rappel} : 1~cmHg = 10~mmHg. Les valeurs 16/10 et 15/9 sont exprimées en cmHg (usage courant au Cameroun), équivalentes à \SI{160}{\mmHg}/\SI{100}{\mmHg} et \SI{150}{\mmHg}/\SI{90}{\mmHg}.
    
    \item \textbf{Fréquence cardiaque au repos (adulte)} : Normale entre \textbf{60 et 100~batt/min}. $> 100$ = tachycardie ; $< 60$ = bradycardie. Une FC entre 80 et 100~batt/min est un marqueur de risque cardiovasculaire indépendant.
    
    \item \textbf{Principaux facteurs de risque cardiovasculaire (score de Framingham / OMS)} :
    \begin{itemize}
        \item \textit{Non modifiables} : Âge ($\ge$ 50~ans homme, $\ge$ 60~ans femme), Sexe (homme $>$ femme avant ménopause), Hérédité (antécédent familial précoce : père $<$ 55~ans, mère $<$ 65~ans).
        \item \textit{Modifiables} : Hypertension artérielle, Tabagisme, Dyslipidémie (LDL-c élevé, HDL-c bas), Diabète / Prédiabète, Sédentarité, Surpoids / Obésité (IMC $\ge$ 25), Alimentation déséquilibrée (excès sel, graisses saturées, alcool), Stress psychosocial.
    \end{itemize}
    
    \item \textbf{Mesures de prévention primaire (OMS \emph{Best Buys})} : Réduction sel $<$ \SI{5}{g/j} ; 5 portions fruits/légumes/j ; Activité physique $\ge$ \SI{150}{min/sem} (intensité modérée) ; Arrêt tabac ; Limitation alcool ; Contrôle pondéral (IMC 18,5--24,9) ; Gestion du stress ; Dépistage régulier (TA, glycémie, lipidogramme).
\end{itemize}

\courssub{Démarche et corrigé commenté}
\begin{exoresolu}[Corrigé commenté de l'activité d'intégration -- Profil de M.~Fotso]
\textbf{Tâche 1 : Interprétation des mesures physiologiques}

\textit{Analyse de la Tension Artérielle (TA)} :
Les deux mesures successives (\SI{160}{\mmHg}/\SI{100}{\mmHg} puis \SI{150}{\mmHg}/\SI{90}{\mmHg}) sont toutes deux supérieures aux seuils de la normale haute (\SI{139}{\mmHg}/\SI{89}{\mmHg}).
\begin{itemize}
    \item La première mesure (\SI{160}{\mmHg}/\SI{100}{\mmHg}) place M.~Fotso en \textbf{HTA Stade 1} (systolique et diastolique dans l'intervalle 140--159 / 90--99 pour la diastolique \SI{100}{\mmHg} qui est en Stade 2, mais la classification retient le stade le plus élevé atteint par l'une des deux valeurs ; ici on note une HTA systo-diastolique stade 1/2).
    \item La deuxième mesure (\SI{150}{\mmHg}/\SI{90}{\mmHg}) confirme une \textbf{HTA Stade 1} (systolique dans l'intervalle 140--159, diastolique à la limite 90).
\end{itemize}
\textbf{Conclusion} : M.~Fotso présente une \textbf{hypertension artérielle confirmée au stade 1} (selon les recommandations, le diagnostic nécessite 2 mesures élevées lors de 2 consultations distinctes, mais en dépistage, deux mesures élevées successives orientent fortement vers le diagnostic et justifient l'adressage). La valeur systolique (\SIrange{150}{160}{\mmHg}) est plus élevée que la diastolique (\SIrange{90}{100}{\mmHg}), ce qui est fréquent après 50~ans du fait de la rigidité artérielle (\cle{hypertension systolique prédominante}).

\textit{Analyse de la Fréquence Cardiaque (FC)} :
FC = \SI{92}{batt/min}. Elle se situe dans la norme (60--100) mais dans la \textbf{zone haute de la normale}. Une FC de repos $> 80$~batt/min est associée à une augmentation du risque cardiovasculaire et de mortalité toutes causes (études épidémiologiques type Framingham, Paris Prospective Study). Elle témoigne ici d'une \cle{prédominance du tone orthosympathique} (stress, tabagisme, sédentarité, surpoids probable) et/ou d'un déconditionnement cardiovasculaire.

\vspace{0.5cm}
\textbf{Tâche 2 : Identification et classification des facteurs de risque}

\begin{center}
\begin{tabularx}{\linewidth}{|>{\bfseries}X|X|}\hline
\textbf{FACTEURS MODIFIABLES (Comportementaux / Métaboliques)} & \textbf{FACTEURS NON MODIFIABLES (Constitutionnels)} \\ \hline
\begin{itemize}
    \item \textbf{HTA installée} (TA $>$ \SI{140}{\mmHg}/\SI{90}{\mmHg}) : Mécanisme : surcharge de pression $\rightarrow$ lésion endothéliale $\rightarrow$ athérosclérose + hypertrophie ventriculaire gauche.
    \item \textbf{Tabagisme actif} (15~cig/j $\times$ 20~ans = 15~paquet-années) : Mécanisme : CO $\rightarrow$ hypoxie ; Nicotine $\rightarrow$ vasoconstriction, agrégation plaquettaire, dyslipidémie (baisse HDL-c) ; Oxydation LDL-c.
    \item \textbf{Alimentation déséquilibrée} : Excès sel (NaCl) $\rightarrow$ rétention hydrosodée $\rightarrow$ augmentation volémie $\rightarrow$ HTA. Excès graisses saturées (huile de palme, abats, bouillons) $\rightarrow$ augmentation LDL-c $\rightarrow$ athérosclérose. Faible apport fruits/légumes (antioxydants, K$^+$, fibres).
    \item \textbf{Sédentarité / Inactivité physique} : Déconditionnement $\rightarrow$ augmentation FC repos, résistance à l'insuline, prise de poids, dyslipidémie.
    \item \textbf{Surpoids / Obésité probable} (non mesuré mais inféré : IMC probable $>$ \SI{25}{kg/m^2} vu alimentation + inactivité) : Mécanisme : tissu adipeux viscéral inflammatoire (adipokines), résistance à l'insuline, HTA, dyslipidémie.
    \item \textbf{Consommation d'alcool} (régulière week-end) : Effet dose-dépendant sur TA et triglycérides.
    \item \textbf{Stress psychosocial} (charge de travail, insécurité économique) : Activation axe cortico-surrénalien + orthosympathique $\rightarrow$ HTA réactionnelle, troubles du rythme.
\end{itemize}
&
\begin{itemize}
    \item \textbf{Âge} : 50~ans (seuil de risque majeur chez l'homme).
    \item \textbf{Sexe} : Masculin (risque plus précoce que la femme avant ménopause).
    \item \textbf{Hérédité cardiovasculaire} : Père décédé d'AVC à 60~ans (antécédent familial positif ; la précocité est définie à $<$ 55~ans pour le père, $<$ 65~ans pour la mère -- ici non précoce pour le père mais facteur de risque établi). Mère hypertendue. Risque multiplié par 2 à 3.
\end{itemize}
\\ \hline
\end{tabularx}
\end{center}

\textbf{Commentaire pédagogique} : L'association de \textbf{3 facteurs de risque majeurs modifiables} (HTA, Tabagisme, Sédentarité/Alimentation) sur un terrain \textbf{non modifiable défavorable} (Âge, Sexe, Hérédité) place M.~Fotso à un \textbf{risque cardiovasculaire global très élevé} (score de Framingham ou SCORE OMS estimé $>$ 20~\% à 10~ans pour un événement coronarien ou cérébrovasculaire majeur). C'est un candidat prioritaire pour la prévention primaire d'événement majeur.

\vspace{0.5cm}
\textbf{Tâche 3 : Tableau de prévention personnalisé (Facteur $\rightarrow$ Mesure concrète)}

\begin{center}
\begin{tabularx}{\linewidth}{|l|X|}\hline
\textbf{Facteur de risque identifié} & \textbf{Mesure de prévention concrète et contextualisée (Douala, commerçant)} \\ \hline
HTA Stade 1 confirmée & \textbf{Consulter le médecin du district cette semaine} pour confirmation (MAPA ou automesures selon la règle des 3 : 3 matins, 3 soirs, 3 jours), bilan d'atteinte d'organes (ECG, fond d'œil, créatininémie, protéinurie, glycémie, lipidogramme) et \textbf{initiation probable d'un traitement médicamenteux d'emblée} (risque global très élevé) associé aux règles hygiéno-diététiques. Respecter scrupuleusement l'ordonnance. \\ \hline
Tabagisme (15~cig/j) & \textbf{Arrêt total du tabac}. Recourir au programme national de lutte contre le tabagisme (CNLCT) : consultation d'aide à l'arrêt (hôpital Laquintinie ou district), substituts nicotiniques (patchs, gommes) remboursés/gratuits, soutien psychologique. Objectif : 0 cigarette à J+30. \\ \hline
Excès de sel (plats restaurant, bouillons cubes type Maggi/Knorr) & \textbf{Réduire l'apport sodique à $<$ \SI{5}{g/j} (1 c. à café rase)}. Actions concrètes : Ne pas ajouter de sel à table ; Remplacer bouillons cubes par épices locales (penja, gingembre, ail, oignon, céleri) ; Demander au restaurant \og moins de sel, pas de cube \fg ; Privilégier \emph{ndolé} maison sans cube, \emph{kwem} sans sel ajouté. \\ \hline
Excès graisses saturées (huile de palme rouge, peau poulet, abats, fritures) & \textbf{Limiter huile de palme à 1 c. à soupe/jour}. Alterner avec huiles riches en insaturés (arachide, tournesol, soja). Retirer peau volaille, gras visible viandes. Cuisson : braisage, vapeur, papillote $\rightarrow$ pas de friture. 2 à 3 portions poisson/sem (source oméga-3). \\ \hline
Faible apport fruits/légumes & \textbf{5 portions/jour} (\SI{400}{g}). Exemple local : 1 mangue/papaye/ananas au petit-déj ; salade crudités (tomate, concombre, carotte) au déjeuner ; légumes verts cuits (feuilles de manioc, \emph{folong}, \emph{erkang}) au dîner ; jus naturel (bissap, ananas) sans sucre ajouté. \\ \hline
Sédentarité (assise $>$ 10~h, moto-taxi) & \textbf{Activité physique modérée $\ge$ \SI{150}{min/sem}}. Intégrer dans routine : Marche rapide \SI{30}{min} $\times$ 5~j/sem (ex. : aller-retour marché-domicile à pied si $<$ \SI{3}{km}, ou tour du marché avant ouverture). Monter escaliers au lieu d'ascenseur. Étirements \SI{5}{min} toutes les 2~h assis. \\ \hline
Surpoids probable & \textbf{Objectif IMC $<$ \SI{25}{kg/m^2}}. Perte de 5--10~\% du poids initial sur 6~mois (ex. : si \SI{90}{kg} $\rightarrow$ -5 à -9~kg). Pesée hebdomadaire. Journal alimentaire 3~jours/sem. \\ \hline
Alcool (bière week-end) & \textbf{Limiter à $\le$ 2 verres/jour, pas tous les jours}. Alterner avec eau. Jours sans alcool $\ge$ 2/sem. \\ \hline
Stress / Charge mentale & \textbf{Gestion du stress} : Respiration abdominale \SI{5}{min} midi/soir. Temps de pause réel (pas téléphone). Sommeil 7--8~h/nuit. Activité de loisir (lecture, prière, musique) \SI{30}{min/j}. \\ \hline
Hérédité / Âge / Sexe & \textbf{Surveillance renforcée} : Contrôle TA mensuel (automesure ou pharmacie/dispensaire). Bilan lipidique et glycémie à jeun annuels. ECG annuel. Informer frères/sœurs/enfants du risque familial. \\ \hline
\end{tabularx}
\end{center}

\vspace{0.5cm}
\textbf{Tâche 4 : Message de sensibilisation (5 lignes max)}

\textit{Proposition de message SMS / Affiche :}

\begin{quote}
\textbf{« M.~Fotso, votre tension est élevée (15/9). Votre cœur fatigue. Arrêtez le tabac, réduisez le sel et l'huile, marchez 30~min/jour. Venez au centre de santé cette semaine pour protéger votre avenir. Votre famille compte sur vous. »}
\end{quote}

\textit{Analyse du message} : Ton direct (\og vous \fg), chiffre clé (15/9), lien cause-effet simple, 3 actions concrètes (tabac, sel/huile, marche), appel à l'action immédiate (centre de santé), levier motivationnel (famille). Respecte la limite de 5 lignes.
\end{exoresolu}

\courssub{Bilan de l'activité}
Cette activité d'intégration a permis de mettre en évidence plusieurs points clés du programme :

\begin{enumerate}
    \item \textbf{L'articulation savoirs / savoir-faire} : L'interprétation des mesures (TA, FC) n'est pas une simple lecture de chiffres ; elle nécessite la connaissance des seuils diagnostiques (OMS), de l'unité (cmHg vs mmHg, piège fréquent), et de la signification physiopathologique (hypertension systolique prédominante de l'adulte de 50~ans, FC haute = risque).
    \item \textbf{L'approche globale du risque} : Le risque cardiovasculaire ne s'évalue pas facteur par facteur mais de façon \textbf{globale et multiplicative}. M.~Fotso cumule des facteurs \textit{modifiables} (HTA, tabac, nutrition, sédentarité) sur un terrain \textit{non modifiable} défavorable (âge, sexe, hérédité). Cette synergie explique le risque très élevé.
    \item \textbf{La prévention doit être \og sur mesure \fg} : Les recommandations génériques (\og mangez mieux, bougez plus \fg) sont inefficaces. Le tableau construit montre qu'une mesure de prévention doit être \textbf{spécifique} (quoi ?), \textbf{chiffrée} (combien ?), \textbf{contextualisée} (comment au marché de Douala ?) et \textbf{progressivement réalisable} (objectifs SMART).
    \item \textbf{Le rôle de l'éducation à la santé} : Le message de sensibilisation illustre la compétence \og Communiquer \fg. Il ne s'agit pas d'effrayer (discours culpabilisant) mais d'\textbf{empowerment} (autonomisation) : information claire + action concrète + soutien + valorisation de la personne.
    \item \textbf{Le parcours de soins au Cameroun} : L'activité ancre la biologie dans le système de santé réel : dépistage communautaire (infirmier) $\rightarrow$ orientation médecin district $\rightarrow$ confirmation (MAPA, bilan) $\rightarrow$ prise en charge (règles hygiéno-diététiques $\pm$ médicamenteuse) $\rightarrow$ suivi long terme. L'élève comprend qu'il est un acteur de ce parcours (pour lui-même, sa famille, sa communauté).
\end{enumerate}

\begin{aretenir}[Ce qu'il faut retenir de cette intégration]
\begin{itemize}
    \item Une TA $\ge$ \SI{140}{\mmHg}/\SI{90}{\mmHg} (ou 14/9~cmHg) confirmée à deux reprises = \textbf{HTA avérée}.
    \item Une FC de repos $>$ 80~batt/min = marqueur de risque, même si $<$ 100.
    \item Les facteurs de risque s'additionnent et se potentialisent : \textbf{risque global} $>$ somme des risques individuels.
    \item La prévention efficace = \textbf{mesures hygiéno-diététiques prioritaires} (sel, tabac, activité, poids) + \textbf{traitement médicamenteux si indiqué} + \textbf{suivi régulier}.
    \item L'élève de Première C/TI doit savoir \textbf{argumenter scientifiquement} (mécanismes : endothélial, LDL-c, orthosympathique) et \textbf{prescrire concrètement} (contexte local).
\end{itemize}
\end{aretenir}

\begin{attention}[Pièges à éviter au Probatoire]
\begin{itemize}
    \item \textbf{Confondre cmHg et mmHg} : Ne pas écrire \og TA = 16/10 mmHg \fg (c'est \SI{160}{\mmHg}/\SI{100}{\mmHg}). Préciser l'unité utilisée.
    \item \textbf{Oublier de classer les facteurs} : La consigne \og classer en modifiables / non modifiables \fg est systématique. Ne pas mélanger.
    \item \textbf{Proposer des mesures vagues} : \og Faire du sport \fg $\rightarrow$ \textbf{Zéro point}. Exiger : \og Marche rapide 30~min, 5 fois/sem \fg.
    \item \textbf{Méconnaître la définition de la précocité héréditaire} : Père $<$ 55~ans, Mère $<$ 65~ans. Un père AVC à 60~ans est un antécédent familial positif mais \textbf{non précoce} au sens strict. Ne pas écrire \og hérédité précoce \fg pour un père à 60~ans.
    \item \textbf{Séparer prévention et traitement} : Pour une HTA Stade 1 confirmée à risque global très élevé, les règles hygiéno-diététiques sont \textbf{indispensables mais souvent insuffisantes} ; le traitement médicamenteux est fréquemment nécessaire d'emblée. Ne pas dire \og il suffit de marcher pour guérir \fg.
\end{itemize}
\end{attention}

\newpage

\coursec{Méthodes et exercices résolus}

\coursec{Lutte contre les accidents cardiovasculaires}

\courssub{Le système cardiovasculaire : structure et fonctionnement simplifié}

\begin{definition}[Appareil cardiovasculaire]
L'appareil cardiovasculaire assure le transport du sang dans l'organisme. Il comprend :
\begin{itemize}
\item le \textbf{cœur}, pompe musculaire creuse à quatre cavités (deux oreillettes, deux ventricules) située dans le médiastin ;
\item les \textbf{vaisseaux sanguins} : artères (fortes parois, haute pression), veines (parois fines, valves, basse pression), capillaires (parois mono-cellulaires, site des échanges) ;
\item le \textbf{sang}, fluide vecteur (globules rouges, blancs, plaquettes, plasma).
\end{itemize}
\end{definition}

\begin{propriete}[Double circulation et cycle cardiaque]
Le sang effectue une \textbf{double circulation} à chaque tour complet :
\begin{enumerate}
\item \textbf{Circulation pulmonaire (petite circulation)} : Ventricule droit $\rightarrow$ Artère pulmonaire $\rightarrow$ Poumons (hématose : $\ce{CO2}$ sorti, $\ce{O2}$ entré) $\rightarrow$ Veines pulmonaires $\rightarrow$ Oreillette gauche.
\item \textbf{Circulation systémique (grande circulation)} : Ventricule gauche $\rightarrow$ Aorte $\rightarrow$ Organes (échanges capillaires) $\rightarrow$ Veines caves $\rightarrow$ Oreillette droite.
\end{enumerate}
Le \textbf{cycle cardiaque} alterne \textbf{diastole} (relâchement, remplissage ventriculaire, pression basse) et \textbf{systole} (contraction, éjection, pression haute). Les valves (mitrale, tricuspide, aortique, pulmonaire) assurent l'unidirectionnalité du flux.
\end{propriete}

\begin{experience}[Observation de la structure cardiaque et mesure du pouls]
\textbf{Matériel} : Cœur de bœuf frais (ou modèle anatomique), gants, scalpel, plateau ; tensiomètre manuel (baumanomètre) et stéthoscope.
\textbf{Protocole} :
\begin{enumerate}
\item \textbf{Dissection/Observation} : Identifier les 4 cavités, l'épaisseur des parois (ventricule gauche $\gg$ ventricule droit), les valves, les artères coronaires (sur la surface du cœur).
\item \textbf{Mesure de la tension artérielle} : Sujet au repos 5 min, assis, brassard au niveau du cœur. Gonfler au-dessus de la disparition du pouls radial, dégonfler lentement (2--3 mmHg/s). Noter la 1ère bruit de Korotkoff (systolique) et la disparition (diastolique).
\item \textbf{Mesure de la fréquence cardiaque} : Palper l'artère radiale, compter 60 s (ou 15 s $\times$ 4).
\end{enumerate}
\textbf{Observations} : Paroi ventriculaire gauche très épaisse (pression systémique $\approx$ 120 mmHg) vs droite fine (pression pulmonaire $\approx$ 25 mmHg). Bruits cardiaques « tou-toup » (fermeture valves AV puis valves sigmoïdes). Pouls radial synchrone de la systole.
\textbf{Interprétation} : La structure anatomique (épaisseur, valves, coronaires) est adaptée à la fonction de pompe à double pression. La tension artérielle reflète la pression exercée par le sang sur les parois artérielles lors de la systole et de la diastole.
\end{experience}

\begin{aretenir}[Vocabulaire clé du Probatoire]
\begin{itemize}
\item \cle{Systole} : Contraction ventriculaire $\rightarrow$ Éjection du sang $\rightarrow$ Pression maximale (systolique).
\item \cle{Diastole} : Relâchement ventriculaire $\rightarrow$ Remplissage $\rightarrow$ Pression minimale (diastolique).
\item \cle{Artères coronaires} : Irriguent le myocarde ; leur obstruction cause l'infarctus.
\item \cle{Débit cardiaque} ($Q_c$) = Fréquence cardiaque ($f_c$) $\times$ Volume d'éjection systolique ($V_{es}$). Unité : L/min.
\end{itemize}
\end{aretenir}

\courssub{La tension artérielle et sa mesure}

\begin{methode}[Lire et interpréter une tension artérielle]
La tension artérielle (TA) s'exprime en \textbf{centimètres de mercure (cmHg)} : exemple 12/8.
\begin{enumerate}
\item \textbf{Identifier les deux valeurs} :
\begin{itemize}
\item \textbf{Premier chiffre} = \cle{Tension systolique} (pression maximale, systole ventriculaire) ;
\item \textbf{Second chiffre} = \cle{Tension diastolique} (pression résiduelle, diastole ventriculaire).
\end{itemize}
\item \textbf{Comparer aux valeurs de référence} (adulte au repos, moyenne de plusieurs mesures) :
\begin{itemize}
\item \textbf{Normale} : Systolique 11--13 \textbf{ET} Diastolique 7--9 (ex. 12/8) ;
\item \textbf{Hypertension artérielle (HTA)} : Systolique $\geq 14$ \textbf{OU} Diastolique $\geq 9$ (confirmée à 3 mesures rapprochées) ;
\item \textbf{Hypotension} : Systolique $< 10$.
\end{itemize}
\item \textbf{Vérifier les conditions de validité} : Repos 5 min, position assise, brassard au niveau du cœur, pas de café/tabac/effort 30 min avant. \textbf{Une seule mesure ne suffit jamais pour diagnostiquer une HTA} (effet « blouse blanche », stress, variabilité physiologique).
\item \textbf{Conclure} : Préciser l'état tensionnel (normotendu, hypertendu, hypotendu) et la conduite à tenir (surveillance, hygiène de vie, consultation).
\end{enumerate}
\textbf{Réflexe rédaction} : Écrire « TA = 14/9 cmHg » et non « 14 sur 9 ». Préciser toujours l'unité (cmHg). 1 cmHg = 10 mmHg.
\end{methode}

\begin{exoresolu}[Interprétation de mesures tensionnelles en campagne de dépistage]
Lors d'une campagne au Lycée de Biyem-Assi (Yaoundé), trois adultes au repos donnent : Mme A : 12/7 ; M. B : 15/9 ; M. C : 9/6 (en cmHg).
\textbf{Questions :} 1) Identifier systolique et diastolique. 2) Interpréter. 3) Conduite pour M. B.
\tcblower
\textbf{1) Identification}
\begin{center}
\begin{tabularx}{\linewidth}{|l|X|X|X|}
\hline
\rowcolor{popblueL} \textbf{Personne} & \textbf{Systolique (cmHg)} & \textbf{Diastolique (cmHg)} & \textbf{Interprétation} \\
\hline
Mme A & 12 & 7 & Normale \\
\hline
M. B & 15 & 9 & Hypertension artérielle \\
\hline
M. C & 9 & 6 & Hypotension \\
\hline
\end{tabularx}
\end{center}

\textbf{2) Interprétation détaillée}
\begin{itemize}
\item \textbf{Mme A (12/7)} : Dans les normes (S 11--13, D 7--9). Perfusion tissulaire optimale sans surcharge pariétale.
\item \textbf{M. B (15/9)} : S $\geq 14$ \textbf{et} D $\geq 9$ : \textbf{HTA stade 1}. Risque majoré d'athérosclérose, infarctus, AVC, néphropathie. Nécessite confirmation.
\item \textbf{M. C (9/6)} : S $< 10$ : \textbf{Hypotension}. Souvent asymptomatique, mais risque de malaise vagal (hypoperfusion cérébrale) à la station debout prolongée.
\end{itemize}

\textbf{3) Conduite pour M. B}
Confirmer par \textbf{automesures} (règle des 3 : 3 mesures le matin, 3 le soir, 3 jours de suite) ou MAPA (Mesure Ambulatoire de la Pression Artérielle). Si confirmé : consultation médicale, régime hyposodé ($< 5$ g sel/j), activité physique 30 min/j, arrêt tabac/alcool, traitement antihypertenseur si indiqué (observance stricte).
\end{exoresolu}

\courssub{Les principaux accidents cardiovasculaires}

\begin{definition}[Accident cardiovasculaire]
Survient lorsque l'irrigation sanguine d'un organe vital est brutalement interrompue (ischémie) ou qu'une artère se rompt (hémorragie), entraînant une \textbf{nécrose} (mort cellulaire irréversible par anoxie) du territoire en aval.
\end{definition}

\begin{propriete}[Mécanismes des trois accidents majeurs]
\begin{center}
\begin{tabularx}{\linewidth}{|l|X|X|X|}
\hline
\rowcolor{popblueL} \textbf{Accident} & \textbf{Siège \& Mécanisme} & \textbf{Signes d'alerte (Urgence vitale)} & \textbf{Conséquences} \\
\hline
\textbf{Infarctus du myocarde (IDM)} & Obstruction (thrombus sur plaque d'athérome) d'une \textbf{artère coronaire} $\rightarrow$ nécrose du \textbf{myocarde}. & Douleur thoracique intense, constrictive, irradiant bras gauche/mâchoire/dos, sueurs froides, angoisse, essoufflement. & Arythmie, insuffisance cardiaque, mort subite, cicatrice fibreuse (séquelle). \\
\hline
\textbf{Accident Vasculaire Cérébral (AVC)} & \textbf{Ischémique (80\%)} : Thrombose/embolie artère cérébrale. \textbf{Hémorragique (20\%)} : Rupture anévrisme/HTA. $\rightarrow$ Nécrose tissu cérébral. & \textbf{FAST} : \textbf{F}ace (bouche déviée), \textbf{A}rm (paralysie hémicorps), \textbf{S}peech (aphasie), \textbf{T}ime (appeler 112/119 SAMU). & Hémiparésie, aphasie, troubles cognitifs, décès. Séquelles lourdes si prise en charge $> 4$h30 (thrombolyse). \\
\hline
\textbf{Insuffisance cardiaque (IC)} & Cœur incapable d'assurer le débit nécessaire. Cause fréquente : \textbf{HTA chronique} (surcharge de pression) ou séquelle d'IDM. & Dyspnée d'effort puis repos, œdèmes membres inférieurs, fatigue, orthopnée, prise de poids rapide. & Décompensations aiguës (œdème pulmonaire aigu), mortalité élevée à 5 ans sans traitement. \\
\hline
\end{tabularx}
\end{center}
\end{propriete}

\begin{methode}[Expliquer le mécanisme d'un infarctus ou AVC (Chaîne causale)]
Toujours structurer la réponse en \textbf{4 étapes logiques} avec connecteurs (\emph{donc, par conséquent, ce qui entraîne}) :
\begin{enumerate}
\item \textbf{Terrain (Facteurs de risque)} : Tabac, HTA, dyslipidémie, diabète $\rightarrow$ lésion de l'\textbf{endothélium} (paroi interne artères).
\item \textbf{Lésion (Athérosclérose)} : Dépôt de cholestérol LDL $\rightarrow$ \textbf{plaque d'athérome} (lipides, cellules inflammatoires, fibres) $\rightarrow$ rétrécissement lumière (sténose) + fragilisation paroi.
\item \textbf{Événement aigu (Accident)} : \textbf{Rupture de la plaque} (fissuration calotte fibreuse) $\rightarrow$ contact sang/colonne lipidique $\rightarrow$ activation plaquettes/coagulation $\rightarrow$ \textbf{thrombus} (caillot) occlusif \textbf{OU} rupture artérielle (hémorragie).
\item \textbf{Conséquence (Nécrose)} : Arrêt brutal perfusion $\rightarrow$ anoxie $\rightarrow$ arrêt respiration cellulaire ($\ce{ATP}$ chute) $\rightarrow$ \textbf{nécrose} (myocarde ou neurones) en quelques heures. Irréversible.
\end{enumerate}
\end{methode}

\begin{exoresolu}[Mécanisme d'un AVC ischémique]
Mme K., 65 ans, diabétique, hypertendue mal équilibrée, fait un AVC ischémique (hémiparésie droite, aphasie).
\textbf{Question :} Expliquer la chaîne causale.
\tcblower
\begin{enumerate}
\item \textbf{Terrain} : Diabète (hyperglycémie toxique pour endothélium) + HTA (contrainte mécanique pariétale) $\rightarrow$ lésion endothéliale artères carotides/cérébrales.
\item \textbf{Lésion} : Athérosclérose carotidienne : plaques d'athérome rétrécissant la lumière, source d'emboles.
\item \textbf{Accident} : Rupture d'une plaque carotidienne $\rightarrow$ formation thrombus $\rightarrow$ migration embolique (embole) dans circulation cérébrale $\rightarrow$ obstruction artère cérébrale moyenne (gauche).
\item \textbf{Conséquence} : Ischémie territoire artère cérébrale moyenne gauche $\rightarrow$ nécrose neurones hémisphère gauche (aire motrice controlatérale + aire de Broca) $\rightarrow$ hémiparésie droite + aphasie expressive.
\end{enumerate}
\textbf{Urgence} : Thrombolyse (rt-PA) possible $< 4$h30 ; thrombectomie mécanique $< 6$h (jusqu'à 24h si imagerie favorable). « Time is brain ».
\end{exoresolu}

\courssub{Les facteurs de risque cardiovasculaire}

\begin{definition}[Facteur de risque cardiovasculaire]
Caractéristique ou mode de vie associé statistiquement à une augmentation de la probabilité de survenue d'un accident cardiovasculaire. Ils se \textbf{cumulent de façon multiplicative} (synergie), non additive.
\end{definition}

\begin{propriete}[Classification des facteurs de risque]
\begin{center}
\begin{tabularx}{\linewidth}{|l|X|X|}
\hline
\rowcolor{popblueL} \textbf{Catégorie} & \textbf{Facteurs} & \textbf{Mécanisme pathophysiologique principal} \\
\hline
\textbf{Non modifiables} & \textbf{Âge} ($\uparrow$ risque apr. 50 ans H, 60 ans F) & Vieillissement vasculaire (rigidité artérielle, athérosclérose sénile). \\
& \textbf{Sexe} (H $>$ F avant ménopause) & Protection œstrogénique pré-ménopause (vasodilatation, HDL $\uparrow$). \\
& \textbf{Hérédité} (antécédent IDM/AVC parent $<$ 55 ans H, $<$ 65 ans F) & Polymorphismes génétiques (métabolisme lipides, coagulation, tension). \\
\hline
\textbf{Modifiables (Cibles prévention)} & \textbf{Tabagisme} (actif/passif) & Endommagement endothélium, $\uparrow$ agrégation plaquettaire, $\downarrow$ HDL, vasoconstriction (nicotine), $\uparrow$ CO (hypoxie). \\
& \textbf{Hypertension artérielle} & Surcharge pression $\rightarrow$ hypertrophie ventriculaire gauche + lésion endothéliale $\rightarrow$ athérosclérose accélérée. \\
& \textbf{Dyslipidémie} (LDL-cholestérol $\uparrow$, HDL $\downarrow$) & LDL $\rightarrow$ intima artère $\rightarrow$ oxydation $\rightarrow$ plaque d'athérome. HDL $\rightarrow$ transport inverse cholestérol (protecteur). \\
& \textbf{Diabète type 2} & Glycation protéines (produits de Maillard) $\rightarrow$ rigidité artérielle, dysfonction endothéliale, état pro-thrombotique. \\
& \textbf{Obésité abdominale} (Tour de taille H $> 94$ cm, F $> 80$ cm) & Adiposité viscérale $\rightarrow$ cytokines inflammatoires, résistance insuline, HTA, dyslipidémie. \\
& \textbf{Sédentarité} & $\downarrow$ dépense énergétique, $\downarrow$ sensibilité insuline, $\downarrow$ HDL, $\uparrow$ rigidité artérielle. \\
& \textbf{Alcool} (excès $> 2$ verres/j H, $> 1$ F) & $\uparrow$ TA, $\uparrow$ triglycérides, cardiomyopathie dilatée, arythmies. \\
& \textbf{Stress psychosocial chronique} & Activation sympathique/axe corticosurrénalien $\rightarrow$ $\uparrow$ TA, $\uparrow$ coagulabilité, troubles rythme. \\
\hline
\end{tabularx}
\end{center}
\end{propriete}

\begin{methode}[Identifier et classer les facteurs dans une situation (Savoir-faire 1)]
\begin{enumerate}
\item \textbf{Lister exhaustivement} tous les indices du texte (âge, sexe, habitudes, antécédents, biologie).
\item \textbf{Classer} dans les deux colonnes : Non modifiables / Modifiables.
\item \textbf{Relier chaque facteur modifiable à son mécanisme} (ex. : « Tabac $\rightarrow$ lésion endothéliale $\rightarrow$ athérosclérose »).
\item \textbf{Évaluer le risque global} : Cumul = multiplication des risques relatifs. Un patient avec HTA + Tabac + Diabète a un risque $\times 10$ à $\times 20$ vs sans facteur.
\end{enumerate}
\end{methode}

\begin{exoresolu}[Analyse complète : Cas de M. N., 52 ans, Douala]
\textit{Commerçant, 1 paquet/j 30 ans, ndolè/huile palme/sel quotidiens, sédentaire, stress permanent, père IDM 55 ans, TA 16/10 cmHg.}
\tcblower
\textbf{Facteurs non modifiables :} Âge 52 ans (risque $\uparrow$) ; Hérédité (père IDM précoce).
\textbf{Facteurs modifiables (7 majeurs) :}
\begin{enumerate}
\item \textbf{Tabagisme majeur} (30 PA) $\rightarrow$ Endothélium lésé, thrombose, vasoconstriction.
\item \textbf{HTA sévère} (16/10) $\rightarrow$ Surcharge pression, rupture plaques, hypertrophie VG.
\item \textbf{Dyslipidémie probable} (Huile palme/grasses saturées) $\rightarrow$ LDL $\uparrow$, athérome coronarien/carotidien.
\item \textbf{Excès sel} (Plats salés) $\rightarrow$ Rétention Na$^+$/Eau $\rightarrow$ Volémie $\uparrow$ $\rightarrow$ TA $\uparrow$.
\item \textbf{Sédentarité} $\rightarrow$ Prise de poids, résistance insuline, $\downarrow$ HDL.
\item \textbf{Stress chronique} $\rightarrow$ Catécholamines $\uparrow$ $\rightarrow$ TA $\uparrow$, ischémie par vasoconstriction coronarienne.
\item \textbf{Surpoids/Obésité abdominale probable} (alimentation + sédentarité) $\rightarrow$ Inflammation bas bruit.
\end{enumerate}
\textbf{Conclusion :} Risque \textbf{très élevé} (Score SCORE/Frammingham $\gg 10\%$ à 10 ans). Prise en charge globale urgente.
\end{exoresolu}

\begin{methode}[Exploiter des données épidémiologiques (Savoir-faire 6)]
\begin{enumerate}
\item \textbf{Décrire} : Population, effectifs, durée suivi, variable étudiée (incidence, mortalité).
\item \textbf{Calculer} : Fréquence ($\%$), Risque Relatif (RR = Fréquence exposés / Fréquence non-exposés), Risque Attributable.
\item \textbf{Comparer} : Groupes exposés vs non-exposés (fumeurs vs non-fumeurs, HTA vs normo).
\item \textbf{Interpréter biologiquement} : Lier l'excès de risque aux mécanismes (athérose, thrombose).
\item \textbf{Nuancer} : Corrélation $\neq$ Causalité (biais, confusions), mais cohérence avec physiopathologie $\rightarrow$ niveau de preuve fort.
\end{enumerate}
\end{methode}

\begin{exoresolu}[Étude de cohorte : Impact du mode de vie]
Suivi 10 ans, 10 000 adultes.
\begin{center}
\begin{tabularx}{\linewidth}{|l|X|X|}
\hline
\rowcolor{popblueL} \textbf{Groupe} & \textbf{Effectif} & \textbf{Accidents CV (IDM+AVC)} \\
\hline
Non-fumeurs, actifs, équilibré & 4 000 & 40 \\
\hline
Fumeurs, sédentaires, gras/salé & 2 000 & 120 \\
\hline
\end{tabularx}
\end{center}
\textbf{Questions :} Fréquences, Comparaison, Conclusion.
\tcblower
\textbf{Fréquences :} $f_1 = \frac{40}{4000}\times 100 = 1\%$ ; $f_2 = \frac{120}{2000}\times 100 = 6\%$.
\textbf{Risque Relatif :} $RR = \frac{6\%}{1\%} = 6$. Le risque est \textbf{multiplié par 6}.
\textbf{Interprétation :} L'association tabac/sédentarité/alimentation délétère multiplie le risque par 6 via l'athérosclérose accélérée et l'HTA. Cela justifie l'efficacité de la prévention primaire.
\end{exoresolu}

\courssub{La prévention des accidents cardiovasculaires}

\begin{propriete}[Niveaux de prévention]
\begin{itemize}
\item \textbf{Primaire} : Éviter l'apparition des facteurs de risque (éducation santé dès l'école, environnement sans tabac, urbanisme actif, fiscalité sel/sucre/alcool). \emph{Ex: Campagne « Zéro tabac » dans les lycées camerounais.}
\item \textbf{Secondaire} : Dépister et traiter précocement les facteurs de risque (HTA, diabète, dyslipidémie) et maladies asymptomatiques pour éviter l'accident. \emph{Ex: Journée mondiale HTA (17 mai) : dépistage gratuit centres de santé.}
\item \textbf{Tertiaire} : Limiter récidives et séquelles après un accident (réadaptation cardiaque, observance traitements : antiagrégants, statines, IEC/ARA2, bêta-bloquants). \emph{Ex: Unité de réadaptation cardiaque CHU Yaoundé.}
\end{itemize}
\end{propriete}

\begin{methode}[Proposer un plan de prévention personnalisé (Savoir-faire 5)]
\begin{enumerate}
\item \textbf{Cibler les facteurs modifiables} identifiés (ignorer âge/hérédité pour l'action).
\item \textbf{Prescrire des mesures SMART} (Spécifiques, Mesurables, Atteignables, Réalistes, Temporelles) :
\begin{itemize}
\item \textbf{Tabac} : Arrêt total (substitution nicotinique, varénicline, soutien psychologique). Risque IDM divisé par 2 à 1 an.
\item \textbf{Alimentation} : Régime méditerranéen / DASH. Sel $< 5$ g/j (1 c. à café), $\uparrow$ fruits/légumes (5/j), poisson 2$\times$/sem, huiles insaturées (arachide, coton, olive), $\downarrow$ gras saturés (huile palme, beurre, fritures), $\downarrow$ sucres ajoutés.
\item \textbf{Activité physique} : $\geq 150$ min/sem endurance modérée (marche rapide 30 min $\times$ 5/j) + renforcement 2$\times$/sem. Adapter si pathologie.
\item \textbf{Poids} : Viser IMC 18,5--25, Tour taille H $< 94$ cm / F $< 80$ cm.
\item \textbf{Alcool} : $\leq 2$ verres/j (H), $\leq 1$ (F), jours d'abstinence.
\item \textbf{Stress/Sommeil} : 7--8 h/nuit, cohérence cardiaque, loisirs.
\end{itemize}
\item \textbf{Intégrer le dépistage/suivi} : TA $\geq$ 1/an (ou automesure), Glycémie/Cholestérol $\geq$ 1/an si facteurs, ECG repos si symptômes.
\item \textbf{Insister sur l'observance} : Traitement antihypertenseur/hypolipémiant/antidiabétique \textbf{à vie}, même si « je me sens bien ». Arrêt = rebond tensionnel/risque majoré.
\end{enumerate}
\end{methode}

\begin{exoresolu}[Plan de prévention pour M. N. (cas précédent)]
\tcblower
\textbf{Plan d'action priorisé et négocié avec le patient :}
\begin{enumerate}
\item \textbf{Urgence : Arrêt tabac} (J-0). Prescription substituts nicotiniques (patchs/gommes) + counseling. Objectif : 0 cigarette à J+7.
\item \textbf{HTA : Traitement médicamenteux immédiat} (ex. IEC + Diurétique ou ARA2 + CCB) + Automasure 3$\times$3. Objectif TA $< 14/9$ à 3 mois.
\item \textbf{Nutrition :} Consultation diététicienne. Remplacer huile palme friture par huile araide/coton crue ; ndolè moins gras (viande maigre, moins d'huile) ; sel $\rightarrow$ épices/ail/oignon/poivre. Objectif sel $< 5$ g/j.
\item \textbf{Activité :} Marche 30 min matin (avant chaleur) 5$\times$/sem. Intégrer dans trajet marché/travail.
\item \textbf{Suivi :} RDV mensuel 3 mois (TA, poids, observance, effets secondaires), puis trimestriel. Bilan lipidique/glycémie à J+30.
\item \textbf{Éducation thérapeutique} : Expliquer le « pourquoi » de chaque pilule (ex. : « Ce cachet protège vos reins et votre cœur en baissant la pression dans vos artères »).
\end{enumerate}
\textbf{Pronostic :} Si observance stricte, risque IDM/AVC réduit de 50--80\% en 5 ans.
\end{exoresolu}

\courssub{Synthèse et évaluation des acquis}

\begin{aretenir}[Fiche de révision Probatoire - Points clés]
\begin{itemize}
\item \textbf{TA Normale} : 12/8 cmHg (S 11--13, D 7--9). \textbf{HTA} : $\geq 14/9$ (confirmée). \textbf{Unité} : cmHg.
\item \textbf{FC Repos} : 60--100 bpm. \textbf{FC Max théorique} : $220 - \text{âge}$. Calcul pouls 15s $\times$ 4.
\item \textbf{Infarctus} = Obstruction \textbf{coronaire} $\rightarrow$ Nécrose \textbf{myocarde}. \textbf{AVC} = Obstruction/Rupture \textbf{cérébrale} $\rightarrow$ Nécrose \textbf{cerveau}.
\item \textbf{Athérosclérose} : Plaque d'\textbf{athérome} (cholestérol LDL, macrophages, fibres) sur \textbf{intima} artères moyennes/grosses.
\item \textbf{Thrombose} : Caillot (\textbf{thrombus}) sur plaque rompue $\rightarrow$ occlusion aiguë.
\item \textbf{Facteurs modifiables (9)} : Tabac, HTA, LDL$\uparrow$, Diabète, Obésité abdominale, Sédentarité, Alcool, Stress, Alimentation (Sel/Gras saturés).
\item \textbf{Prévention} : Hygiène de vie (Alimentation DASH, Activité, Arrêt tabac) + Dépistage + Observance traitements.
\end{itemize}
\end{aretenir}

\begin{attention}[Les 5 erreurs fatales au Probatoire]
\begin{enumerate}
\item \textbf{Inverser Systolique/Diastolique} : Écrire « TA = 8/12 » ou dire « le petit chiffre c'est la systole ». \textbf{Zéro point}. Systolique = Grand chiffre = Contraction.
\item \textbf{Diagnostiquer HTA sur 1 mesure} : Oublier « confirmée par plusieurs mesures au repos ». Perte point rigueur.
\item \textbf{Confondre Infarctus / AVC} : « Infarctus cérébral » sans préciser AVC ischémique ; situer l'infarctus au cerveau. \textbf{Terminologie exacte exigée} : Artère \textbf{coronaire} $\leftrightarrow$ Myocarde ; Artère \textbf{cérébrale} $\leftrightarrow$ Cerveau.
\item \textbf{Erreur calcul FC} : Compter 18 battements/15s $\rightarrow$ écrire « FC = 18 bpm ». \textbf{Formule} : $FC = \frac{N}{t(s)} \times 60$. Ici $18 \times 4 = 72$ bpm.
\item \textbf{Liste de facteurs sans mécanisme} : « Tabac, sel, gras sont des facteurs ». \textbf{Insuffisant}. Exigé : « Tabac $\rightarrow$ lésion endothélium $\rightarrow$ plaque athérome $\rightarrow$ sténose/thrombose $\rightarrow$ IDM/AVC ». Vocabulaire : \textbf{athérome, thrombus, nécrose, endothélium, LDL, HDL}.
\end{enumerate}
\end{attention}

\begin{exercice}[Entraînement Probatoire - Type Bac]
\textbf{Exercice 1} (Interprétation TA/FC) : Un élève de 17 ans mesure sa TA au repos : 13/8 cmHg. Son pouls radial : 22 battements en 15 s.
\begin{enumerate}
\item Préciser sa tension systolique et diastolique. Interpréter.
\item Calculer sa fréquence cardiaque. Interpréter.
\item Après un effort modéré, sa FC monte à 150 bpm. Calculer sa FC max théorique. Cette FC d'effort est-elle physiologique ? Justifier.
\end{enumerate}

\textbf{Exercice 2} (Chaîne causale) : M. X, 60 ans, fumeur, diabétique, présente une douleur thoracique brutale à l'effort, irradiant au bras gauche, cédant au repos. L'ECG montre un sus-décalage du segment ST.
\begin{enumerate}
\item Nommer l'accident cardiovasculaire probable.
\item En 4 étapes numérotées, expliquer le mécanisme physiopathologique en utilisant les termes : endothélium, plaque d'athérome, thrombus, nécrose, anoxie.
\item Pourquoi la prise en charge doit-elle être immédiate (délai < 2h idéalement) ?
\end{enumerate}

\textbf{Exercice 3} (Prévention \& Épidémiologie) : On compare l'incidence de l'IDM chez 5000 hommes de 50 ans suivis 5 ans.
\begin{center}
\begin{tabular}{|c|c|c|}\hline
\textbf{Groupe} & \textbf{Effectif} & \textbf{Nouveaux IDM} \\ \hline
Non-fumeurs, normotendus, actifs & 2500 & 10 \\ \hline
Fumeurs, hypertendus, sédentaires & 2500 & 75 \\ \hline
\end{tabular}
\end{center}
\begin{enumerate}
\item Calculer l'incidence (pour 1000 personnes-années) dans chaque groupe.
\item Calculer le Risque Relatif. Interpréter.
\item Citer 3 mesures de prévention primaire ciblant le groupe à risque.
\end{enumerate}
\end{exercice}

\begin{corrige}[Exercice 1]
\begin{enumerate}
\item Systolique = 13 cmHg, Diastolique = 8 cmHg. TA \textbf{normale} (dans fourchettes 11--13 / 7--9).
\item $FC = \frac{22}{15} \times 60 = 88$ bpm. \textbf{Normale} (60--100 bpm).
\item $FC_{max} = 220 - 17 = 203$ bpm. $150 < 203$ : \textbf{Oui, physiologique}. Adaptation normale : $\uparrow$ débit cardiaque pour $\uparrow$ apport $\ce{O2}$ muscles (respiration cellulaire $\ce{C6H12O6 + 6O2 -> 6CO2 + 6H2O + ATP}$). Récupération attendue $< 100$ bpm à 5 min post-effort.
\end{enumerate}
\end{corrige}

\begin{corrige}[Exercice 2]
\begin{enumerate}
\item \textbf{Angor de poitrine instable} (ou IDM si marqueurs nécrose $\uparrow$ ; ici cédant au repos + sus-decalage ST $\rightarrow$ ischémie transitoire mais haut risque d'IDM imminent).
\item \textbf{Chaîne causale :}
\begin{enumerate}
\item \textbf{Terrain} : Tabac + Diabète $\rightarrow$ dysfonction endothéliale artères coronaires.
\item \textbf{Lésion} : Infiltration LDL oxydé $\rightarrow$ formation \textbf{plaque d'athérome} sténosante sur intima coronaire.
\item \textbf{Accident} : Effort $\rightarrow$ demande $\ce{O2}$ $\uparrow$ + stress hémodynamique $\rightarrow$ fissuration plaque $\rightarrow$ agrégation plaquettes/coagulation $\rightarrow$ \textbf{thrombus} occlusif partiel/instable.
\item \textbf{Conséquence} : Débit coronaire insuffisant $\rightarrow$ \textbf{anoxie} myocarde en aval $\rightarrow$ ischémie (douleur, sus-decalage ST). Si persiste $\rightarrow$ \textbf{nécrose} (IDM).
\end{enumerate}
\item \textbf{Urgence} : La nécrose myocardique est \textbf{irréversible} et s'étend de l'endocarde vers l'épicarde (« vague de nécrose »). La reperfusion (thrombolyse/angioplastie) dans les \textbf{premières heures} sauve le myocarde en zone pénombre (ischémique mais viable). « Time is muscle ». Risque mortel : fibrillation ventriculaire, rupture paroi libre, insuffisance cardiaque aiguë.
\end{enumerate}
\end{corrige}

\begin{corrige}[Exercice 3]
\begin{enumerate}
\item Incidence = $\frac{\text{Nouveaux cas}}{\text{Effectif} \times \text{Durée (années)}} \times 1000$.
Groupe 1 : $\frac{10}{2500 \times 5} \times 1000 = \frac{10}{12500} \times 1000 = \mathbf{0,8}$ pour 1000 pers-années.
Groupe 2 : $\frac{75}{2500 \times 5} \times 1000 = \frac{75}{12500} \times 1000 = \mathbf{6,0}$ pour 1000 pers-années.
\item $RR = \frac{6,0}{0,8} = \mathbf{7,5}$. Le risque d'IDM est \textbf{7,5 fois plus élevé} chez les fumeurs hypertendus sédentaires.
\item \textbf{Prévention primaire ciblée :} 1) Programme national arrêt tabac (taxation, paquet neutre, consultation sevrage). 2) Dépistage/traitement HTA systématique (accès médicaments génériques gratuits/peu coûteux). 3) Aménagement urbain (pistes cyclables, espaces verts) + Éducation physique scolaire obligatoire pour lutter sédentarité.
\end{enumerate}
\end{corrige}

\newpage

\coursec{Exercices}

\coursec{Lutte contre les accidents cardiovasculaires}

\courssub{Le système cardiovasculaire : structure et fonctionnement}

\begin{experience}[Observation : Le cœur, une pompe à double circuit]
\textbf{Matériel :} Cœur de bœuf frais (ou modèle plastique), gants, plateau, scalpel, sonde, eau colorée (rouge/bleue), seringue.
\textbf{Protocole :}
\begin{enumerate}
    \item Observer le cœur \emph{in situ} : repérer les gros vaisseaux (aorte, artère pulmonaire, veines caves, veines pulmonaires) et les quatre cavités.
    \item Inciser le long du septum interventriculaire pour ouvrir les ventricules. Observer l'épaisseur des parois (ventricule gauche \textgreater\ ventricule droit), les valves (sigmoïdes, mitrale, tricuspide), les piliers et cordages.
    \item Injecter de l'eau rouge dans l'oreillette gauche (via veines pulmonaires) : elle sort par l'aorte. Injecter de l'eau bleue dans l'oreillette droite (via veines caves) : elle sort par l'artère pulmonaire.
\end{enumerate}
\textbf{Observations :} Le sang ne mélange pas. Circuit court (cœur-poumons-cœur) et circuit long (cœur-organes-cœur). Paroi VG très épaisse (pression forte), paroi VD fine (pression faible).
\textbf{Interprétation :} Le cœur assure la double circulation : \textbf{circulation pulmonaire} (petite circulation, sang désoxygéné $\to$ oxygéné) et \textbf{circulation systémique} (grande circulation, sang oxygéné $\to$ désoxygéné). Les valves assurent l'unidirectionnalité du flux.
\end{experience}

\begin{popfigure}
\begin{tikzpicture}[scale=1.1, every node/.style={font=\small}]
% --- Coordinates setup ---
% Right Heart (Anatomical Right = Viewer Left) : x < 0
% Left Heart (Anatomical Left = Viewer Right) : x > 0
% Atria y ~ 2.5 to 3.5, Ventricles y ~ -0.5 to 2.0

% PERICARDIUM / OUTLINE (simplified)
\draw[popdark, thick, fill=popcream] (-3, -1) .. controls (-3.5, 1) and (-3.5, 3) .. (-1.5, 3.8) .. controls (0, 4.2) .. (1.5, 3.8) .. controls (3.5, 3) and (3.5, 1) .. (3, -1) .. controls (3, -1.5) and (-3, -1.5) .. (-3, -1);

% SEPTUM
\draw[popdark, thick] (0, -1) -- (0, 3.5);

% RIGHT ATRIUM (RA) - Viewer Left
\draw[popblue, thick, fill=popblue!10] (-2.8, 2.5) .. controls (-2.8, 3.5) and (-1.2, 3.5) .. (0, 3.5) -- (0, 2.0) .. controls (-1.2, 2.0) and (-2.8, 2.0) .. (-2.8, 2.5);
% RIGHT VENTRICLE (RV) - Viewer Left
\draw[popblue, thick, fill=popblue!10] (-2.5, -0.5) .. controls (-2.8, 1) and (-1.5, 2.0) .. (0, 2.0) -- (0, -0.5) -- (-2.5, -0.5);

% LEFT ATRIUM (LA) - Viewer Right
\draw[poporange, thick, fill=poporange!10] (0, 3.5) .. controls (1.2, 3.5) and (2.8, 3.5) .. (2.8, 2.5) .. controls (2.8, 2.0) and (1.2, 2.0) .. (0, 2.0) -- (0, 3.5);
% LEFT VENTRICLE (LV) - Viewer Right
\draw[poporange, thick, fill=poporange!10] (0, -0.5) -- (0, 2.0) .. controls (1.5, 2.0) and (2.8, 1) .. (2.5, -0.5) -- (0, -0.5);

% VALVES (schematic)
% Tricuspid (RA-RV)
\draw[popdark, dashed] (-1.8, 2.0) -- (-0.2, 2.0);
% Mitral (LA-LV)
\draw[popdark, dashed] (0.2, 2.0) -- (1.8, 2.0);
% Pulmonary Valve (RV exit)
\draw[popdark, thick] (0.2, 2.8) -- (1.0, 3.2);
% Aortic Valve (LV exit)
\draw[popdark, thick] (-0.2, 2.8) -- (-1.0, 3.2);

% VESSELS
% VCS (Superior Vena Cava) -> RA (Top Left)
\draw[popblue, very thick] (-2.0, 4.5) -- (-2.0, 3.6) node[midway, left=2mm, popblue] {\textbf{VCS}};
% IVC (Inferior Vena Cava) -> RA (Bottom Left)
\draw[popblue, very thick] (-2.0, -1.0) -- (-2.0, -0.2) node[midway, left=2mm, popblue] {\textbf{IVC}};
% Pulmonary Artery -> RV (Top Middle-Left)
\draw[popblue, very thick] (0.5, 3.2) -- (1.8, 4.2) node[midway, right=1mm, popblue] {\textbf{Art. Pulmonaire}};
% Pulmonary Veins (x4) -> LA (Top Right) - simplified as two trunks
\draw[poporange, very thick] (2.8, 3.8) -- (3.8, 4.5) node[midway, right=1mm, poporange] {\textbf{Veines Pulm.}};
\draw[poporange, very thick] (2.8, 3.0) -- (3.8, 2.3) node[midway, right=1mm, poporange] {};
% Aorta -> LV (Top Middle-Right, posterior to Pulm Art)
\draw[poporange, very thick] (-0.5, 3.2) -- (-1.8, 4.2) node[midway, left=1mm, poporange] {\textbf{Aorte}};
% Aortic arch hint
\draw[poporange, very thick] (-1.8, 4.2) .. controls (-2.5, 4.5) and (-2.5, 3.5) .. (-1.8, 3.2);

% LABELS CAVITIES
\node[align=center, popblue] at (-1.5, 3.0) {\textbf{Oreillette}\\\textbf{ droite}};
\node[align=center, popblue] at (-1.3, 0.8) {\textbf{Ventricule}\\\textbf{ droit}};
\node[align=center, poporange] at (1.5, 3.0) {\textbf{Oreillette}\\\textbf{ gauche}};
\node[align=center, poporange] at (1.3, 0.8) {\textbf{Ventricule}\\\textbf{ gauche}};

% ARROWS FLOW
% Blue: Deoxygenated
\draw[->, popblue, thick] (-2.0, 4.5) -- (-2.0, 3.6);
\draw[->, popblue, thick] (-2.0, -0.2) -- (-2.0, -1.0);
\draw[->, popblue, thick] (-1.5, 2.5) -- (-0.5, 1.5); % RA -> RV
\draw[->, popblue, thick] (0.5, 2.5) -- (1.5, 3.5); % RV -> Pulm Art

% Red: Oxygenated
\draw[->, poporange, thick] (3.8, 4.5) -- (2.8, 3.8);
\draw[->, poporange, thick] (3.8, 2.3) -- (2.8, 3.0);
\draw[->, poporange, thick] (1.5, 2.5) -- (0.5, 1.5); % LA -> LV
\draw[->, poporange, thick] (-0.5, 2.5) -- (-1.5, 3.5); % LV -> Aorta

% LEGEND BOX
\node[draw=popdark, fill=white, rounded corners, align=left, font=\tiny, anchor=south east] at (3.5, -1.2) {
\textbf{Légende :}\\
\textcolor{popblue}{---} Sang désoxygéné (CO$_2$ riche)\\
\textcolor{poporange}{---} Sang oxygéné (O$_2$ riche)\\
\textcolor{popdark}{- - -} Valves AV (Tricuspide / Mitrale)
};
\end{tikzpicture}
\legende{Coupe antérieure du cœur (vue de face). Côté gauche du dessin = Cœur droit (sang désoxygéné, bleu). Côté droit du dessin = Cœur gauche (sang oxygéné, orange). VCS : Veine Cave Supérieure ; IVC : Veine Cave Inférieure.}
\end{popfigure}

\begin{definition}[Cycle cardiaque]
Le cycle cardiaque comprend une \textbf{systole} (contraction : éjection du sang) et une \textbf{diastole} (relâchement : remplissage). La \textbf{fréquence cardiaque (FC)} est le nombre de cycles par minute (bpm). Au repos : 60--100 bpm. Le \textbf{débit cardiaque (Qc)} = Volume d'éjection systolique (VES) $\times$ FC.
\end{definition}

\begin{propriete}[Double circulation et pressions]
\begin{itemize}
    \item \textbf{Circulation pulmonaire (basse pression) :} VD $\xrightarrow{\text{Art. Pulmonaire}}$ Poumons $\xrightarrow{\text{Veines Pulmonaires}}$ OG. Pression artérielle pulmonaire $\approx$ 25/10 mmHg.
    \item \textbf{Circulation systémique (haute pression) :} VG $\xrightarrow{\text{Aorte}}$ Organes $\xrightarrow{\text{Veines Caves}}$ OD. Pression artérielle systémique (normale) $\approx$ 120/80 mmHg.
\end{itemize}
\end{propriete}

\courssub{La tension artérielle et sa mesure}

\begin{definition}[Pression artérielle]
La \textbf{pression artérielle (PA)} est la pression exercée par le sang sur la paroi des artères. On distingue :
\begin{itemize}
    \item \textbf{Pression Artérielle Systolique (PAS) :} Pression maximale lors de la systole ventriculaire (éjection). Valeur de référence : $\sim$ 120 mmHg.
    \item \textbf{Pression Artérielle Diastolique (PAD) :} Pression minimale lors de la diastole (remplissage, fermeture valves sigmoïdes). Valeur de référence : $\sim$ 80 mmHg.
    \item \textbf{Pression Artérielle Moyenne (PAM) :} $\text{PAM} \approx \text{PAD} + \frac{1}{3}(\text{PAS} - \text{PAD})$. Perfusion des organes.
\end{itemize}
Unité : millimètre de mercure (mmHg) ou kilopascal (kPa) ; 1 kPa $\approx$ 7,5 mmHg.
\end{definition}

\begin{methode}[Mesure de la tension artérielle (Méthode auscultatoire de Korotkoff)]
\begin{enumerate}
    \item \textbf{Conditions :} Sujet au repos depuis 10 min, assis, bras nu au niveau du cœur, pas de café/tabac 30 min avant.
    \item \textbf{Matériel :} Brassard (manchon) adapté à la circonférence du bras, tensiomètre à mercure ou anéroïde, stéthoscope.
    \item \textbf{Poser le brassard :} Bord inférieur 2--3 cm au-dessus du pli du coude. Artère brachiale centrée sous la chambre à air.
    \item \textbf{Gonfler :} Fermer la vis, gonfler jusqu'à disparition des bruits de Korotkoff + 20--30 mmHg.
    \item \textbf{Dégonfler lentement :} 2--3 mmHg/seconde. Écouter au stéthoscope sur l'artère brachiale (pli du coude).
    \item \textbf{Lire :}
    \begin{itemize}
        \item Premier bruit clair (Phase I Korotkoff) = \textbf{PAS}.
        \item Disparition brutale des bruits (Phase V Korotkoff) = \textbf{PAD}.
    \end{itemize}
    \item \textbf{Répéter :} 2 mesures espacées de 1--2 min. Moyenne retenue. Si écart $>$ 5 mmHg, 3\up{ème} mesure.
\end{enumerate}
\end{methode}

\begin{propriete}[Classification de la PA chez l'adulte (OMS / ESH 2023 -- seuils Probatoire)]
\begin{center}
\begin{tabularx}{\linewidth}{|l|X|X|}\hline
\rowcolor{popblueL}
\textbf{Catégorie} & \textbf{PAS (mmHg)} & \textbf{PAD (mmHg)} \\ \hline
Optimale & $<$ 120 & $<$ 80 \\ \hline
Normale & 120--129 & 80--84 \\ \hline
Normale haute (Préhypertension) & 130--139 & 85--89 \\ \hline
HTA Grade 1 (Légère) & 140--159 & 90--99 \\ \hline
HTA Grade 2 (Modérée) & 160--179 & 100--109 \\ \hline
HTA Grade 3 (Sévère) & $\ge$ 180 & $\ge$ 110 \\ \hline
HTA Systolique Isolée & $\ge$ 140 & $<$ 90 \\ \hline
\end{tabularx}
\end{center}
\textbf{Règle diagnostique :} Le diagnostic d'HTA ne peut être posé que sur \textbf{plusieurs mesures effectuées à des occasions différentes} (généralement 3 consultations sur 3--6 mois), hors contexte d'urgence hypertensive.
\end{propriete}

\courssub{Les principaux accidents cardiovasculaires}

\begin{definition}[Infarctus du myocarde (IDM / Crise cardiaque)]
Nécrose (mort) d'une partie du muscle cardiaque (myocarde) par \textbf{ischémie} prolongée (arrêt de l'apport en O$_2$).
\textbf{Cause immédiate :} Occlusion brutale d'une \textbf{artère coronaire} par un \textbf{thrombus} (caillot) se formant sur une \textbf{plaque d'athérome} rompue (athérosclérose).
\textbf{Signes d'alerte (FAIRE) :} Douleur thoracique \textbf{en barreau} (étau, brûlure) $>$ 20 min, irradiant \textbf{bras gauche, mâchoire, dos}, non calmée par le repos/nitroglycérine. Pâleur, sueurs froides, dyspnée, angoisse, nausées. \textbf{Urgence vitale : appeler le 112 (SAMU/Pompiers)}.
\end{definition}

\begin{definition}[Accident Vasculaire Cérébral (AVC / Attaque cérébrale)]
Déficit neurologique brutal par atteinte de la vascularisation cérébrale. Deux types majeurs :
\begin{itemize}
    \item \textbf{AVC Ischémique (80\% des cas) :} Obstruction d'une artère cérébrale par un \textbf{thrombus} (thrombose sur plaque d'athérome) ou un \textbf{embole} (venant du cœur : fibrillation auriculaire, ou des grosses artères). $\to$ Ischémie $\to$ Infarctus cérébral.
    \item \textbf{AVC Hémorragique (20\% des cas) :} Rupture d'une artère cérébrale (souvent sur HTA mal contrôlée, anévrisme, angiome) $\to$ Hématome intra-cérébral $\to$ Compression du tissu cérébral.
\end{itemize}
\textbf{Signes d'alerte (Test \cle{FAST} / \cle{VITE}) :}
\textbf{V} -- Visage déformé (bouche tordue) ; \textbf{I} -- Incapacité à lever un bras ; \textbf{T} -- Trouble de la parole (incompréhensible) ; \textbf{E} -- Extrême urgence : \textbf{Appeler le 112 immédiatement}. \emph{Le temps, c'est du cerveau.}
\end{definition}

\begin{definition}[Hypertension Artérielle (HTA) : Maladie silencieuse]
Élévation persistante des chiffres tensionnels. Souvent asymptomatique (découverte fortuite). \textbf{Complications à long terme si non traitée :}
\begin{itemize}
    \item \textbf{Cardiaques :} Hypertrophie ventriculaire gauche (cœur fatigué) $\to$ Insuffisance cardiaque ; Infarctus du myocarde ; Arythmies (fibrillation auriculaire).
    \item \textbf{Cérébrales :} AVC (ischémique ou hémorragique), Démence vasculaire.
    \item \textbf{Rénales :} Néphro-angiosclérose $\to$ Insuffisance rénale chronique.
    \item \textbf{Vasculaires :} Artériopathie oblitérante des membres inférieurs (AOMI), Anévrisme de l'aorte, Rétinopathie hypertensive.
\end{itemize}
\end{definition}

\courssub{Les facteurs de risque cardiovasculaire}

\begin{aretenir}[Classification des facteurs de risque (OMS / SCORE)]
\begin{center}
\begin{tabularx}{\linewidth}{|X|X|}\hline
\rowcolor{popblueL}
\textbf{Facteurs NON MODIFIABLES (Constitutionnels)} & \textbf{Facteurs MODIFIABLES (Comportementaux / Métaboliques)} \\ \hline
\cle{Âge} : Homme $>$ 50 ans, Femme $>$ 60 ans (ou ménopausée) & \cle{Tabagisme} (actif ou passif) -- \textbf{1er facteur évitable} \\ \hline
\cle{Sexe} : Homme $>$ Femme (avant ménopause) & \cle{Hypertension artérielle} (HTA) \\ \hline
\cle{Antécédents familiaux} précoces (IDM/AVC père $<$ 55 ans, mère $<$ 65 ans) & \cle{Dyslipidémie} : LDL-cholestérol $\uparrow$, HDL-cholestérol $\downarrow$, Triglycérides $\uparrow$ \\ \hline
\cle{Ethnie} (prédisposition génétique) & \cle{Diabète de type 2} (et diabète gestationnel) \\ \hline
& \cle{Obésité abdominale} (Tour de taille H $>$ 94 cm, F $>$ 80 cm) / IMC $\ge$ 25 \\ \hline
& \cle{Sédentarité} ($<$ 150 min/semaine activité modérée) \\ \hline
& \cle{Alimentation} : Sel $\uparrow$ (NaCl), Graisses saturées $\uparrow$, Fibres $\downarrow$, Fruits/Légumes $\downarrow$ \\ \hline
& \cle{Alcoolisation} excessive ($>$ 2 verres/jour H, $>$ 1 F) \\ \hline
& \cle{Stress} psychosocial chronique, manque de sommeil \\ \hline
\end{tabularx}
\end{center}
\textbf{Synergie :} La présence de plusieurs facteurs multiplie le risque (effet exponentiel, non additif). Ex: HTA + Tabagisme + Diabète = Risque $\times$ 10 à $\times$ 20 vs population générale. Scores de risque global (ex: \cle{SCORE2}, \cle{Framingham}) utilisés en pratique.
\end{aretenir}

\courssub{La prévention des accidents cardiovasculaires}

\begin{propriete}[Niveaux de prévention]
\begin{itemize}
    \item \textbf{Primaire (avant la maladie) :} Agir sur les facteurs modifiables chez le sujet sain. \emph{Cible : population générale.}
    \item \textbf{Secondaire (dépistage précoce) :} Dépister et traiter les facteurs de risque (HTA, diabète, dyslipidémie) et maladies asymptomatiques. \emph{Cible : sujets à risque, $\ge$ 40 ans.}
    \item \textbf{Tertiaire (après l'accident) :} Prévenir la récidive, la réadaptation (rééducation cardiaque), observance thérapeutique. \emph{Cible : post-IDM, post-AVC.}
\end{itemize}
\end{propriete}

\begin{methode}[Les 3 piliers de l'hygiène de vie cardiovasculaire (Règles "0-5-30" adaptées)]
\begin{enumerate}
    \item \textbf{Alimentation protectrice (Type Méditerranéen / DASH) :}
    \begin{itemize}
        \item \textbf{Réduire le sel :} $<$ 5 g/jour (1 c. à café) = 2 g Sodium. Attention sel caché (pain, conserves, bouillons cubes, sauces soja, plats rapides).
        \item \textbf{Limiter graisses saturées} (viandes grasses, charcuteries, beurre, huile de palme, fritures) ; \textbf{Privilégier} oméga-3 (poissons gras, noix, huile colza/lin), fibres (céréales complètes, légumineuses, fruits/légumes 5 portions/j).
        \item \textbf{Éviter} sucres ajoutés, boissons sucrées, alcool excessif.
    \end{itemize}
    \item \textbf{Activité physique régulière :}
    \begin{itemize}
        \item $\ge$ 150 min/semaine d'endurance modérée (marche rapide, vélo, natation, danse) OU 75 min intense.
        \item Renforcement musculaire 2$\times$/sem. Limiter la sédentarité (lever toutes les heures).
    \end{itemize}
    \item \textbf{Éviter les toxiques \& Dépistage :}
    \begin{itemize}
        \item \textbf{Zéro tabac} (sevrage = bénéfice immédiat sur risque IDM/AVC).
        \item \textbf{Alcool} : modération stricte.
        \item \textbf{Dépistage systématique :} PA (annuel $\ge$ 18 ans), Glycémie à jeun / HbA1c, Bilan lipidique (LDL/HDL/Triglycérides) selon facteurs de risque. IMC + Tour de taille.
    \end{itemize}
\end{enumerate}
\end{methode}

\begin{savaistu}[Contexte Camerounais : La transition épidémiologique]
Au Cameroun, les maladies cardiovasculaires (MCV) sont la \textbf{1ère cause de mortalité} dans les hôpitaux (registres de Yaoundé, Douala). L'HTA touche $\sim$ 30\% des adultes (souvent méconnue, mal contrôlée). L'urbanisation rapide (Douala, Yaoundé, Bafoussam) favorise la sédentarité, l'alimentation industrielle (nouilles instantanées, conserves, boissons sucrées), le stress et la pollution. Le \textbf{Programme National de Lutte contre les Maladies Non Transmissibles (PNLMNT)} promeut le dépistage intégré (HTA, Diabète, Obésité) dans les formations sanitaires de base (CMA, CS). Le numéro d'urgence unique est le \textbf{112} (gratuit, GSM).
\end{savaistu}

\courssub{Synthèse et évaluation des acquis}

\courssub{Je vérifie mes connaissances}

\begin{exercice}[QCM : Système cardiovasculaire et tension artérielle]
Parmi les affirmations suivantes, une seule est correcte. Laquelle ?
\begin{enumerate}
    \item Le sang artériel est toujours oxygéné et le sang veineux toujours désoxygéné.
    \item La pression artérielle systolique correspond à la pression dans les artères lors du relâchement ventriculaire (diastole).
    \item L'hypertension artérielle (HTA) est définie par une pression artérielle systolique $\ge 140$ mmHg et/ou une pression diastolique $\ge 90$ mmHg, confirmée à plusieurs reprises.
    \item L'infarctus du myocarde résulte systématiquement de la rupture d'une artère coronaire.
\end{enumerate}
\end{exercice}

\begin{exercice}[Vrai ou Faux : Facteurs de risque et accidents vasculaires]
Pour chaque affirmation, indiquez si elle est VRAIE ou FAUSSE et justifiez brièvement votre réponse.
\begin{enumerate}
    \item L'âge et le sexe sont des facteurs de risque cardiovasculaire modifiables.
    \item Un Accident Vasculaire Cérébral (AVC) ischémique est causé par l'obstruction d'une artère cérébrale par un caillot (thrombus ou embolie).
    \item La fréquence cardiaque au repos d'un adulte en bonne santé se situe généralement entre 100 et 120 battements par minute.
    \item La sédentarité, le tabagisme et une alimentation riche en sel et en graisses saturées sont des facteurs de risque majeurs et modifiables.
\end{enumerate}
\end{exercice}

\begin{exercice}[Définitions et valeurs de référence]
Donnez la définition précise des termes suivants et indiquez les valeurs de référence normales pour un adulte au repos :
\begin{enumerate}
    \item Pression artérielle systolique (PAS) et diastolique (PAD).
    \item Fréquence cardiaque (FC).
    \item Infarctus du myocarde.
    \item Accident Vasculaire Cérébral (AVC) hémorragique.
\end{enumerate}
\end{exercice}

\begin{exercice}[Calcul direct : Fréquence cardiaque]
Un élève de Première compte ses pulsations au niveau de l'artère radiale.
\begin{enumerate}
    \item Au repos, il compte 18 battements en 15 secondes. Calculez sa fréquence cardiaque en battements par minute (bpm). Cette valeur est-elle normale ?
    \item Juste après un effort modéré (monter 3 étages), il compte 35 battements en 15 secondes. Calculez sa fréquence cardiaque. Comment évolue-t-elle par rapport au repos ? Expliquez le mécanisme physiologique.
\end{enumerate}
\end{exercice}

\courssub{Je m'entraîne}

\begin{exercice}[Interprétation de mesures tensionnelles]
Trois personnes se présentent à une campagne de dépistage de l'hypertension au Centre de Santé de District. Les mesures moyennes (sur trois prises espacées de quelques minutes) donnent les résultats suivants :
\begin{center}
\begin{tabular}{|c|c|c|}\hline
\textbf{Personne} & \textbf{PAS (mmHg)} & \textbf{PAD (mmHg)} \\ \hline
M. Atangana (52 ans) & 138 & 88 \\ \hline
Mme Mballa (45 ans) & 162 & 98 \\ \hline
M. Essomba (30 ans) & 118 & 75 \\ \hline
\end{tabular}
\end{center}
\begin{enumerate}
    \item Classez chaque mesure selon la classification de l'OMS (Optimale, Normale, Normale haute, HTA Grade 1, Grade 2, Grade 3).
    \item Identifiez la ou les personnes hypertendues (selon les seuils diagnostiques).
    \item Pour chaque personne, proposez une conduite à tenir adaptée (surveillance, hygiène de vie, consultation médicale programmée, consultation urgente, traitement).
\end{enumerate}
\end{exercice}

\begin{exercice}[Schéma du cœur : Anatomie et circulation]
Le schéma ci-dessous représente une coupe antérieure du cœur humain (vue de face), non légendé.
\begin{popfigure}
\begin{tikzpicture}[scale=0.95, every node/.style={font=\small}]
% OUTLINE
\draw[popdark, thick, fill=popcream] (-3, -1) .. controls (-3.5, 1) and (-3.5, 3) .. (-1.5, 3.8) .. controls (0, 4.2) .. (1.5, 3.8) .. controls (3.5, 3) and (3.5, 1) .. (3, -1) .. controls (3, -1.5) and (-3, -1.5) .. (-3, -1);
% SEPTUM
\draw[popdark, thick] (0, -1) -- (0, 3.5);
% CAVITIES OUTLINES (thin)
\draw[popdark, thin] (-2.8, 2.5) .. controls (-2.8, 3.5) and (-1.2, 3.5) .. (0, 3.5) -- (0, 2.0) .. controls (-1.2, 2.0) and (-2.8, 2.0) .. (-2.8, 2.5); % RA
\draw[popdark, thin] (-2.5, -0.5) .. controls (-2.8, 1) and (-1.5, 2.0) .. (0, 2.0) -- (0, -0.5) -- (-2.5, -0.5); % RV
\draw[popdark, thin] (0, 3.5) .. controls (1.2, 3.5) and (2.8, 3.5) .. (2.8, 2.5) .. controls (2.8, 2.0) and (1.2, 2.0) .. (0, 2.0) -- (0, 3.5); % LA
\draw[popdark, thin] (0, -0.5) -- (0, 2.0) .. controls (1.5, 2.0) and (2.8, 1) .. (2.5, -0.5) -- (0, -0.5); % LV

% VESSELS (Numbered 1 to 5)
% 1: VCS (Top Left - Blue)
\draw[popblue, very thick] (-2.0, 4.5) -- (-2.0, 3.6) node[midway, left=2mm, popblue] {\textbf{1}};
% 2: IVC (Bottom Left - Blue)
\draw[popblue, very thick] (-2.0, -1.0) -- (-2.0, -0.2) node[midway, left=2mm, popblue] {\textbf{2}};
% 3: Pulmonary Artery (Top Middle-Left - Blue)
\draw[popblue, very thick] (0.5, 3.2) -- (1.8, 4.2) node[midway, right=1mm, popblue] {\textbf{3}};
% 4: Pulmonary Veins (Top Right - Red/Orange)
\draw[poporange, very thick] (2.8, 3.8) -- (3.8, 4.5) node[midway, right=1mm, poporange] {\textbf{4}};
\draw[poporange, very thick] (2.8, 3.0) -- (3.8, 2.3);
% 5: Aorta (Top Middle-Right - Red/Orange)
\draw[poporange, very thick] (-0.5, 3.2) -- (-1.8, 4.2) node[midway, left=1mm, poporange] {\textbf{5}};
\draw[poporange, very thick] (-1.8, 4.2) .. controls (-2.5, 4.5) and (-2.5, 3.5) .. (-1.8, 3.2);

% CAVITY LABELS (A, B, C, D)
\node[align=center] at (-1.5, 3.0) {\textbf{A}};
\node[align=center] at (1.5, 3.0) {\textbf{B}};
\node[align=center] at (-1.3, 0.8) {\textbf{C}};
\node[align=center] at (1.3, 0.8) {\textbf{D}};
\end{tikzpicture}
\legende{Coupe antérieure schématique du cœur (vue de face). Les numéros 1 à 5 indiquent les grands vaisseaux. Les lettres A, B, C, D indiquent les cavités.}
\end{popfigure}
\begin{enumerate}
    \item Nommez les quatre cavités cardiaques (A, B, C, D) et les cinq vaisseaux (1 à 5) en précisant pour chaque vaisseau s'il transporte du sang oxygéné ou désoxygéné.
    \item Tracez, sur une copie du schéma, le trajet du sang désoxygéné (flèches bleues) et du sang oxygéné (flèches rouges) à l'intérieur du cœur.
    \item Rappellez le rôle des valves situées entre les oreillettes et les ventricules (valves atrio-ventriculaires).
\end{enumerate}
\end{exercice}

\begin{exercice}[Étude de cas : Conduite à tenir face à un infarctus]
M. Kouam, 58 ans, chef d'entreprise à Douala, ressent brutalement au cours d'un repas d'affaires une douleur violente en barreau au niveau du thorax, irradiant vers l'épaule et le bras gauche, ainsi que la mâchoire. Il est pâle, transpire abondamment (sueurs froides), a du mal à respirer et ressent une angoisse intense. Il nie la douleur par peur de l'hôpital.
\begin{enumerate}
    \item Nommez l'accident cardiovasculaire suspecté. Justifiez votre réponse par au moins trois signes cliniques présents dans le texte.
    \item Quelle est la cause immédiate de cet accident au niveau des artères coronaires ?
    \item En attendant l'arrivée des secours (SAMU / Pompiers : \textbf{112}), détaillez les 4 gestes de premiers secours à effectuer dans l'ordre chronologique.
    \item Pourquoi ne doit-on \textbf{pas} donner à boire ni à manger à la victime ?
\end{enumerate}
\end{exercice}

\begin{exercice}[Classement des facteurs de risque]
Voici une liste de facteurs de risque cardiovasculaire : \emph{Âge $> 50$ ans, Tabagisme, Diabète de type 2, Sexe masculin, Obésité abdominale, Antécédents familiaux précoces, Hypertension artérielle, Sédentarité, Dyslipidémie (excès LDL-cholestérol), Stress chronique, Alcoolisation excessive, Femme ménopausée.}
\begin{enumerate}
    \item Réalisez un tableau à deux colonnes pour classer ces facteurs en \textbf{Modifiables} (liés au mode de vie ou traitables) et \textbf{Non modifiables} (constitutionnels).
    \item Parmi les facteurs modifiables, lesquels relèvent d'une prise en charge médicale (thérapeutique) et lesquels relèvent principalement d'une modification de l'hygiène de vie (comportementale) ? (Certains peuvent relever des deux).
    \item Expliquez pourquoi la combinaison de plusieurs facteurs (ex: HTA + Tabagisme + Diabète) augmente le risque global de façon exponentielle (multiplicative) et non additive.
\end{enumerate}
\end{exercice}

\courssub{Je me prépare au Probatoire}

\begin{exercice}[Sujet type Probatoire : Dépistage et prise en charge de l'HTA]
Lors d'une journée « Cœur en santé » organisée au Lycée de Biyem-Assi, l'infirmière scolaire mesure la tension artérielle de trois élèves de Première C au repos, après 10 minutes assises. Les résultats (moyenne de 2 mesures) sont :
\begin{center}
\begin{tabular}{|c|c|c|c|}\hline
\textbf{Élève} & \textbf{Âge} & \textbf{PAS (mmHg)} & \textbf{PAD (mmHg)} \\ \hline
A (Garçon) & 16 ans & 125 & 78 \\ \hline
B (Fille) & 17 ans & 138 & 88 \\ \hline
C (Garçon) & 16 ans & 145 & 92 \\ \hline
\end{tabular}
\end{center}
\textbf{Rappel :} Chez l'adolescent $> 13$ ans, les seuils de l'HTA rejoignent ceux de l'adulte (HTA si PAS $\ge 140$ et/ou PAD $\ge 90$). Une « normale haute » (préhypertension) est définie par PAS 130-139 et/ou PAD 85-89.
\begin{enumerate}
    \item Interprétez la tension artérielle de chaque élève en utilisant les termes précis : « Normale », « Normale haute (Préhypertension) », « Hypertension artérielle (HTA) ».
    \item L'élève C présente une HTA. L'infirmière doit-elle poser le diagnostic définitif d'HTA sur cette unique mesure ? Justifiez en citant la règle diagnostique standard.
    \item L'élève C est en surpoids (IMC = 28 kg/m$^2$), ne pratique pas de sport et consomme beaucoup de plats rapides riches en sel (nouilles instantanées, conserves). Citez \textbf{trois} mesures hygiéno-diététiques prioritaires à lui conseiller pour faire baisser sa tension sans médicament dans un premier temps.
    \item L'élève B a une tension « normale haute ». Expliquez pourquoi elle constitue une population à risque et quelle surveillance lui proposer.
    \item Rappellez les deux complications majeures, l'une cardiaque, l'autre cérébrale, d'une HTA non traitée au long cours.
\end{enumerate}
\end{exercice}

\begin{exercice}[Sujet type Probatoire : Prévention et éducation à la santé]
Vous êtes délégué santé de votre classe. Le Proviseur vous demande de rédiger un message de sensibilisation (affichage ou réseau social de l'établissement) destiné aux parents d'élèves et aux personnels, sur le thème : « \emph{Prévenir les accidents cardiovasculaires : 3 gestes quotidiens qui sauvent des vies} ».
\begin{enumerate}
    \item Rédigez ce message (entre 8 et 10 lignes maximum). Il doit être clair, percutant, scientifiquement exact et adapter le vocabulaire au grand public.
    \item Votre message doit obligatoirement citer \textbf{trois mesures de prévention distinctes} (une alimentaire, une liée à l'activité physique, une liée au dépistage ou aux toxiques) et expliquer \textbf{brièvement le mécanisme} par lequel chacune protège le système cardiovasculaire.
    \item Indiquez le numéro d'urgence à composer au Cameroun en cas de signes d'infarctus ou d'AVC (douleur thoracique, paralysie brutale, trouble de la parole).
\end{enumerate}
\end{exercice}

\begin{exercice}[Sujet type Probatoire : Physiopathologie de l'AVC]
Le document ci-dessous présente deux coupes schématiques d'une artère cérébrale.
\begin{popfigure}
\begin{tikzpicture}[scale=1.3, every node/.style={font=\small}]
% Coupe 1 : Normale (Left)
\begin{scope}[xshift=-5.5cm]
% Adventitia / Media
\draw[popdark, thick, fill=popmuted!20] (-1.5, -1.2) rectangle (1.5, 1.2);
% Intima / Lumen
\draw[poporange, thick, fill=poporange!30] (-1.2, -0.9) rectangle (1.2, 0.9);
% Endothelium line
\draw[popdark, very thick] (-1.2, 0.9) -- (1.2, 0.9);
\draw[popdark, very thick] (-1.2, -0.9) -- (1.2, -0.9);
\node at (0, -1.6) {\textbf{Artère saine}};
\node[poporange, font=\bfseries] at (0, 0) {Lumière (Sang)};
\node[popdark, align=center, font=\tiny] at (0, 1.5) {Intima intacte\\ Média élastique};
\end{scope}

% Coupe 2 : Athéromateuse + Thrombus (Right)
\begin{scope}[xshift=5.5cm]
% Outer wall
\draw[popdark, thick, fill=popmuted!20] (-1.5, -1.2) rectangle (1.5, 1.2);
% Atheroma Plaque (Pink) - protruding into lumen
\fill[poppink!70] (-1.2, -0.1) .. controls (-0.6, 0.4) and (0.6, 0.4) .. (1.2, -0.1) -- (1.2, -0.9) -- (-1.2, -0.9) -- cycle;
\draw[popdark, thick] (-1.2, -0.1) .. controls (-0.6, 0.4) and (0.6, 0.4) .. (1.2, -0.1);
% Thrombus (Red) - on top of plaque
\fill[popred!80] (-0.4, -0.1) .. controls (-0.1, 0.2) and (0.5, 0.2) .. (0.9, -0.1) -- (0.9, -0.25) .. controls (0.5, 0.05) and (-0.1, 0.05) .. (-0.4, -0.25) -- cycle;
\draw[popred, thick] (-0.4, -0.1) .. controls (-0.1, 0.2) and (0.5, 0.2) .. (0.9, -0.1);
% Residual Lumen (Orange) - only lower part
\fill[poporange!30] (-1.2, -0.9) rectangle (1.2, -0.35);
% Endothelium remnants
\draw[popdark, very thick] (-1.2, -0.9) -- (1.2, -0.9);

\node at (0, -1.6) {\textbf{Artère athéromateuse}};
\node[poppink, font=\bfseries, align=center] at (0, 0.4) {Plaque\\ d'athérome};
\node[popred, font=\bfseries] at (0.3, -0.15) {Thrombus};
\node[poporange, font=\bfseries, align=center] at (0, -0.6){Lumière\\ rétrécie};
\end{scope}
\end{tikzpicture}
\legende{Coupes transversales schématiques d'une artère cérébrale : saine (gauche) et atteinte d'athérosclérose compliquée d'une thrombose sur plaque rompue (droite).}
\end{popfigure}
\begin{enumerate}
    \item Sur le document de droite, nommez les structures indiquées par les zones colorées : \textbf{zone rose} et \textbf{zone rouge}.
    \item Expliquez le mécanisme de formation de la plaque d'athérome (zone rose) en reliant les étapes clés : excès de LDL-cholestérol, endothelium, monocytes/macrophages, cellules spumeuses, cellules musculaires lisses, cap fibreux.
    \item Quelle est la conséquence hémodynamique immédiate de la présence du thrombus (zone rouge) sur le débit sanguin cérébral en aval de l'obstacle ?
    \item Différenciez l'AVC \textbf{ischémique} (infarctus cérébral) et l'AVC \textbf{hémorragique} par leur mécanisme lésionnel initial et leur prise en charge thérapeutique d'urgence concernant la thrombolyse (OUI / NON ?).
    \item Citez le test mnémotechnique international (acronyme anglais) utilisé pour reconnaître rapidement les signes d'alerte d'un AVC auprès du grand public et donnez la signification de chaque lettre en français.
\end{enumerate}
\end{exercice}

\newpage

\coursec{Corrigés des exercices}

\begin{corrige}[Exercice 1 — QCM : Système cardiovasculaire et tension artérielle]
\emph{Analyse des propositions :}
\begin{enumerate}
    \item \textbf{Fausse.} Le sang artériel est défini par le sens du flux (sortie du cœur), le sang veineux par le sens du flux (retour au cœur). Dans la \textbf{circulation pulmonaire}, l'artère pulmonaire transporte du sang \textbf{désoxygéné} (bleu) et les veines pulmonaires transportent du sang \textbf{oxygéné} (rouge). Seule la circulation systémique respecte la règle « artère = oxygéné, veine = désoxygéné ».
    \item \textbf{Fausse.} La \textbf{Pression Artérielle Systolique (PAS)} correspond à la pression maximale lors de la \textbf{systole ventriculaire} (contraction/éjection). La diastole correspond à la PAD (pression minimale).
    \item \textbf{Vraie.} C'est la définition diagnostique de référence (OMS/ESH) : HTA si $\text{PAS} \ge 140$ mmHg \textbf{et/ou} $\text{PAD} \ge 90$ mmHg, \textbf{confirmée à plusieurs reprises} (généralement 3 consultations sur 3--6 mois) hors urgence hypertensive.
    \item \textbf{Fausse.} L'infarctus du myocarde résulte de l'\textbf{occlusion} (obstruction) d'une artère coronaire par un thrombus sur plaque d'athérome rompue. La \textbf{rupture de l'artère} elle-même est une complication rare (rupture du myocarde, rupture de septum, rupture de muscle papillaire), pas la cause initiale de l'infarctus.
\end{enumerate}
\textbf{Bonne réponse : Proposition 3.}
\par\textbf{Attention [Piège fréquent] :} Ne pas confondre « sang artériel / veineux » (topographie : sens du flux) et « sang oxygéné / désoxygéné » (composition en gaz). La circulation pulmonaire inverse la correspondance habituelle.
\end{corrige}

\begin{corrige}[Exercice 2 — Vrai ou Faux : Facteurs de risque et accidents vasculaires]
\begin{enumerate}
    \item \textbf{FAUX.} L'âge et le sexe sont des facteurs de risque \textbf{non modifiables} (constitutionnels). On ne peut pas rajeunir ni changer son sexe génétique. Seuls les facteurs liés au mode de vie (tabagisme, alimentation, activité physique, alcool) ou métaboliques traitables (HTA, diabète, dyslipidémie) sont modifiables.
    \item \textbf{VRAI.} L'AVC ischémique (80\% des cas) est causé par l'obstruction brutale d'une artère cérébrale par un \textbf{thrombus} (thrombose sur plaque d'athérome \emph{in situ}) ou un \textbf{embole} (embolie d'origine cardiaque -- fibrillation auriculaire -- ou artérielle). Cela entraîne une ischémie puis un infarctus cérébral en aval.
    \item \textbf{FAUX.} La fréquence cardiaque au repos normale chez l'adulte est comprise entre \textbf{60 et 100 bpm}. Une FC au repos entre 100 et 120 bpm correspond à une \textbf{tachycardie} (pathologique ou réactionnelle : fièvre, anxiété, thyrotoxicose, insuffisance cardiaque, etc.).
    \item \textbf{VRAI.} Ce sont les trois piliers majeurs des facteurs de risque \textbf{modifiables} (comportementaux et métaboliques). La sédentarité, le tabagisme (1er facteur évitable) et une alimentation déséquilibrée (excès de sel $\to$ HTA, excès de graisses saturées $\to$ dyslipidémie/athérosclérose) potentialisent leurs effets (synergie multiplicative).
\end{enumerate}
\end{corrige}

\begin{corrige}[Exercice 3 — Définitions et valeurs de référence]
\begin{enumerate}
    \item \textbf{Pression Artérielle Systolique (PAS) :} Pression maximale exercée par le sang sur la paroi des artères lors de la \textbf{systole ventriculaire} (éjection du sang par le ventricule gauche dans l'aorte). \textbf{Valeur de référence :} $\sim 120$ mmHg (Optimale $< 120$, Normale 120--129).
    \textbf{Pression Artérielle Diastolique (PAD) :} Pression minimale exercée par le sang sur la paroi des artères lors de la \textbf{diastole ventriculaire} (remplissage, valves sigmoïdes fermées, maintien de la pression par l'élasticité artérielle). \textbf{Valeur de référence :} $\sim 80$ mmHg (Optimale $< 80$, Normale 80--84).
    \item \textbf{Fréquence Cardiaque (FC) :} Nombre de cycles cardiaques (systole + diastole) par minute. \textbf{Valeur de référence au repos :} \textbf{60 à 100 bpm} (souvent 60--80 bpm chez l'adulte entraîné).
    \item \textbf{Infarctus du myocarde (IDM) :} Nécrose (mort cellulaire irréversible) d'une portion du muscle cardiaque (myocarde) secondaire à une \textbf{ischémie prolongée} (arrêt de l'apport sanguin) due à l'occlusion brutale d'une artère coronaire par un thrombus sur plaque d'athérome.
    \item \textbf{AVC Hémorragique :} Type d'accident vasculaire cérébral (20\% des cas) causé par la \textbf{rupture} spontanée d'une artère intracérébrale (souvent sur HTA mal contrôlée, anévrisme, angiome), entraînant un \textbf{hématome intra-parenchymateux} qui compresse et détruit le tissu cérébral environnant.
\end{enumerate}
\end{corrige}

\begin{corrige}[Exercice 4 — Calcul direct : Fréquence cardiaque]
\begin{enumerate}
    \item \textbf{Au repos :} 18 battements en 15 secondes.
    \[
    \text{FC} = 18 \times \frac{60}{15} = 18 \times 4 = \textbf{72 bpm}.
    \]
    \textbf{Interprétation :} 72 bpm se situe dans l'intervalle de référence [60--100 bpm]. \textbf{C'est une valeur normale.}
    \item \textbf{Après effort :} 35 battements en 15 secondes.
    \[
    \text{FC} = 35 \times 4 = \textbf{140 bpm}.
    \]
    \textbf{Évolution :} La FC augmente de 72 à 140 bpm (quasi doublement).
    \textbf{Mécanisme physiologique :} Lors de l'effort, les muscles squelettiques ont des besoins accrus en $\text{O}_2$ et en nutriments, et produisent plus de $\text{CO}_2$ et $\text{H}^+$.
    \begin{itemize}
        \item \textbf{Système nerveux autonome :} Diminution du tonus \textbf{parasympathique} (vagal) + augmentation du tonus \textbf{sympathique} (noradrénaline/adrénaline) $\to$ accélération du nœud sinusal (chronotropie positive).
        \item \textbf{Rétrocontrôle :} Chémorécepteurs ($\uparrow \text{CO}_2, \downarrow \text{pH}, \downarrow \text{O}_2$) et barorécepteurs (baisse relative de PA diastolique) stimulent le centre cardio-accélérateur bulbaire.
        \item \textbf{Adaptation :} L'augmentation du débit cardiaque ($Q_c = \text{VES} \times \text{FC}$) assure l'hyperémie musculaire active.
    \end{itemize}
\end{enumerate}
\end{corrige}

\begin{corrige}[Exercice 5 — Interprétation de mesures tensionnelles]
\begin{center}
\begin{tabularx}{\linewidth}{|l|c|X|X|}\hline
\rowcolor{popblueL}
\textbf{Personne} & \textbf{PAS / PAD} & \textbf{Classification OMS/ESH} & \textbf{Conduite à tenir} \\ \hline
M. Atangana (52 ans) & 138 / 88 & \textbf{Normale haute (Préhypertension)} & \textbf{Surveillance rapprochée + Hygiène de vie stricte.} Conseil : réduire sel ($< 5$ g/j), activité physique régulière (150 min/sem), perte de poids si IMC $\ge 25$, arrêt tabac si fumeur, limiter alcool. Contrôle TA dans 3--6 mois. Dépistage diabète/dyslipidémie. \\ \hline
Mme Mballa (45 ans) & 162 / 98 & \textbf{HTA Grade 2 (Modérée)} & \textbf{Consultation médicale programmée rapide (semaine).} Confirmation diagnostic (3 mesures sur 3 consultations). Bilan étiologique (rein, endocrinien) et bilan lésionnel (ECG, fond d'œil, créatininémie, protéinurie, bilan lipidique, glycémie). \textbf{Traitement médicamenteux probable d'emblée} (Grade 2 + âge) associé aux mesures hygiéno-diététiques. \\ \hline
M. Essomba (30 ans) & 118 / 75 & \textbf{Optimale} & \textbf{Rassurer. Maintenir hygiène de vie.} Contrôle TA annuel. Promouvoir activité physique, alimentation équilibrée, non-tabagisme. \\ \hline
\end{tabularx}
\end{center}
\begin{enumerate}
    \item Voir tableau ci-dessus (colonne Classification).
    \item \textbf{Personne hypertendue :} \textbf{Mme Mballa} (PAS $\ge 160$ et/ou PAD $\ge 100$ $\to$ Grade 2). M. Atangana est en « normale haute » (préhypertension), pas encore HTA.
    \item Voir tableau ci-dessus (colonne Conduite à tenir).
\end{enumerate}
\par\textbf{Attention [Règle diagnostique] :} Le diagnostic d'HTA \textbf{ne se pose jamais sur une seule séance de mesures} (même avec 3 prises). Il nécessite \textbf{plusieurs consultations} (généralement 3 sur 3 à 6 mois) pour exclure l'effet « blouse blanche » et confirmer le caractère permanent. Seule l'urgence hypertensive (PAS $\ge 180$ et/ou PAD $\ge 110$ avec signes de lésion d'organe aiguë) justifie une prise en charge immédiate.
\end{corrige}

\begin{corrige}[Exercice 6 — Schéma du cœur : Anatomie et circulation]
\begin{enumerate}
    \item \textbf{Cavités :}
    \begin{itemize}
        \item \textbf{A : Oreillette droite} (reçoit VCS et IVC, sang désoxygéné).
        \item \textbf{B : Oreillette gauche} (reçoit veines pulmonaires, sang oxygéné).
        \item \textbf{C : Ventricule droit} (envoie sang vers artère pulmonaire, paroi fine).
        \item \textbf{D : Ventricule gauche} (envoie sang vers aorte, paroi très épaisse).
    \end{itemize}
    \textbf{Vaisseaux :}
    \begin{itemize}
        \item \textbf{1 : Veine Cave Supérieure (VCS)} -- Sang \textbf{désoxygéné} (retour circulation systémique supérieure).
        \item \textbf{2 : Veine Cave Inférieure (IVC)} -- Sang \textbf{désoxygéné} (retour circulation systémique inférieure).
        \item \textbf{3 : Artère Pulmonaire} -- Sang \textbf{désoxygéné} (vers poumons pour oxygénation). \emph{Seule artère transportant du sang désoxygéné.}
        \item \textbf{4 : Veines Pulmonaires} (4 troncs) -- Sang \textbf{oxygéné} (retour poumons $\to$ oreillette gauche). \emph{Seules veines transportant du sang oxygéné.}
        \item \textbf{5 : Aorte} -- Sang \textbf{oxygéné} (vers circulation systémique).
    \end{itemize}
    \item \textbf{Trajet (à tracer sur la copie) :}
    \begin{itemize}
        \item \textbf{Sang désoxygéné (Flèches BLEUES) :} VCS (1) / IVC (2) $\to$ Oreillette droite (A) $\xrightarrow{\text{Valve tricuspide}}$ Ventricule droit (C) $\xrightarrow{\text{Valve sigmoïde pulmonaire}}$ Artère Pulmonaire (3) $\to$ Poumons.
        \item \textbf{Sang oxygéné (Flèches ROUGES) :} Veines Pulmonaires (4) $\to$ Oreillette gauche (B) $\xrightarrow{\text{Valve mitrale}}$ Ventricule gauche (D) $\xrightarrow{\text{Valve sigmoïde aortique}}$ Aorte (5) $\to$ Organes.
    \end{itemize}
    \item \textbf{Rôle des valves atrio-ventriculaires (Tricuspide à droite, Mitrale à gauche) :} Elles assurent l'\textbf{unidirectionnalité} du flux sanguin de l'oreillette vers le ventricule. Elles s'\textbf{ouvrent} pendant la diastole ventriculaire (remplissage) et se \textbf{ferment} hermétiquement au début de la systole ventriculaire (grâce aux piliers et cordages tendineux qui empêchent leur retournement) pour empêcher le reflux du sang vers les oreillettes (régurgitation).
\end{enumerate}
\end{corrige}

\begin{corrige}[Exercice 7 — Étude de cas : Conduite à tenir face à un infarctus]
\begin{enumerate}
    \item \textbf{Accident suspecté :} \textbf{Infarctus du myocarde (IDM) / Crise cardiaque.}
    \textbf{Justification (3 signes cliniques) :}
    \begin{enumerate}
        \item Douleur thoracique \textbf{violente en barreau} (étau, brûlure) survenant \textbf{brutalement} au repos (repas).
        \item \textbf{Irradiation} typique : épaule/bras gauche + mâchoire.
        \item Signes d'instabilité hémodynamique et neurovégétative : \textbf{pâleur, sueurs froides abondantes, dyspnée (difficulté respiratoire), angoisse intense} (sentiment de mort imminente).
    \end{enumerate}
    \item \textbf{Cause immédiate :} \textbf{Occlusion brutale (thrombose) d'une artère coronaire} par un \textbf{thrombus} (caillot sanguin) se formant sur une \textbf{plaque d'athérome rompue} (athérosclérose coronarienne). Cela entraîne une ischémie aiguë puis nécrose du territoire vasculaire en aval.
    \item \textbf{Gestes de premiers secours (ordre chronologique) :}
    \begin{enumerate}
        \item \textbf{Appeler / Faire appeler le 112 (SAMU / Pompiers)} immédiatement. Préciser : « Douleur thoracique type infarctus, homme 58 ans, conscient, instable ». \emph{Ne pas raccrocher avant l'opérateur.}
        \item \textbf{Mettre au repos absolu :} Position \textbf{demi-assise} (tronc à 30--45°, jambes pliées) pour diminuer le retour veineux (précharge) et le travail cardiaque, faciliter la respiration. Interdire tout effort (marcher, monter escaliers).
        \item \textbf{Surveiller en continu :} Conscience, respiration, pouls (radial/carotidien). Rassurer la victime (calme = moins d'adrénaline = moins de travail cardiaque). Desserer vêtements serrants (cravate, ceinture, col).
        \item \textbf{Préparer l'arrivée des secours :} Dégager l'accès, désigner quelqu'un pour guider l'ambulance. Si arrêt cardiaque survient (inconscience, absence respiration/pouls) : \textbf{Démarrer RCP (massage cardiaque 100--120/min) + DAE si disponible} jusqu'à l'arrivée du SAMU.
    \end{enumerate}
    \item \textbf{Pourquoi ne rien donner à boire/manger :}
    \begin{itemize}
        \item Risque majeur de \textbf{fausse route} (inhalation broncho-pulmonaire) si la victime vomit (fréquent par stimulation vagale/douleur) ou perd connaissance (perte des réflexes de déglutition/toux).
        \item Nécessité d'un \textbf{estomac vide} en vue d'une éventuelle anesthésie générale urgente (coronarographie/angioplastie primaire) ou intubation.
        \item La nitroglycérine (souvent donnée par le SAMU) est \textbf{sublinguale} (spray/comprimé), pas orale.
    \end{itemize}
\end{enumerate}
\end{corrige}

\begin{corrige}[Exercice 8 — Classement des facteurs de risque]
\begin{enumerate}
    \item \textbf{Tableau de classification :}
    \begin{center}
    \begin{tabularx}{\linewidth}{|X|X|}\hline
    \rowcolor{popblueL}
    \textbf{Facteurs NON MODIFIABLES (Constitutionnels)} & \textbf{Facteurs MODIFIABLES (Comportementaux / Métaboliques)} \\ \hline
    \begin{itemize}
        \item Âge $> 50$ ans (H) / $> 60$ ans (F)
        \item Sexe masculin
        \item Antécédents familiaux précoces
        \item Femme ménopausée (perte protection œstrogènes)
        \item Ethnie (prédisposition génétique)
    \end{itemize} &
    \begin{itemize}
        \item \textbf{Tabagisme} (actif/passif)
        \item \textbf{Hypertension artérielle}
        \item \textbf{Diabète de type 2}
        \item \textbf{Obésité abdominale} (Tour taille H $> 94$, F $> 80$)
        \item \textbf{Sédentarité}
        \item \textbf{Dyslipidémie} (LDL $\uparrow$, HDL $\downarrow$, TG $\uparrow$)
        \item \textbf{Alcoolisation excessive}
        \item \textbf{Stress chronique} / manque sommeil
        \item \textbf{Alimentation} déséquilibrée (Sel $\uparrow$, Graisses sat. $\uparrow$, Fibres $\downarrow$)
    \end{itemize} \\ \hline
    \end{tabularx}
    \end{center}
    \item \textbf{Prise en charge médicale (Thérapeutique) vs Hygiène de vie (Comportementale) :}
    \begin{center}
    \begin{tabularx}{\linewidth}{|X|X|X|}\hline
    \rowcolor{popblueL}
    \textbf{Facteur Modifiable} & \textbf{Principalement Thérapeutique (Médical)} & \textbf{Principalement Comportementale (Hygiène de vie)} \\ \hline
    HTA & \checkmark (Antihypertenseurs) & \checkmark (Sel, poids, sport, alcool, stress) \\ \hline
    Diabète T2 & \checkmark (Antidiabétiques oraux/insuline) & \checkmark (Régime, sport, poids) \\ \hline
    Dyslipidémie & \checkmark (Statines, fibrates, ézétimibe) & \checkmark (Graisses, fibres, sport, poids) \\ \hline
    Obésité abdominale & (Chirurgie bariatrique si IMC $\ge 40$ ou $\ge 35$ + comorbidités) & \checkmark \textbf{Essentiel :} Régime hypocalorique + Sport \\ \hline
    Tabagisme & \checkmark (Substitution nicotinique, Varenicline, Bupropion, TCC) & \checkmark \textbf{Arrêt volontaire} (sevrage) \\ \hline
    Sédentarité & -- & \checkmark \textbf{Activité physique} (150 min/sem) \\ \hline
    Alcoolisation excessive & \checkmark (Sevrage assisté, Baclofène, Acamprosate) & \checkmark \textbf{Modération / Arrêt} \\ \hline
    Stress chronique & \checkmark (Anxiolytiques, Antidépresseurs si trouble psy) & \checkmark (Relaxation, sommeil, organisation, TCC) \\ \hline
    Alimentation & -- & \checkmark \textbf{Régime Méditerranéen / DASH} \\ \hline
    \end{tabularx}
    \end{center}
    \emph{Note : La plupart relèvent des DEUX (prise en charge globale). Le médecin traite les chiffres (PA, HbA1c, LDL) ; le patient agit sur les causes (mode de vie).}
    \item \textbf{Pourquoi effet multiplicatif (exponentiel) et non additif ?}
    Les facteurs de risque agissent sur des \textbf{voies physiopathologiques communes et interconnectées} (endothélium vasculaire, inflammation, stress oxydatif, thrombose, remodelage cardiaque).
    \begin{itemize}
        \item \textbf{HTA} : Contrainte mécanique (cisaillement, étirement) $\to$ lésion endothéliale + hypertrophie média.
        \item \textbf{Tabagisme} : Monoxyde de carbone (hypoxie), nicotine (vasoconstriction, agrégation plaquettaire), radicaux libres (oxydation LDL) $\to$ dysfonction endothéliale majeure + état pro-thrombotique.
        \item \textbf{Diabète} : Glycation non enzymatique des protéines (collagène, LDL) $\to$ rigidité artérielle + dysfonction endothéliale (moins de NO) + état pro-inflammatoire/pro-thrombotique.
    \end{itemize}
    La coexistence de ces trois facteurs crée un \textbf{cercle vicieux synergique} : l'endothélium est agressé par 3 mécanismes simultanés, la réparation est dépassée, l'athérosclérose progresse exponentiellement, le risque thrombotique est maximal. Les scores de risque (SCORE2, Framingham) intègrent des \textbf{termes d'interaction} (produits croisés) et non de simples sommes. Cliniquement : Risque IDM/AVC multiplié par 10 à 20 vs sujet sans facteur.
\end{enumerate}
\end{corrige}

\begin{corrige}[Exercice 9 — Sujet type Probatoire : Dépistage et prise en charge de l'HTA]
\textbf{Barème indicatif : 10 points}
\begin{enumerate}
    \item \textbf{Interprétation (3 pts -- 1 par élève) :}
    \begin{itemize}
        \item \textbf{Élève A (125/78) :} \textbf{Normale} (PAS 120--129, PAD 80--84). \emph{Point de vigilance : PAS 125 = limite haute normale, surveiller IMC/poids.}
        \item \textbf{Élève B (138/88) :} \textbf{Normale haute (Préhypertension)} (PAS 130--139 \textbf{et/ou} PAD 85--89). Ici les deux sont dans l'intervalle.
        \item \textbf{Élève C (145/92) :} \textbf{Hypertension Artérielle (HTA)} (PAS $\ge 140$ \textbf{et/ou} PAD $\ge 90$). Seuil adulte atteint ($> 13$ ans).
    \end{itemize}
    \item \textbf{Diagnostic définitif ? (2 pts) :} \textbf{NON.} L'infirmière \textbf{ne doit pas poser le diagnostic définitif} sur cette unique mesure (même si moyenne de 2 prises).
    \textbf{Justification (Règle diagnostique standard) :} Le diagnostic d'HTA nécessite la confirmation de chiffres tensionnels élevés \textbf{à plusieurs reprises, lors de consultations distinctes} (généralement \textbf{3 consultations espacées de quelques semaines à 3--6 mois}), en dehors de tout contexte aigu (douleur, stress, urgence). Une seule séance ne permet pas d'éliminer l'\textbf{effet « blouse blanche »} (HTA de consultation).
    \item \textbf{Mesures hygiéno-diététiques prioritaires pour Élève C (3 pts -- 1 par mesure justifiée) :}
    \begin{enumerate}
        \item \textbf{Réduction drastique de l'apport sodé (Sel) :} $< 5$ g/j (1 c. à café). \textbf{Mécanisme :} Diminue la volémie (rétention d'eau) $\to$ baisse du débit cardiaque et de la résistance vasculaire périphérique $\to$ baisse de la PA. \emph{Action concrète :} Arrêter nouilles instantanées, conserves, bouillons cubes, sauces soja, saler à table ; cuisiner frais, épices/herbes.
        \item \textbf{Perte de poids (Surpoids IMC 28) + Activité physique régulière :} Objectif IMC $< 25$. $\ge$ 150 min/sem marche rapide/vélo/natation. \textbf{Mécanisme :} Améliore la sensibilité à l'insuline (souvent résistante chez l'obèse), diminue l'activité sympathique et le système rénine-angiotensine-aldostérone (SRAA), améliore la fonction endothéliale (vasodilatation) $\to$ baisse PA.
        \item \textbf{Adoption régime type DASH / Méditerranéen :} Riches en fruits/légumes (K$^+$, Mg$^{2+}$, fibres, antioxydants), céréales complètes, poissons, huiles végétales ; pauvres en graisses saturées, sucres ajoutés. \textbf{Mécanisme :} Apport en potassium favorise la natriurèse (excrétion Na$^+$) ; amélioration profil lipidique et inflammation vasculaire.
    \end{enumerate}
    \item \textbf{Élève B : Population à risque \& Surveillance (1 pt) :} La « normale haute » (préhypertension) n'est pas une maladie mais un \textbf{état à haut risque de progression vers l'HTA avérée} (risque $\times 2$ à $\times 3$ vs PA optimale) et de survenue précoce de complications cardiovasculaires. \textbf{Surveillance :} Contrôle TA \textbf{tous les 6--12 mois}, renforcement mesures hygiéno-diététiques (sel, poids, sport), dépistage facteurs associés (glycémie, lipides).
    \item \textbf{Complications majeures HTA non traitée (1 pt) :}
    \begin{itemize}
        \item \textbf{Cardiaque :} \textbf{Insuffisance cardiaque} (secondaire à l'hypertrophie ventriculaire gauche $\to$ dysfonction diastolique puis systolique) \textbf{OU} \textbf{Infarctus du myocarde} (athérosclérose coronarienne accélérée).
        \item \textbf{Cérébrale :} \textbf{Accident Vasculaire Cérébral (AVC)} (Ischémique par athérosclérose carotidienne/cérébrale ou Hémorragique par rupture micro-anévrismes / lipohyalinose artérioles).
    \end{itemize}
\end{enumerate}
\par\textbf{Attention [Points de vigilance Probatoire] :}
\begin{itemize}
    \item Citer explicitement la règle des « 3 consultations » pour le diagnostic HTA.
    \item Distinguer « Normale haute » (pas de traitement médicamenteux, hygiène de vie + surveillance) et « HTA » (traitement possible selon grade/risque global).
    \item Justifier le mécanisme physiologique de chaque conseil hygiéno-diététique (pas de liste brute).
    \item Préciser « AVC ischémique OU hémorragique » pour la complication cérébrale.
\end{itemize}
\end{corrige}

\begin{corrige}[Exercice 10 — Sujet type Probatoire : Prévention et éducation à la santé]
\textbf{Barème indicatif : 8 points}
\begin{enumerate}
    \item \textbf{Message de sensibilisation (Proposition) :}
    \begin{center}
    \begin{minipage}{0.9\linewidth}
    \centering
    \fbox{\parbox{0.95\linewidth}{
    \textbf{PRÉVENIR LES ACCIDENTS CARDIOVASCULAIRES : 3 GESTES QUOTIDIENS QUI SAUVENT DES VIES}\par\medskip
    \textbf{1. Assiette protectrice :} \emph{Moins de sel (5 g/j max), plus de fruits/légumes (5 portions), poissons gras, huiles colza/noix.} $\to$ \textbf{Baisse la tension artérielle et le « mauvais » cholestérol (LDL), protège les artères de l'athérosclérose.}\par\medskip
    \textbf{2. Bougez chaque jour :} \emph{30 min de marche rapide, vélo, danse... au moins 5 jours/semaine.} $\to$ \textbf{Renforce le cœur, fluidifie le sang, régule la glycémie et le poids, abaisse la tension au repos.}\par\medskip
    \textbf{3. Dépistez-vous, arrêtez le tabac :} \emph{Contrôlez votre tension, glycémie, cholestérol chaque année dès 40 ans (ou avant si risque). Zéro cigarette.} $\to$ \textbf{Détecte l'hypertension et le diabète silencieux ; le sevrage tabagique divise le risque d'infarctus par 2 en 1 an.}\par\medskip
    \textbf{URGENCE : Douleur thoracique, paralysie brutale, trouble de la parole ? $\to$ APPELEZ LE \textbf{112} IMMÉDIATEMENT (SAMU/Pompiers, gratuit). LE TEMPS, C'EST DU CŒUR ET DU CERVEAU.}
    }}
    \end{minipage}
    \end{center}
    \item \textbf{Vérification des contraintes :}
    \begin{itemize}
        \item \textbf{Longueur :} $\approx$ 9 lignes (conforme 8--10).
        \item \textbf{3 mesures distinctes :} 1. Alimentaire (Sel/Fruits/Légumes/Oméga-3), 2. Activité physique (30 min/5j), 3. Dépistage + Tabac (Zéro tabac).
        \item \textbf{Mécanismes expliqués brièvement :} Baisse PA/LDL ; Renforce cœur/fluidifie/Poids/Glycémie ; Détecte silencieux/Sevrage divise risque.
        \item \textbf{Vocabulaire grand public :} « Mauvais cholestérol », « artères bouchées », « silencieux », « divise par 2 ».
    \end{itemize}
    \item \textbf{Numéro d'urgence au Cameroun :} \textbf{112} (numéro unique d'urgence, gratuit, accessible depuis GSM et fixes).
\end{enumerate}
\par\textbf{Attention [Points de vigilance Probatoire] :}
\begin{itemize}
    \item Respecter la contrainte de longueur (8--10 lignes) : être synthétique.
    \item Ne pas oublier le \textbf{mécanisme} pour chaque geste (ex: « sel $\to$ volémie $\to$ PA » ; « sport $\to$ cœur/insuline » ; « dépistage $\to$ silencieux »).
    \item Le \textbf{112} est obligatoire et doit être mis en évidence.
    \item Ton impératif, positif, non culpabilisant.
\end{itemize}
\end{corrige}

\begin{corrige}[Exercice 11 — Sujet type Probatoire : Physiopathologie de l'AVC]
\textbf{Barème indicatif : 12 points}
\begin{enumerate}
    \item \textbf{Nommage des structures (2 pts) :}
    \begin{itemize}
        \item \textbf{Zone rose :} \textbf{Plaque d'athérome (ou atherome)}.
        \item \textbf{Zone rouge :} \textbf{Thrombus (caillot sanguin)} sur plaque rompue.
    \end{itemize}
    \item \textbf{Mécanisme de formation de la plaque d'athérome (4 pts -- étapes clés attendues) :}
    \begin{enumerate}
        \item \textbf{Excès de LDL-cholestérol (LDL-c) circulant :} Pénétration et rétention des LDL dans l'\textbf{intima} artérielle (perméabilité endothéliale augmentée par facteurs de risque : HTA, tabac, diabète).
        \item \textbf{Modification oxydative :} LDL $\xrightarrow{\text{oxydation (radicaux libres)}}$ \textbf{LDL oxydées (ox-LDL)}. Pro-inflammatoires, chimiotactiques.
        \item \textbf{Recrutement monocytaire :} ox-LDL $\to$ expression molécules d'adhésion (VCAM-1, ICAM-1) sur endothélium $\to$ adhésion et migration des \textbf{monocytes} (sang) $\to$ différenciation en \textbf{macrophages} (intima).
        \item \textbf{Formation des cellules spumeuses :} Macrophages phagocytent massivement ox-LDL via récepteurs « scavengers » (non régulés par feedback négatif) $\to$ surcharge lipidique cytoplasmique $\to$ \textbf{cellules spumeuses} (stade de \textbf{stries lipidiques}, réversible).
        \item \textbf{Évolution vers plaque fibreuse :} Macrophages libèrent cytokines (TNF-$\alpha$, IL-1, PDGF) $\to$ prolifération/migration \textbf{cellules musculaires lisses (CML)} de la média vers l'intima + synthèse matrice extracellulaire (collagène, élastine) $\to$ formation d'un \textbf{cap fibreux} recouvrant le noyau lipidique nécrotique (cellules spumeuses mortes, cholestérol cristallisé). Plaque stable (athérome fibrolipidique).
        \item \textbf{Complication (Rupture/Érosion) :} Inflammation (MMPs -- métalloprotéinases) amincit le cap fibreux $\to$ \textbf{Rupture de la plaque} $\to$ exposition du noyau thrombogène (facteur tissulaire, collagène) au sang $\to$ activation plaquettes + coagulation $\to$ \textbf{Thrombus} (zone rouge) occlusif ou sous-occlusif.
    \end{enumerate}
    \item \textbf{Conséquence hémodynamique du thrombus (1 pt) :}
    \textbf{Réduction brutale (sténose sévère) ou occlusion complète de la lumière artérielle} $\to$ \textbf{Chute brutale du débit sanguin (ischémie) en aval de l'obstacle} (territoire cérébral vascularisé par cette artère). Si occlusion complète $\to$ infarctus cérébral (mort neuronale en minutes si pas de circulation collatérale suffisante).
    \item \textbf{Différenciation AVC Ischémique vs Hémorragique (3 pts) :}
    \begin{center}
    \begin{tabularx}{\linewidth}{|l|X|X|}\hline
    \rowcolor{popblueL}
    \textbf{Critère} & \textbf{AVC Ischémique (Infarctus cérébral)} & \textbf{AVC Hémorragique} \\ \hline
    \textbf{Mécanisme lésionnel initial} & \textbf{Obstruction} artérielle (Thrombose sur athérome \textbf{OU} Embolie artérielle/cardiaque) $\to$ \textbf{Ischémie} $\to$ Nécrose (infarctus) tissu cérébral en aval. & \textbf{Rupture} artérielle (HTA, anévrisme, angiome, amyloidose) $\to$ \textbf{Hématome} intra-cérébral $\to$ Compression + toxicité sang (thrombine, fer) + ischémie par compression vaisseaux adjacents. \\ \hline
    \textbf{Thrombolyse (rt-PA / Alteplase) en urgence ?} & \textbf{OUI} (si délai $< 4$h$30$, pas de contre-indication : HTA sévère non contrôlée, hémorragie active, AVC hémorragique récent, chirurgie récente, anticoagulants efficaces, thrombopénie...). \textbf{Objectif :} Recanalisation précoce = sauver pénombre ischémique. & \textbf{NON ABSOLU (Contre-indication formelle).} Risque majeur d'aggravation de l'hématome et d'hémorragie intracérébrale fatale. \textbf{Prise en charge :} Contrôle PA strict, neurochirurgie si hématome volumineux/signe d'engagement, correction coagulopathie. \\ \hline
    \end{tabularx}
    \end{center}
    \item \textbf{Test mnémotechnique international (2 pts) :}
    \textbf{FAST} (anglais) -- Utilisé mondialement (aussi \textbf{VITE} en français).
    \begin{itemize}
        \item \textbf{F -- Face (Visage) :} Demander de sourire $\to$ \textbf{Bouche tordue, asymétrie faciale} (paralysie faciale centrale).
        \item \textbf{A -- Arm (Bras) :} Demander de lever les deux bras $\to$ \textbf{Incapacité à lever un bras} (hémiplégie/hémiparésie) ou dérive du bras.
        \item \textbf{S -- Speech (Parole) :} Demander de répéter une phrase simple $\to$ \textbf{Trouble de la parole} (aphasie : ne trouve pas les mots ; dysarthrie : articulation difficile, « patate chaude ») ou incompréhension.
        \item \textbf{T -- Time (Temps / Urgence Extrême) :} \textbf{Appeler le 112 IMMÉDIATEMENT}. Noter l'heure de début des signes. \emph{« Time is brain » (Le temps, c'est du cerveau).}
    \end{itemize}
\end{enumerate}
\par\textbf{Attention [Points de vigilance Probatoire] :}
\begin{itemize}
    \item Pour la physiopathologie de l'athérome, citer \textbf{LDL, oxydation, monocytes/macrophages, cellules spumeuses, CML, cap fibreux, rupture, thrombus} dans l'ordre chronologique.
    \item Distinguer clairement \textbf{Thrombose} (sur place) et \textbf{Embolie} (venant d'ailleurs) comme causes d'AVC ischémique.
    \item \textbf{Thrombolyse :} OUI pour ischémique (fenêtre thérapeutique), NON formel pour hémorragique. Citer au moins une contre-indication majeure.
    \item \textbf{FAST} : Donner la lettre, le mot anglais, la traduction française et le signe clinique testé.
\end{itemize}
\end{corrige}

\newpage

\coursec{Fiche bilan}

\begin{popfigure}
\begin{tikzpicture}[
  mindmap,
  concept color=popblue,
  level 1 concept/.append style={level distance=4.5cm, sibling angle=72},
  level 2 concept/.append style={level distance=3cm, sibling angle=45},
  every node/.style={font=\small, text=white},
  grow cyclic,
  root concept/.append style={concept color=popdark, minimum size=2.8cm},
  c1/.style={concept color=popgreen},
  c2/.style={concept color=poporange},
  c3/.style={concept color=poppurple},
  c4/.style={concept color=poppink},
  c5/.style={concept color=popblue!70!black}
]
\node[root concept, align=center]{Lutte contre les\\ accidents cardiovasculaires}
  child[c1] { node[concept] {Système CV}
    child { node[concept, scale=0.7] {Cœur : pompe à 4 cavités} }
    child { node[concept, scale=0.7] {Vaisseaux : artères, veines, capillaires} }
    child { node[concept, scale=0.7] {Double circulation} }
  }
  child[c2] { node[concept] {Tension artérielle}
    child { node[concept, scale=0.7] {Mesure : brassard (tensiomètre)} }
    child { node[concept, scale=0.7] {Systolique / Diastolique} }
    child { node[concept, scale=0.7] {Normale : environ 120/80 mmHg} }
    child { node[concept, scale=0.7] {HTA $\ge$ 140/90 mmHg} }
  }
  child[c3] { node[concept] {Accidents majeurs}
    child { node[concept, scale=0.7] {Infarctus du myocarde} }
    child { node[concept, scale=0.7] {AVC (ischémique / hémorragique)} }
    child { node[concept, scale=0.7] {HTA (maladie chronique)} }
    child { node[concept, scale=0.7] {Insuffisance cardiaque} }
  }
  child[c4] { node[concept] {Facteurs de risque}
    child { node[concept, scale=0.7] {Non modifiables : âge, sexe, hérédité} }
    child { node[concept, scale=0.7] {Modifiables : tabac, alcool, excès de sel, sédentarité, obésité, diabète, cholestérol, stress} }
  }
  child[c5] { node[concept] {Prévention}
    child { node[concept, scale=0.7] {Alimentation équilibrée (peu de sel et de graisses)} }
    child { node[concept, scale=0.7] {Activité physique régulière} }
    child { node[concept, scale=0.7] {Arrêt du tabac / modération de l'alcool} }
    child { node[concept, scale=0.7] {Dépistage : TA, glycémie, cholestérol} }
    child { node[concept, scale=0.7] {Observance des traitements} }
  };
\end{tikzpicture}
\legende{Carte mentale récapitulative de la leçon NCT03 : structure du système cardiovasculaire, tension artérielle, accidents majeurs, facteurs de risque et piliers de la prévention.}
\end{popfigure}

\courssub{Définitions clés}
\begin{itemize}
  \item \cle{Système cardiovasculaire} : ensemble formé du \textbf{cœur} (pompe musculaire à 4 cavités : 2 oreillettes et 2 ventricules) et des \textbf{vaisseaux sanguins} (artères, veines, capillaires) assurant la circulation du sang. La circulation est \textbf{double} : la \emph{petite circulation} (ou circulation pulmonaire : ventricule droit $\to$ poumons $\to$ oreillette gauche) et la \emph{grande circulation} (ou circulation systémique : ventricule gauche $\to$ organes $\to$ oreillette droite).
  \item \cle{Tension artérielle (TA)} : pression exercée par le sang sur la paroi des artères. \cle{TA systolique (PAS)} : pression maximale lors de la contraction des ventricules (systole). \cle{TA diastolique (PAD)} : pression minimale lors du relâchement du cœur (diastole). Unité : le millimètre de mercure (\SI{}{mmHg}).
  \item \cle{Hypertension artérielle (HTA)} : élévation persistante de la TA, avec PAS $\ge$ \SI{140}{mmHg} et/ou PAD $\ge$ \SI{90}{mmHg}, confirmée sur plusieurs mesures prises au repos. C'est une maladie chronique, longtemps silencieuse, d'où son surnom de \emph{tueur silencieux}.
  \item \cle{Infarctus du myocarde} : destruction (nécrose) d'une partie du muscle cardiaque par obstruction d'une artère coronaire (dépôt de plaque d'athérome + formation d'un caillot ou thrombose). Signe typique : douleur thoracique intense, constrictive, irradiant vers le bras gauche ou la mâchoire, ne cédant pas au repos.
  \item \cle{Accident vasculaire cérébral (AVC)} : déficit neurologique brutal par interruption de l'irrigation d'une zone du cerveau — soit par obstruction d'une artère (\emph{AVC ischémique}, le plus fréquent, environ 80\,\% des cas), soit par rupture d'une artère (\emph{AVC hémorragique}, environ 20\,\% des cas). Signe d'alerte : \cle{VITE} (Visage asymétrique, Incapacité à lever un membre, Trouble de la parole, Extrême urgence).
  \item \cle{Facteur de risque cardiovasculaire} : toute caractéristique augmentant la probabilité de survenue d'un accident cardiovasculaire. On distingue les facteurs \textbf{non modifiables} (âge, sexe masculin, antécédents familiaux) et les facteurs \textbf{modifiables} (tabagisme, excès de sel, alcool, sédentarité, obésité, diabète, excès de cholestérol, HTA, stress).
\end{itemize}

\courssub{Valeurs de référence et classifications (à connaître)}
\begin{center}
\begin{tabularx}{\linewidth}{|l|X|X|}\hline
\rowcolor{popblueL}
\textbf{Paramètre} & \textbf{Valeurs normales / Seuils} & \textbf{Commentaire} \\ \hline
TA optimale & $<$ 120 / $<$ 80 mmHg & Cible de la prévention \\ \hline
TA normale & 120--129 / 80--84 mmHg & Surveillance simple \\ \hline
TA normale élevée & 130--139 / 85--89 mmHg & Hygiène de vie stricte \\ \hline
\cle{HTA Grade 1} & 140--159 / 90--99 mmHg & Prise en charge médicale \\ \hline
\cle{HTA Grade 2} & 160--179 / 100--109 mmHg & Traitement médicamenteux + hygiène de vie \\ \hline
\cle{HTA Grade 3} & $\ge$ 180 / $\ge$ 110 mmHg & Urgence / hospitalisation possible \\ \hline
FC au repos (adulte) & 60--100 battements/min & Bradycardie $<$ 60 ; tachycardie $>$ 100 \\ \hline
FC maximale théorique & 220 $-$ âge (en années) & Effort modéré conseillé : 50--70\,\% de la FC max \\ \hline
Tour de taille à risque & H $>$ 94 cm ; F $>$ 80 cm & Obésité abdominale \\ \hline
Sel alimentaire & $<$ 5 g/jour (OMS) & Environ 1 cuillère à café \\ \hline
\end{tabularx}
\end{center}

\begin{attention}[Interpréter une mesure de TA : erreurs fréquentes]
\begin{itemize}
  \item Une mesure s'écrit toujours \textbf{PAS/PAD} (ex. 158/96 mmHg) : la valeur la plus élevée est la systolique. On retient le grade le plus sévère : 158/96 correspond à une HTA de grade 1 par la PAS (158) mais de grade 2 par la PAD (96) $\Rightarrow$ on conclut \textbf{HTA grade 2}.
  \item Une \textbf{seule} mesure élevée ne suffit pas à diagnostiquer une HTA : il faut des mesures répétées, au repos, sur plusieurs jours (l'émotion ou l'effort font monter transitoirement la TA — c'est l'\emph{effet blouse blanche} au cabinet médical).
  \item Ne pas confondre \textbf{tension artérielle} (pression, en mmHg) et \textbf{fréquence cardiaque} (nombre de battements par minute) : ce sont deux grandeurs différentes.
\end{itemize}
\end{attention}

\courssub{Méthodes express}

\begin{methode}[Mesurer la TA au brassard électronique (automesure)]
\begin{enumerate}
  \item Repos de 5 min, assis, dos soutenu, pieds au sol, sans avoir fumé ni bu de café dans l'heure.
  \item Placer le brassard adapté à la circonférence du bras, sur le bras nu, à hauteur du cœur.
  \item Effectuer 3 mesures espacées de 1 à 2 min ; retenir la moyenne des 2 dernières.
  \item Interpréter le résultat selon la grille de classification (tableau ci-dessus) et consulter si les valeurs sont élevées de façon répétée.
\end{enumerate}
\end{methode}

\begin{methode}[Identifier les facteurs de risque dans une situation donnée]
\begin{enumerate}
  \item Relever les antécédents personnels et familiaux : âge, sexe, maladies connues (diabète, HTA, excès de cholestérol).
  \item Relever les habitudes de vie : tabac, alcool, alimentation (sel, graisses, fruits et légumes), activité physique, poids et tour de taille.
  \item Classer chaque facteur : \textbf{modifiable} (cible d'une action de prévention) ou \textbf{non modifiable} (il ne peut pas être changé, mais il aggrave le risque global).
  \item Conclure sur le niveau de risque et proposer les mesures de prévention adaptées aux facteurs modifiables identifiés.
\end{enumerate}
\end{methode}

\begin{methode}[Gestes d'urgence face à un infarctus ou un AVC suspecté]
\begin{enumerate}
  \item \textbf{Alerter immédiatement les secours} : au Cameroun, SAMU 119, pompiers 118 (numéro d'urgence général 112).
  \item Installer la victime au repos, en position demi-assise (infarctus) ; en position latérale de sécurité si elle est inconsciente mais respire (AVC).
  \item Ne rien donner à manger ni à boire. Rassurer la victime et surveiller sa respiration.
  \item Si la victime ne respire plus (arrêt cardiaque) : commencer le massage cardiaque (100 à 120 compressions par minute) jusqu'à l'arrivée des secours.
\end{enumerate}
\end{methode}

\begin{aretenir}[Checklist avant le Probatoire]
\begin{itemize}
  \item[$\square$] Je sais décrire le trajet du sang dans la double circulation (cœur droit $\to$ poumons $\to$ cœur gauche $\to$ organes $\to$ cœur droit) et nommer les 4 cavités du cœur.
  \item[$\square$] Je définis la PAS, la PAD et je donne les valeurs de référence : TA normale environ 120/80 mmHg ; HTA si $\ge$ 140/90 mmHg.
  \item[$\square$] Je sais lire et interpréter une mesure de TA (ex. 158/96 = HTA grade 2) et une fréquence cardiaque (normale : 60--100 battements/min).
  \item[$\square$] Je différencie les signes d'un infarctus (douleur thoracique oppressive irradiant vers le bras gauche ou la mâchoire, sueurs, angoisse) et d'un AVC (signe VITE : Visage asymétrique, Incapacité à lever un bras, Trouble de la parole, Extrême urgence).
  \item[$\square$] Je cite au moins 5 facteurs de risque \textbf{modifiables} (tabac, excès de sel, sédentarité, alcool, obésité, diabète, excès de cholestérol, stress) et 3 \textbf{non modifiables} (âge, sexe, hérédité).
  \item[$\square$] J'explique la chaîne : \textbf{athérosclérose $\to$ plaque d'athérome $\to$ thrombose (caillot) $\to$ ischémie $\to$ nécrose} (infarctus ou AVC).
  \item[$\square$] Je propose 4 mesures de prévention concrètes : 30 min de marche par jour, 5 fruits et légumes par jour, moins de 5 g de sel par jour, arrêt du tabac, contrôle régulier de la TA.
  \item[$\square$] Je connais les numéros d'urgence au Cameroun : \textbf{119 (SAMU)}, \textbf{118 (Pompiers)}, \textbf{112 (urgences)}.
  \item[$\square$] Je sais construire un conseil d'éducation à la santé pour un hypertendu : observance du traitement, automesure de la TA, hygiène de vie, consultations régulières.
  \item[$\square$] Je distingue prévention \textbf{primaire} (éviter l'apparition de la maladie chez le sujet sain), \textbf{secondaire} (dépistage et traitement précoce) et \textbf{tertiaire} (limiter les complications et les récidives après un accident).
  \item[$\square$] Je sais analyser un document (courbe d'évolution de la TA, tableau de facteurs de risque, affiche de campagne) pour en dégager le message de prévention.
\end{itemize}
\end{aretenir}

\findecours

\end{document}
